\documentclass[aspectratio=169]{article}
\usepackage[paperwidth=25.4cm,paperheight=14.2875cm,margin=0cm]{geometry}
\usepackage{fontspec}
\usepackage{unicode-math}
\usepackage[brazilian]{babel}
\usepackage{graphicx}
\usepackage{adjustbox}
\usepackage{tikz}
\usetikzlibrary{arrows.meta,positioning,calc,shapes.geometric,fit}
\usepackage{pgfplots}
\pgfplotsset{compat=1.18}
\usepackage{qrcode}
\usepackage[hidelinks]{hyperref}
\setmainfont{Arial}
\setsansfont{Arial}
\setmathfont{latinmodern-math.otf}
\hypersetup{pdftitle={Algoritmos Exatos Para o Problema de Visita de Polígonos},pdfauthor={Gabriel Freire Ushijima; Ernesto G. Birgin},pdfsubject={34º SIICUSP, IME-USP}}
\pagestyle{empty}
\setlength{\parindent}{0pt}
\definecolor{Ink}{HTML}{142D38}
\definecolor{Muted}{HTML}{47616A}
\definecolor{Accent}{HTML}{C2501A}
\definecolor{Pale}{HTML}{EEF4F1}
\definecolor{GoldText}{HTML}{9A6B12}
\definecolor{Rule}{HTML}{C5D3CF}
\definecolor{Teal}{HTML}{1F7A73}
\definecolor{MapDark}{HTML}{102C38}
\definecolor{Route}{HTML}{F59A4A}


% Tamanhos escritos em pt de um slide de 33,87 cm (padrão PowerPoint 16:9); a página
% tem 25,4 cm, então escalamos por 0,75 para que a legibilidade equivalha ao guia do IME.
\newcommand{\sz}[2]{\fontsize{\dimexpr0.75pt*#1\relax}{\dimexpr0.75pt*#2\relax}\selectfont}
\newcommand{\ink}{\color{Ink}}
\newcommand{\muted}{\color{Muted}}
\newcommand{\hi}{\color{Accent}}
\newcommand{\appurl}{https://yushipy.github.io/TouringPolygons/apps/siicusp34/}
% Rodapé comum: financiamento e número do slide
\newcommand{\footerbar}[1]{%
  \draw[Rule,line width=0.8pt] ([xshift=1.2cm,yshift=1.15cm]current page.south west) -- ([xshift=-1.2cm,yshift=1.15cm]current page.south east);
  % três blocos com espaços iguais entre eles (\hfill reparte o excedente igualmente)
  \node[anchor=south west,inner sep=0] at ([xshift=1.2cm,yshift=0.3cm]current page.south west) {\sffamily\sz{16}{19}\muted\hbox to 23cm{Gabriel Freire Ushijima\hfill Algoritmos Exatos Para o Problema de Visita de Polígonos\hfill #1\,/\,5}};}
\newcommand{\slidetitle}[1]{%
  \node[anchor=north west,inner sep=0,text width=21.5cm] at ([xshift=1.2cm,yshift=-0.75cm]current page.north west) {\sffamily\bfseries\sz{36}{40}\ink #1};}
\begin{document}
\sffamily\raggedright\hyphenpenalty=10000\exhyphenpenalty=10000\tolerance=9999

%%%%%%%%%%%%%%%%%%%%%%%%%%%%%%%% 1. Capa
% O título é uma instância de TPP: cada letra é uma região; a rota é o caminho
% mínimo que toca todas (make_title_instance.py).
\begin{tikzpicture}[remember picture,overlay]
  \fill[white] (current page.south west) rectangle (current page.north east);
  \node[anchor=north,inner sep=0] at ([yshift=-0.9cm]current page.north) {%
    \begin{tikzpicture}[x=1.95cm,y=1.95cm]\begin{scope}[shift={(0.4700,3.7504)}]
\fill[color={rgb,255:red,20;green,48;blue,58},even odd rule] (0.5610,0.0000) -- (0.7183,0.0000) -- (0.4316,0.7158) -- (0.2788,0.7158) -- (0.0000,0.0000) -- (0.1533,0.0000) -- (0.2124,0.1626) -- (0.4985,0.1626) -- cycle (0.3535,0.5488) -- (0.4521,0.2832) -- (0.2568,0.2832) -- cycle;
\fill[color={rgb,255:red,21;green,51;blue,61},even odd rule] (0.7939,0.0000) -- (0.9312,0.0000) -- (0.9312,0.7158) -- (0.7939,0.7158) -- cycle;
\fill[color={rgb,255:red,21;green,54;blue,63},even odd rule] (1.2158,-0.0532) -- (1.0591,-0.0342) -- (1.0640,-0.0918) -- (1.0896,-0.1412) -- (1.1138,-0.1653) -- (1.1372,-0.1806) -- (1.2001,-0.2018) -- (1.3008,-0.2104) -- (1.3983,-0.2008) -- (1.4379,-0.1878) -- (1.4704,-0.1704) -- (1.4959,-0.1485) -- (1.5135,-0.1245) -- (1.5369,-0.0653) -- (1.5459,0.0102) -- (1.5469,0.5186) -- (1.4185,0.5186) -- (1.4185,0.4458) -- (1.3761,0.4903) -- (1.3272,0.5184) -- (1.2717,0.5299) -- (1.2150,0.5260) -- (1.1741,0.5131) -- (1.1370,0.4915) -- (1.1038,0.4614) -- (1.0763,0.4232) -- (1.0567,0.3772) -- (1.0471,0.3377) -- (1.0413,0.2783) -- (1.0427,0.2244) -- (1.0502,0.1777) -- (1.0636,0.1354) -- (1.0831,0.0972) -- (1.1133,0.0583) -- (1.1506,0.0280) -- (1.1821,0.0124) -- (1.2167,0.0031) -- (1.2660,0.0003) -- (1.3098,0.0082) -- (1.3497,0.0266) -- (1.3855,0.0554) -- (1.4097,0.0840) -- (1.4067,-0.0412) -- (1.3967,-0.0684) -- (1.3804,-0.0865) -- (1.3609,-0.0964) -- (1.3319,-0.1032) -- (1.2630,-0.1027) -- (1.2289,-0.0864) -- cycle (1.1848,0.3152) -- (1.1968,0.3591) -- (1.2177,0.3916) -- (1.2499,0.4157) -- (1.2888,0.4251) -- (1.3305,0.4198) -- (1.3658,0.3995) -- (1.3934,0.3641) -- (1.4058,0.3292) -- (1.4109,0.2954) -- (1.4072,0.2141) -- (1.3856,0.1606) -- (1.3639,0.1368) -- (1.3389,0.1207) -- (1.3108,0.1122) -- (1.2754,0.1122) -- (1.2493,0.1204) -- (1.2261,0.1361) -- (1.2061,0.1594) -- (1.1941,0.1841) -- (1.1848,0.2228) -- cycle;
\fill[color={rgb,255:red,23;green,66;blue,72},even odd rule] (1.6510,0.2751) -- (1.6542,0.2134) -- (1.6640,0.1655) -- (1.6805,0.1229) -- (1.7087,0.0790) -- (1.7456,0.0436) -- (1.7830,0.0205) -- (1.8408,-0.0015) -- (1.8924,-0.0106) -- (1.9471,-0.0105) -- (1.9998,-0.0008) -- (2.0476,0.0185) -- (2.0905,0.0475) -- (2.1279,0.0857) -- (2.1564,0.1294) -- (2.1754,0.1782) -- (2.1848,0.2320) -- (2.1849,0.2893) -- (2.1755,0.3426) -- (2.1567,0.3909) -- (2.1285,0.4341) -- (2.0914,0.4718) -- (2.0485,0.5004) -- (2.0003,0.5195) -- (1.9467,0.5291) -- (1.8808,0.5282) -- (1.8284,0.5171) -- (1.7800,0.4966) -- (1.7377,0.4671) -- (1.7032,0.4290) -- (1.6767,0.3829) -- (1.6573,0.3256) -- cycle (1.7916,0.2498) -- (1.7967,0.3121) -- (1.8121,0.3547) -- (1.8425,0.3914) -- (1.8806,0.4127) -- (1.9316,0.4178) -- (1.9734,0.4055) -- (2.0129,0.3722) -- (2.0358,0.3276) -- (2.0426,0.2961) -- (2.0448,0.2506) -- (2.0334,0.1841) -- (2.0209,0.1578) -- (2.0039,0.1361) -- (1.9788,0.1161) -- (1.9563,0.1059) -- (1.9251,0.1003) -- (1.8927,0.1027) -- (1.8634,0.1131) -- (1.8425,0.1272) -- (1.8195,0.1519) -- (1.8058,0.1770) -- (1.7951,0.2146) -- cycle;
\fill[color={rgb,255:red,23;green,69;blue,75},even odd rule] (2.4248,0.0000) -- (2.4259,0.2340) -- (2.4302,0.2992) -- (2.4377,0.3391) -- (2.4583,0.3818) -- (2.4900,0.4033) -- (2.5130,0.4072) -- (2.5357,0.4050) -- (2.5811,0.3848) -- (2.6235,0.5044) -- (2.5792,0.5238) -- (2.5291,0.5302) -- (2.4839,0.5193) -- (2.4498,0.4923) -- (2.4150,0.4448) -- (2.4150,0.5186) -- (2.2876,0.5186) -- (2.2876,0.0000) -- cycle;
\fill[color={rgb,255:red,24;green,72;blue,77},even odd rule] (2.6826,0.5889) -- (2.8198,0.5889) -- (2.8198,0.7158) -- (2.6826,0.7158) -- cycle;
\fill[color={rgb,255:red,24;green,75;blue,79},even odd rule] (2.6826,0.0000) -- (2.8198,0.0000) -- (2.8198,0.5186) -- (2.6826,0.5186) -- cycle;
\fill[color={rgb,255:red,25;green,78;blue,81},even odd rule] (3.1982,0.5186) -- (3.1045,0.5186) -- (3.1045,0.7017) -- (2.9668,0.6216) -- (2.9668,0.5186) -- (2.9038,0.5186) -- (2.9038,0.4092) -- (2.9668,0.4092) -- (2.9672,0.1424) -- (2.9754,0.0667) -- (2.9856,0.0409) -- (3.0007,0.0213) -- (3.0236,0.0057) -- (3.0535,-0.0058) -- (3.0874,-0.0113) -- (3.1306,-0.0104) -- (3.1718,-0.0037) -- (3.2095,0.0088) -- (3.1978,0.1152) -- (3.1532,0.1029) -- (3.1244,0.1059) -- (3.1080,0.1237) -- (3.1049,0.1600) -- (3.1045,0.4092) -- (3.1982,0.4092) -- cycle;
\fill[color={rgb,255:red,25;green,81;blue,84},even odd rule] (3.2832,0.5186) -- (3.2832,0.0000) -- (3.4204,0.0000) -- (3.4206,0.2700) -- (3.4288,0.3502) -- (3.4417,0.3809) -- (3.4617,0.4037) -- (3.4912,0.4199) -- (3.5249,0.4253) -- (3.5511,0.4219) -- (3.5704,0.4117) -- (3.5852,0.3921) -- (3.5926,0.3675) -- (3.5966,0.3023) -- (3.5967,0.0000) -- (3.7339,0.0000) -- (3.7341,0.2667) -- (3.7425,0.3471) -- (3.7555,0.3794) -- (3.7756,0.4031) -- (3.8046,0.4197) -- (3.8359,0.4253) -- (3.8627,0.4212) -- (3.8837,0.4088) -- (3.8962,0.3930) -- (3.9042,0.3659) -- (3.9086,0.3058) -- (3.9087,0.0000) -- (4.0459,0.0000) -- (4.0424,0.3955) -- (4.0351,0.4300) -- (4.0237,0.4550) -- (4.0034,0.4826) -- (3.9822,0.5009) -- (3.9518,0.5172) -- (3.9164,0.5270) -- (3.8760,0.5303) -- (3.8377,0.5263) -- (3.8015,0.5146) -- (3.7721,0.4982) -- (3.7178,0.4478) -- (3.6872,0.4878) -- (3.6484,0.5146) -- (3.6011,0.5283) -- (3.5483,0.5290) -- (3.5046,0.5187) -- (3.4740,0.5042) -- (3.4360,0.4758) -- (3.4097,0.4478) -- (3.4097,0.5186) -- cycle;
\fill[color={rgb,255:red,26;green,90;blue,91},even odd rule] (4.1510,0.2751) -- (4.1542,0.2134) -- (4.1640,0.1655) -- (4.1805,0.1229) -- (4.2087,0.0790) -- (4.2456,0.0436) -- (4.2830,0.0205) -- (4.3408,-0.0015) -- (4.3924,-0.0106) -- (4.4471,-0.0105) -- (4.4998,-0.0008) -- (4.5476,0.0185) -- (4.5905,0.0475) -- (4.6279,0.0857) -- (4.6564,0.1294) -- (4.6754,0.1782) -- (4.6848,0.2320) -- (4.6849,0.2893) -- (4.6755,0.3426) -- (4.6567,0.3909) -- (4.6285,0.4341) -- (4.5914,0.4718) -- (4.5485,0.5004) -- (4.5003,0.5195) -- (4.4467,0.5291) -- (4.3808,0.5282) -- (4.3284,0.5171) -- (4.2800,0.4966) -- (4.2377,0.4671) -- (4.2032,0.4290) -- (4.1767,0.3829) -- (4.1573,0.3256) -- cycle (4.2916,0.2498) -- (4.2967,0.3121) -- (4.3121,0.3547) -- (4.3425,0.3914) -- (4.3806,0.4127) -- (4.4316,0.4178) -- (4.4734,0.4055) -- (4.5129,0.3722) -- (4.5358,0.3276) -- (4.5426,0.2961) -- (4.5448,0.2506) -- (4.5334,0.1841) -- (4.5209,0.1578) -- (4.5039,0.1361) -- (4.4788,0.1161) -- (4.4563,0.1059) -- (4.4251,0.1003) -- (4.3927,0.1027) -- (4.3634,0.1131) -- (4.3425,0.1272) -- (4.3195,0.1519) -- (4.3058,0.1770) -- (4.2951,0.2146) -- cycle;
\fill[color={rgb,255:red,27;green,93;blue,93},even odd rule] (4.7451,0.1479) -- (4.7654,0.0970) -- (4.7977,0.0547) -- (4.8421,0.0212) -- (4.8884,0.0019) -- (4.9555,-0.0102) -- (5.0354,-0.0099) -- (5.1069,0.0045) -- (5.1541,0.0275) -- (5.2013,0.0735) -- (5.2258,0.1277) -- (5.2295,0.1621) -- (5.2268,0.1930) -- (5.2188,0.2203) -- (5.2053,0.2440) -- (5.1746,0.2730) -- (5.1265,0.2975) -- (4.9154,0.3582) -- (4.8969,0.3721) -- (4.8917,0.3908) -- (4.8990,0.4092) -- (4.9171,0.4216) -- (4.9485,0.4295) -- (4.9849,0.4316) -- (5.0229,0.4276) -- (5.0480,0.4175) -- (5.0685,0.3987) -- (5.0825,0.3691) -- (5.2119,0.3931) -- (5.1843,0.4526) -- (5.1589,0.4814) -- (5.1268,0.5037) -- (5.0849,0.5193) -- (5.0436,0.5269) -- (4.9824,0.5303) -- (4.9183,0.5258) -- (4.8655,0.5123) -- (4.8316,0.4951) -- (4.7893,0.4536) -- (4.7675,0.4034) -- (4.7663,0.3442) -- (4.7759,0.3120) -- (4.7882,0.2905) -- (4.8257,0.2544) -- (4.8653,0.2349) -- (5.0731,0.1754) -- (5.0887,0.1594) -- (5.0907,0.1327) -- (5.0799,0.1124) -- (5.0587,0.0984) -- (5.0243,0.0893) -- (4.9819,0.0877) -- (4.9461,0.0940) -- (4.9152,0.1108) -- (4.8940,0.1376) -- (4.8828,0.1689) -- cycle;
\fill[color={rgb,255:red,36;green,129;blue,120},even odd rule] (5.6284,0.0000) -- (6.1729,0.0000) -- (6.1729,0.1206) -- (5.7729,0.1206) -- (5.7729,0.3154) -- (6.1323,0.3154) -- (6.1323,0.4360) -- (5.7729,0.4360) -- (5.7729,0.5947) -- (6.1592,0.5947) -- (6.1592,0.7158) -- (5.6284,0.7158) -- cycle;
\fill[color={rgb,255:red,37;green,130;blue,121},even odd rule] (6.2285,0.0000) -- (6.3892,0.0000) -- (6.4956,0.1606) -- (6.6011,0.0000) -- (6.7695,0.0000) -- (6.5776,0.2729) -- (6.7534,0.5186) -- (6.5923,0.5186) -- (6.4956,0.3760) -- (6.4038,0.5186) -- (6.2363,0.5186) -- (6.4155,0.2671) -- cycle;
\fill[color={rgb,255:red,38;green,131;blue,122},even odd rule] (6.8286,0.3828) -- (6.9531,0.3604) -- (6.9699,0.3941) -- (6.9917,0.4143) -- (7.0221,0.4239) -- (7.0626,0.4247) -- (7.0977,0.4187) -- (7.1172,0.4084) -- (7.1305,0.3894) -- (7.1360,0.3661) -- (7.1367,0.3384) -- (6.9153,0.2816) -- (6.8723,0.2597) -- (6.8439,0.2323) -- (6.8197,0.1832) -- (6.8152,0.1254) -- (6.8291,0.0742) -- (6.8449,0.0488) -- (6.8665,0.0265) -- (6.9000,0.0053) -- (6.9315,-0.0056) -- (6.9673,-0.0110) -- (7.0092,-0.0107) -- (7.0467,-0.0036) -- (7.0822,0.0099) -- (7.1158,0.0300) -- (7.1475,0.0566) -- (7.1655,0.0000) -- (7.3013,0.0000) -- (7.2806,0.0531) -- (7.2726,0.0990) -- (7.2714,0.3415) -- (7.2669,0.4059) -- (7.2532,0.4534) -- (7.2242,0.4873) -- (7.1780,0.5136) -- (7.1155,0.5272) -- (7.0287,0.5297) -- (6.9553,0.5188) -- (6.9145,0.5026) -- (6.8771,0.4736) -- (6.8486,0.4337) -- cycle (7.1367,0.2490) -- (7.1367,0.2156) -- (7.1326,0.1597) -- (7.1212,0.1329) -- (7.1017,0.1117) -- (7.0641,0.0911) -- (7.0278,0.0850) -- (6.9942,0.0916) -- (6.9658,0.1141) -- (6.9538,0.1385) -- (6.9521,0.1629) -- (6.9602,0.1855) -- (6.9753,0.2017) -- (7.0108,0.2174) -- cycle;
\fill[color={rgb,255:red,40;green,133;blue,124},even odd rule] (7.6445,0.5186) -- (7.5508,0.5186) -- (7.5508,0.7017) -- (7.4131,0.6216) -- (7.4131,0.5186) -- (7.3501,0.5186) -- (7.3501,0.4092) -- (7.4131,0.4092) -- (7.4135,0.1424) -- (7.4217,0.0667) -- (7.4318,0.0409) -- (7.4470,0.0213) -- (7.4699,0.0057) -- (7.4998,-0.0058) -- (7.5337,-0.0113) -- (7.5768,-0.0104) -- (7.6181,-0.0037) -- (7.6558,0.0088) -- (7.6440,0.1152) -- (7.5995,0.1029) -- (7.5707,0.1059) -- (7.5542,0.1237) -- (7.5512,0.1600) -- (7.5508,0.4092) -- (7.6445,0.4092) -- cycle;
\fill[color={rgb,255:red,41;green,134;blue,125},even odd rule] (7.7081,0.2751) -- (7.7113,0.2134) -- (7.7212,0.1655) -- (7.7376,0.1229) -- (7.7658,0.0790) -- (7.8027,0.0436) -- (7.8401,0.0205) -- (7.8980,-0.0015) -- (7.9496,-0.0106) -- (8.0043,-0.0105) -- (8.0570,-0.0008) -- (8.1048,0.0185) -- (8.1476,0.0475) -- (8.1850,0.0857) -- (8.2135,0.1294) -- (8.2325,0.1782) -- (8.2420,0.2320) -- (8.2420,0.2893) -- (8.2326,0.3426) -- (8.2138,0.3909) -- (8.1856,0.4341) -- (8.1485,0.4718) -- (8.1056,0.5004) -- (8.0574,0.5195) -- (8.0039,0.5291) -- (7.9379,0.5282) -- (7.8855,0.5171) -- (7.8372,0.4966) -- (7.7948,0.4671) -- (7.7603,0.4290) -- (7.7338,0.3829) -- (7.7145,0.3256) -- cycle (7.8488,0.2498) -- (7.8538,0.3121) -- (7.8692,0.3547) -- (7.8996,0.3914) -- (7.9377,0.4127) -- (7.9887,0.4178) -- (8.0305,0.4055) -- (8.0701,0.3722) -- (8.0930,0.3276) -- (8.0998,0.2961) -- (8.1019,0.2506) -- (8.0905,0.1841) -- (8.0780,0.1578) -- (8.0611,0.1361) -- (8.0360,0.1161) -- (8.0134,0.1059) -- (7.9822,0.1003) -- (7.9498,0.1027) -- (7.9206,0.1131) -- (7.8996,0.1272) -- (7.8767,0.1519) -- (7.8629,0.1770) -- (7.8522,0.2146) -- cycle;
\fill[color={rgb,255:red,42;green,136;blue,126},even odd rule] (8.3022,0.1479) -- (8.3225,0.0970) -- (8.3548,0.0547) -- (8.3992,0.0212) -- (8.4456,0.0019) -- (8.5126,-0.0102) -- (8.5925,-0.0099) -- (8.6640,0.0045) -- (8.7113,0.0275) -- (8.7585,0.0735) -- (8.7829,0.1277) -- (8.7866,0.1621) -- (8.7839,0.1930) -- (8.7759,0.2203) -- (8.7625,0.2440) -- (8.7318,0.2730) -- (8.6837,0.2975) -- (8.4726,0.3582) -- (8.4541,0.3721) -- (8.4488,0.3908) -- (8.4561,0.4092) -- (8.4742,0.4216) -- (8.5057,0.4295) -- (8.5420,0.4316) -- (8.5800,0.4276) -- (8.6052,0.4175) -- (8.6256,0.3987) -- (8.6396,0.3691) -- (8.7690,0.3931) -- (8.7415,0.4526) -- (8.7160,0.4814) -- (8.6840,0.5037) -- (8.6420,0.5193) -- (8.6007,0.5269) -- (8.5396,0.5303) -- (8.4754,0.5258) -- (8.4227,0.5123) -- (8.3887,0.4951) -- (8.3464,0.4536) -- (8.3246,0.4034) -- (8.3235,0.3442) -- (8.3331,0.3120) -- (8.3453,0.2905) -- (8.3828,0.2544) -- (8.4225,0.2349) -- (8.6303,0.1754) -- (8.6458,0.1594) -- (8.6478,0.1327) -- (8.6370,0.1124) -- (8.6158,0.0984) -- (8.5815,0.0893) -- (8.5390,0.0877) -- (8.5032,0.0940) -- (8.4723,0.1108) -- (8.4511,0.1376) -- (8.4399,0.1689) -- cycle;
\fill[color={rgb,255:red,20;green,45;blue,56},even odd rule] (0.5442,-0.4642) -- (0.5442,-1.1800) -- (0.6887,-1.1800) -- (0.6887,-0.9100) -- (0.8291,-0.9093) -- (0.9118,-0.9032) -- (0.9613,-0.8911) -- (1.0078,-0.8658) -- (1.0524,-0.8214) -- (1.0812,-0.7651) -- (1.0922,-0.7005) -- (1.0867,-0.6247) -- (1.0692,-0.5739) -- (1.0403,-0.5321) -- (0.9980,-0.4963) -- (0.9556,-0.4771) -- (0.9140,-0.4693) -- (0.8496,-0.4652) -- cycle (0.6887,-0.7884) -- (0.6887,-0.5853) -- (0.8625,-0.5902) -- (0.9017,-0.6055) -- (0.9285,-0.6319) -- (0.9422,-0.6677) -- (0.9413,-0.7091) -- (0.9272,-0.7420) -- (0.9019,-0.7670) -- (0.8740,-0.7798) -- (0.8391,-0.7856) -- cycle;
\fill[color={rgb,255:red,22;green,57;blue,65},even odd rule] (1.1882,-0.7972) -- (1.3127,-0.8196) -- (1.3295,-0.7859) -- (1.3513,-0.7657) -- (1.3818,-0.7561) -- (1.4222,-0.7553) -- (1.4573,-0.7613) -- (1.4768,-0.7716) -- (1.4902,-0.7906) -- (1.4957,-0.8139) -- (1.4963,-0.8416) -- (1.2749,-0.8984) -- (1.2319,-0.9203) -- (1.2035,-0.9477) -- (1.1793,-0.9968) -- (1.1748,-1.0546) -- (1.1887,-1.1058) -- (1.2045,-1.1312) -- (1.2261,-1.1535) -- (1.2596,-1.1747) -- (1.2911,-1.1856) -- (1.3269,-1.1910) -- (1.3688,-1.1907) -- (1.4063,-1.1836) -- (1.4418,-1.1701) -- (1.4754,-1.1500) -- (1.5071,-1.1234) -- (1.5251,-1.1800) -- (1.6609,-1.1800) -- (1.6402,-1.1269) -- (1.6322,-1.0810) -- (1.6310,-0.8385) -- (1.6265,-0.7741) -- (1.6128,-0.7266) -- (1.5838,-0.6927) -- (1.5376,-0.6664) -- (1.4751,-0.6528) -- (1.3884,-0.6503) -- (1.3149,-0.6612) -- (1.2741,-0.6774) -- (1.2368,-0.7064) -- (1.2082,-0.7463) -- cycle (1.4963,-0.9310) -- (1.4963,-0.9644) -- (1.4922,-1.0203) -- (1.4809,-1.0471) -- (1.4613,-1.0683) -- (1.4237,-1.0889) -- (1.3875,-1.0950) -- (1.3538,-1.0884) -- (1.3255,-1.0659) -- (1.3134,-1.0415) -- (1.3117,-1.0171) -- (1.3198,-0.9945) -- (1.3349,-0.9783) -- (1.3704,-0.9626) -- cycle;
\fill[color={rgb,255:red,22;green,60;blue,68},even odd rule] (1.8977,-1.1800) -- (1.8988,-0.9460) -- (1.9031,-0.8808) -- (1.9106,-0.8409) -- (1.9312,-0.7982) -- (1.9629,-0.7767) -- (1.9859,-0.7728) -- (2.0086,-0.7750) -- (2.0540,-0.7952) -- (2.0964,-0.6756) -- (2.0521,-0.6562) -- (2.0020,-0.6498) -- (1.9568,-0.6607) -- (1.9227,-0.6877) -- (1.8879,-0.7352) -- (1.8879,-0.6614) -- (1.7605,-0.6614) -- (1.7605,-1.1800) -- cycle;
\fill[color={rgb,255:red,23;green,63;blue,70},even odd rule] (2.1335,-0.7972) -- (2.2581,-0.8196) -- (2.2748,-0.7859) -- (2.2966,-0.7657) -- (2.3271,-0.7561) -- (2.3675,-0.7553) -- (2.4026,-0.7613) -- (2.4221,-0.7716) -- (2.4355,-0.7906) -- (2.4410,-0.8139) -- (2.4417,-0.8416) -- (2.2202,-0.8984) -- (2.1772,-0.9203) -- (2.1488,-0.9477) -- (2.1246,-0.9968) -- (2.1201,-1.0546) -- (2.1340,-1.1058) -- (2.1498,-1.1312) -- (2.1714,-1.1535) -- (2.2049,-1.1747) -- (2.2365,-1.1856) -- (2.2723,-1.1910) -- (2.3141,-1.1907) -- (2.3516,-1.1836) -- (2.3871,-1.1701) -- (2.4207,-1.1500) -- (2.4524,-1.1234) -- (2.4705,-1.1800) -- (2.6062,-1.1800) -- (2.5856,-1.1269) -- (2.5775,-1.0810) -- (2.5763,-0.8385) -- (2.5718,-0.7741) -- (2.5581,-0.7266) -- (2.5291,-0.6927) -- (2.4829,-0.6664) -- (2.4204,-0.6528) -- (2.3337,-0.6503) -- (2.2602,-0.6612) -- (2.2194,-0.6774) -- (2.1821,-0.7064) -- (2.1535,-0.7463) -- cycle (2.4417,-0.9310) -- (2.4416,-0.9644) -- (2.4375,-1.0203) -- (2.4262,-1.0471) -- (2.4066,-1.0683) -- (2.3690,-1.0889) -- (2.3328,-1.0950) -- (2.2991,-1.0884) -- (2.2708,-1.0659) -- (2.2587,-1.0415) -- (2.2570,-1.0171) -- (2.2651,-0.9945) -- (2.2802,-0.9783) -- (2.3157,-0.9626) -- cycle;
\fill[color={rgb,255:red,26;green,84;blue,86},even odd rule] (2.9579,-0.9049) -- (2.9611,-0.9666) -- (2.9709,-1.0145) -- (2.9874,-1.0571) -- (3.0155,-1.1010) -- (3.0524,-1.1364) -- (3.0898,-1.1595) -- (3.1477,-1.1815) -- (3.1993,-1.1906) -- (3.2540,-1.1905) -- (3.3067,-1.1808) -- (3.3545,-1.1615) -- (3.3974,-1.1325) -- (3.4348,-1.0943) -- (3.4633,-1.0506) -- (3.4822,-1.0018) -- (3.4917,-0.9480) -- (3.4917,-0.8907) -- (3.4823,-0.8374) -- (3.4635,-0.7891) -- (3.4353,-0.7459) -- (3.3983,-0.7082) -- (3.3554,-0.6796) -- (3.3071,-0.6605) -- (3.2536,-0.6509) -- (3.1877,-0.6518) -- (3.1353,-0.6629) -- (3.0869,-0.6834) -- (3.0445,-0.7129) -- (3.0101,-0.7510) -- (2.9836,-0.7971) -- (2.9642,-0.8544) -- cycle (3.0985,-0.9302) -- (3.1035,-0.8679) -- (3.1190,-0.8253) -- (3.1493,-0.7886) -- (3.1875,-0.7673) -- (3.2385,-0.7622) -- (3.2803,-0.7745) -- (3.3198,-0.8078) -- (3.3427,-0.8524) -- (3.3495,-0.8839) -- (3.3517,-0.9294) -- (3.3403,-0.9959) -- (3.3278,-1.0222) -- (3.3108,-1.0439) -- (3.2857,-1.0639) -- (3.2632,-1.0741) -- (3.2320,-1.0797) -- (3.1996,-1.0773) -- (3.1703,-1.0669) -- (3.1493,-1.0528) -- (3.1264,-1.0281) -- (3.1127,-1.0030) -- (3.1020,-0.9654) -- cycle;
\fill[color={rgb,255:red,26;green,87;blue,88},even odd rule] (3.8792,-0.4642) -- (3.8792,-1.1800) -- (4.0237,-1.1800) -- (4.0237,-0.9100) -- (4.1641,-0.9093) -- (4.2468,-0.9032) -- (4.2962,-0.8911) -- (4.3428,-0.8658) -- (4.3873,-0.8214) -- (4.4162,-0.7651) -- (4.4271,-0.7005) -- (4.4217,-0.6247) -- (4.4041,-0.5739) -- (4.3753,-0.5321) -- (4.3330,-0.4963) -- (4.2906,-0.4771) -- (4.2490,-0.4693) -- (4.1845,-0.4652) -- cycle (4.0237,-0.7884) -- (4.0237,-0.5853) -- (4.1975,-0.5902) -- (4.2367,-0.6055) -- (4.2634,-0.6319) -- (4.2771,-0.6677) -- (4.2763,-0.7091) -- (4.2622,-0.7420) -- (4.2369,-0.7670) -- (4.2090,-0.7798) -- (4.1741,-0.7856) -- cycle;
\fill[color={rgb,255:red,27;green,96;blue,95},even odd rule] (4.6765,-1.1800) -- (4.6776,-0.9460) -- (4.6819,-0.8808) -- (4.6895,-0.8409) -- (4.7100,-0.7982) -- (4.7417,-0.7767) -- (4.7647,-0.7728) -- (4.7874,-0.7750) -- (4.8328,-0.7952) -- (4.8752,-0.6756) -- (4.8309,-0.6562) -- (4.7808,-0.6498) -- (4.7356,-0.6607) -- (4.7015,-0.6877) -- (4.6667,-0.7352) -- (4.6667,-0.6614) -- (4.5393,-0.6614) -- (4.5393,-1.1800) -- cycle;
\fill[color={rgb,255:red,28;green,99;blue,98},even odd rule] (4.9027,-0.9049) -- (4.9059,-0.9666) -- (4.9157,-1.0145) -- (4.9322,-1.0571) -- (4.9604,-1.1010) -- (4.9973,-1.1364) -- (5.0347,-1.1595) -- (5.0925,-1.1815) -- (5.1442,-1.1906) -- (5.1988,-1.1905) -- (5.2516,-1.1808) -- (5.2993,-1.1615) -- (5.3422,-1.1325) -- (5.3796,-1.0943) -- (5.4081,-1.0506) -- (5.4271,-1.0018) -- (5.4366,-0.9480) -- (5.4366,-0.8907) -- (5.4272,-0.8374) -- (5.4084,-0.7891) -- (5.3802,-0.7459) -- (5.3431,-0.7082) -- (5.3002,-0.6796) -- (5.2520,-0.6605) -- (5.1984,-0.6509) -- (5.1325,-0.6518) -- (5.0801,-0.6629) -- (5.0317,-0.6834) -- (4.9894,-0.7129) -- (4.9549,-0.7510) -- (4.9284,-0.7971) -- (4.9090,-0.8544) -- cycle (5.0434,-0.9302) -- (5.0484,-0.8679) -- (5.0638,-0.8253) -- (5.0942,-0.7886) -- (5.1323,-0.7673) -- (5.1833,-0.7622) -- (5.2251,-0.7745) -- (5.2647,-0.8078) -- (5.2875,-0.8524) -- (5.2944,-0.8839) -- (5.2965,-0.9294) -- (5.2851,-0.9959) -- (5.2726,-1.0222) -- (5.2556,-1.0439) -- (5.2305,-1.0639) -- (5.2080,-1.0741) -- (5.1768,-1.0797) -- (5.1444,-1.0773) -- (5.1151,-1.0669) -- (5.0942,-1.0528) -- (5.0713,-1.0281) -- (5.0575,-1.0030) -- (5.0468,-0.9654) -- cycle;
\fill[color={rgb,255:red,33;green,124;blue,117},even odd rule] (5.5393,-0.4642) -- (5.5393,-1.1800) -- (5.6667,-1.1800) -- (5.6667,-1.1038) -- (5.6919,-1.1334) -- (5.7204,-1.1571) -- (5.7520,-1.1747) -- (5.7847,-1.1862) -- (5.8510,-1.1906) -- (5.8928,-1.1816) -- (5.9219,-1.1690) -- (5.9662,-1.1368) -- (5.9972,-1.1010) -- (6.0210,-1.0575) -- (6.0368,-1.0064) -- (6.0447,-0.9476) -- (6.0435,-0.8693) -- (6.0302,-0.8019) -- (6.0108,-0.7564) -- (5.9836,-0.7183) -- (5.9417,-0.6822) -- (5.8933,-0.6594) -- (5.8501,-0.6508) -- (5.8055,-0.6509) -- (5.7553,-0.6636) -- (5.7097,-0.6904) -- (5.6765,-0.7220) -- (5.6765,-0.4642) -- cycle (5.6756,-0.9194) -- (5.6837,-0.8413) -- (5.6971,-0.8088) -- (5.7167,-0.7841) -- (5.7451,-0.7643) -- (5.7724,-0.7561) -- (5.8085,-0.7561) -- (5.8361,-0.7645) -- (5.8601,-0.7807) -- (5.8767,-0.7991) -- (5.8898,-0.8234) -- (5.8989,-0.8539) -- (5.9045,-0.9007) -- (5.9039,-0.9519) -- (5.8970,-0.9945) -- (5.8870,-1.0223) -- (5.8606,-1.0577) -- (5.8278,-1.0779) -- (5.7884,-1.0831) -- (5.7484,-1.0729) -- (5.7151,-1.0472) -- (5.6923,-1.0127) -- (5.6804,-0.9720) -- cycle;
\fill[color={rgb,255:red,34;green,126;blue,118},even odd rule] (6.1560,-1.1800) -- (6.2932,-1.1800) -- (6.2932,-0.4642) -- (6.1560,-0.4642) -- cycle;
\fill[color={rgb,255:red,35;green,127;blue,119},even odd rule] (6.8708,-1.0379) -- (6.7341,-1.0150) -- (6.7202,-1.0487) -- (6.6981,-1.0737) -- (6.6695,-1.0867) -- (6.6351,-1.0886) -- (6.6033,-1.0807) -- (6.5758,-1.0632) -- (6.5536,-1.0366) -- (6.5415,-1.0085) -- (6.5344,-0.9603) -- (6.8782,-0.9603) -- (6.8732,-0.8696) -- (6.8546,-0.7949) -- (6.8223,-0.7360) -- (6.7967,-0.7082) -- (6.7676,-0.6858) -- (6.7235,-0.6644) -- (6.6864,-0.6545) -- (6.6316,-0.6497) -- (6.5823,-0.6543) -- (6.5375,-0.6680) -- (6.4971,-0.6908) -- (6.4612,-0.7227) -- (6.4383,-0.7521) -- (6.4151,-0.7973) -- (6.4004,-0.8495) -- (6.3941,-0.9087) -- (6.3970,-0.9763) -- (6.4100,-1.0343) -- (6.4276,-1.0755) -- (6.4535,-1.1140) -- (6.4925,-1.1499) -- (6.5404,-1.1748) -- (6.5970,-1.1886) -- (6.6662,-1.1911) -- (6.7236,-1.1819) -- (6.7729,-1.1616) -- (6.8141,-1.1305) -- (6.8468,-1.0893) -- cycle (6.7419,-0.8763) -- (6.5369,-0.8763) -- (6.5438,-0.8260) -- (6.5621,-0.7916) -- (6.5943,-0.7651) -- (6.6343,-0.7548) -- (6.6742,-0.7606) -- (6.7112,-0.7857) -- (6.7335,-0.8239) -- cycle;
\fill[color={rgb,255:red,44;green,139;blue,128},even odd rule] (6.9797,-0.6614) -- (6.9797,-1.1800) -- (7.1169,-1.1800) -- (7.1171,-0.9100) -- (7.1254,-0.8298) -- (7.1382,-0.7991) -- (7.1582,-0.7763) -- (7.1877,-0.7601) -- (7.2214,-0.7547) -- (7.2476,-0.7581) -- (7.2669,-0.7683) -- (7.2817,-0.7879) -- (7.2891,-0.8125) -- (7.2931,-0.8777) -- (7.2932,-1.1800) -- (7.4304,-1.1800) -- (7.4306,-0.9133) -- (7.4390,-0.8329) -- (7.4521,-0.8006) -- (7.4721,-0.7769) -- (7.5011,-0.7603) -- (7.5325,-0.7547) -- (7.5592,-0.7588) -- (7.5803,-0.7712) -- (7.5928,-0.7870) -- (7.6007,-0.8141) -- (7.6052,-0.8742) -- (7.6052,-1.1800) -- (7.7424,-1.1800) -- (7.7389,-0.7845) -- (7.7316,-0.7500) -- (7.7202,-0.7250) -- (7.7000,-0.6974) -- (7.6788,-0.6791) -- (7.6483,-0.6628) -- (7.6129,-0.6530) -- (7.5725,-0.6497) -- (7.5342,-0.6537) -- (7.4980,-0.6654) -- (7.4686,-0.6818) -- (7.4143,-0.7322) -- (7.3838,-0.6922) -- (7.3449,-0.6654) -- (7.2977,-0.6517) -- (7.2448,-0.6510) -- (7.2011,-0.6613) -- (7.1705,-0.6758) -- (7.1326,-0.7042) -- (7.1062,-0.7322) -- (7.1062,-0.6614) -- cycle;
\fill[color={rgb,255:red,43;green,137;blue,127},even odd rule] (7.8572,-0.7972) -- (7.9817,-0.8196) -- (7.9984,-0.7859) -- (8.0203,-0.7657) -- (8.0507,-0.7561) -- (8.0911,-0.7553) -- (8.1262,-0.7613) -- (8.1458,-0.7716) -- (8.1591,-0.7906) -- (8.1646,-0.8139) -- (8.1653,-0.8416) -- (7.9439,-0.8984) -- (7.9008,-0.9203) -- (7.8725,-0.9477) -- (7.8483,-0.9968) -- (7.8437,-1.0546) -- (7.8576,-1.1058) -- (7.8735,-1.1312) -- (7.8951,-1.1535) -- (7.9285,-1.1747) -- (7.9601,-1.1856) -- (7.9959,-1.1910) -- (8.0378,-1.1907) -- (8.0753,-1.1836) -- (8.1108,-1.1701) -- (8.1443,-1.1500) -- (8.1760,-1.1234) -- (8.1941,-1.1800) -- (8.3298,-1.1800) -- (8.3092,-1.1269) -- (8.3012,-1.0810) -- (8.3000,-0.8385) -- (8.2955,-0.7741) -- (8.2817,-0.7266) -- (8.2527,-0.6927) -- (8.2065,-0.6664) -- (8.1440,-0.6528) -- (8.0573,-0.6503) -- (7.9838,-0.6612) -- (7.9430,-0.6774) -- (7.9057,-0.7064) -- (7.8772,-0.7463) -- cycle (8.1653,-0.9310) -- (8.1653,-0.9644) -- (8.1612,-1.0203) -- (8.1498,-1.0471) -- (8.1302,-1.0683) -- (8.0927,-1.0889) -- (8.0564,-1.0950) -- (8.0228,-1.0884) -- (7.9944,-1.0659) -- (7.9823,-1.0415) -- (7.9806,-1.0171) -- (7.9887,-0.9945) -- (8.0038,-0.9783) -- (8.0394,-0.9626) -- cycle;
\fill[color={rgb,255:red,58;green,156;blue,142},even odd rule] (2.7810,-2.3600) -- (2.9084,-2.3600) -- (2.9084,-1.6442) -- (2.7712,-1.6442) -- (2.7712,-1.9020) -- (2.7203,-1.8580) -- (2.6729,-1.8368) -- (2.6208,-1.8297) -- (2.5544,-1.8394) -- (2.4972,-1.8683) -- (2.4566,-1.9072) -- (2.4263,-1.9587) -- (2.4082,-2.0225) -- (2.4023,-2.1147) -- (2.4143,-2.2007) -- (2.4379,-2.2601) -- (2.4822,-2.3166) -- (2.5172,-2.3436) -- (2.5555,-2.3616) -- (2.5970,-2.3706) -- (2.6408,-2.3704) -- (2.6844,-2.3595) -- (2.7219,-2.3410) -- (2.7556,-2.3140) -- (2.7810,-2.2838) -- cycle (2.5423,-2.0997) -- (2.5455,-2.0450) -- (2.5578,-2.0015) -- (2.5794,-1.9688) -- (2.6071,-1.9470) -- (2.6450,-1.9353) -- (2.6867,-1.9385) -- (2.7223,-1.9568) -- (2.7505,-1.9905) -- (2.7637,-2.0258) -- (2.7697,-2.0611) -- (2.7716,-2.1125) -- (2.7671,-2.1571) -- (2.7463,-2.2123) -- (2.7256,-2.2367) -- (2.7018,-2.2532) -- (2.6696,-2.2627) -- (2.6365,-2.2614) -- (2.6050,-2.2499) -- (2.5834,-2.2331) -- (2.5652,-2.2096) -- (2.5531,-2.1819) -- cycle;
\fill[color={rgb,255:red,57;green,155;blue,141},even odd rule] (3.4807,-2.2179) -- (3.3440,-2.1950) -- (3.3301,-2.2287) -- (3.3080,-2.2537) -- (3.2793,-2.2667) -- (3.2450,-2.2686) -- (3.2131,-2.2607) -- (3.1857,-2.2432) -- (3.1635,-2.2166) -- (3.1514,-2.1885) -- (3.1443,-2.1403) -- (3.4880,-2.1403) -- (3.4830,-2.0496) -- (3.4644,-1.9749) -- (3.4322,-1.9160) -- (3.4066,-1.8882) -- (3.3775,-1.8658) -- (3.3334,-1.8444) -- (3.2963,-1.8345) -- (3.2415,-1.8297) -- (3.1922,-1.8343) -- (3.1473,-1.8480) -- (3.1070,-1.8708) -- (3.0710,-1.9027) -- (3.0481,-1.9321) -- (3.0250,-1.9773) -- (3.0102,-2.0295) -- (3.0039,-2.0887) -- (3.0069,-2.1563) -- (3.0199,-2.2143) -- (3.0375,-2.2555) -- (3.0633,-2.2940) -- (3.1024,-2.3299) -- (3.1502,-2.3548) -- (3.2069,-2.3686) -- (3.2760,-2.3711) -- (3.3334,-2.3619) -- (3.3828,-2.3416) -- (3.4240,-2.3105) -- (3.4566,-2.2693) -- cycle (3.3518,-2.0563) -- (3.1467,-2.0563) -- (3.1537,-2.0060) -- (3.1720,-1.9716) -- (3.2042,-1.9451) -- (3.2442,-1.9348) -- (3.2841,-1.9406) -- (3.3210,-1.9657) -- (3.3434,-2.0039) -- cycle;
\fill[color={rgb,255:red,28;green,102;blue,100},even odd rule] (4.0613,-2.3600) -- (4.2156,-2.3600) -- (4.4719,-1.6442) -- (4.3186,-1.6442) -- (4.1433,-2.1740) -- (3.9622,-1.6442) -- (3.8054,-1.6442) -- cycle;
\fill[color={rgb,255:red,29;green,105;blue,102},even odd rule] (4.5447,-1.7711) -- (4.6819,-1.7711) -- (4.6819,-1.6442) -- (4.5447,-1.6442) -- cycle;
\fill[color={rgb,255:red,29;green,108;blue,105},even odd rule] (4.5447,-2.3600) -- (4.6819,-2.3600) -- (4.6819,-1.8414) -- (4.5447,-1.8414) -- cycle;
\fill[color={rgb,255:red,29;green,111;blue,107},even odd rule] (4.7742,-2.2121) -- (4.7945,-2.2630) -- (4.8268,-2.3053) -- (4.8711,-2.3388) -- (4.9175,-2.3581) -- (4.9845,-2.3702) -- (5.0645,-2.3699) -- (5.1359,-2.3555) -- (5.1832,-2.3325) -- (5.2304,-2.2865) -- (5.2548,-2.2323) -- (5.2585,-2.1979) -- (5.2559,-2.1670) -- (5.2478,-2.1397) -- (5.2344,-2.1160) -- (5.2037,-2.0870) -- (5.1556,-2.0625) -- (4.9445,-2.0018) -- (4.9260,-1.9879) -- (4.9207,-1.9692) -- (4.9280,-1.9508) -- (4.9461,-1.9384) -- (4.9776,-1.9305) -- (5.0139,-1.9284) -- (5.0519,-1.9324) -- (5.0771,-1.9425) -- (5.0975,-1.9613) -- (5.1116,-1.9909) -- (5.2410,-1.9669) -- (5.2134,-1.9074) -- (5.1880,-1.8786) -- (5.1559,-1.8563) -- (5.1139,-1.8407) -- (5.0727,-1.8331) -- (5.0115,-1.8297) -- (4.9473,-1.8342) -- (4.8946,-1.8477) -- (4.8607,-1.8649) -- (4.8184,-1.9064) -- (4.7965,-1.9566) -- (4.7954,-2.0158) -- (4.8050,-2.0480) -- (4.8172,-2.0695) -- (4.8547,-2.1056) -- (4.8944,-2.1251) -- (5.1022,-2.1846) -- (5.1178,-2.2006) -- (5.1197,-2.2273) -- (5.1089,-2.2476) -- (5.0877,-2.2616) -- (5.0534,-2.2707) -- (5.0110,-2.2723) -- (4.9751,-2.2660) -- (4.9442,-2.2492) -- (4.9230,-2.2224) -- (4.9119,-2.1911) -- cycle;
\fill[color={rgb,255:red,30;green,114;blue,109},even odd rule] (5.3787,-1.7711) -- (5.5159,-1.7711) -- (5.5159,-1.6442) -- (5.3787,-1.6442) -- cycle;
\fill[color={rgb,255:red,30;green,117;blue,112},even odd rule] (5.3787,-2.3600) -- (5.5159,-2.3600) -- (5.5159,-1.8414) -- (5.3787,-1.8414) -- cycle;
\fill[color={rgb,255:red,31;green,120;blue,114},even odd rule] (5.8943,-1.8414) -- (5.8005,-1.8414) -- (5.8005,-1.6583) -- (5.6628,-1.7384) -- (5.6628,-1.8414) -- (5.5999,-1.8414) -- (5.5999,-1.9508) -- (5.6628,-1.9508) -- (5.6633,-2.2176) -- (5.6714,-2.2933) -- (5.6816,-2.3191) -- (5.6968,-2.3387) -- (5.7196,-2.3543) -- (5.7496,-2.3658) -- (5.7835,-2.3713) -- (5.8266,-2.3704) -- (5.8678,-2.3637) -- (5.9055,-2.3512) -- (5.8938,-2.2448) -- (5.8492,-2.2571) -- (5.8204,-2.2541) -- (5.8040,-2.2363) -- (5.8009,-2.2000) -- (5.8005,-1.9508) -- (5.8943,-1.9508) -- cycle;
\fill[color={rgb,255:red,32;green,123;blue,116},even odd rule] (5.9675,-1.9772) -- (6.0920,-1.9996) -- (6.1088,-1.9659) -- (6.1306,-1.9457) -- (6.1611,-1.9361) -- (6.2015,-1.9353) -- (6.2366,-1.9413) -- (6.2561,-1.9516) -- (6.2695,-1.9706) -- (6.2749,-1.9939) -- (6.2756,-2.0216) -- (6.0542,-2.0784) -- (6.0112,-2.1003) -- (5.9828,-2.1277) -- (5.9586,-2.1768) -- (5.9541,-2.2346) -- (5.9680,-2.2858) -- (5.9838,-2.3112) -- (6.0054,-2.3335) -- (6.0389,-2.3547) -- (6.0704,-2.3656) -- (6.1062,-2.3710) -- (6.1481,-2.3707) -- (6.1856,-2.3636) -- (6.2211,-2.3501) -- (6.2547,-2.3300) -- (6.2864,-2.3034) -- (6.3044,-2.3600) -- (6.4402,-2.3600) -- (6.4195,-2.3069) -- (6.4115,-2.2610) -- (6.4103,-2.0185) -- (6.4058,-1.9541) -- (6.3921,-1.9066) -- (6.3631,-1.8727) -- (6.3169,-1.8464) -- (6.2544,-1.8328) -- (6.1676,-1.8303) -- (6.0942,-1.8412) -- (6.0534,-1.8574) -- (6.0161,-1.8864) -- (5.9875,-1.9263) -- cycle (6.2756,-2.1110) -- (6.2756,-2.1444) -- (6.2715,-2.2003) -- (6.2602,-2.2271) -- (6.2406,-2.2483) -- (6.2030,-2.2689) -- (6.1667,-2.2750) -- (6.1331,-2.2684) -- (6.1048,-2.2459) -- (6.0927,-2.2215) -- (6.0910,-2.1971) -- (6.0991,-2.1745) -- (6.1142,-2.1583) -- (6.1497,-2.1426) -- cycle;
\fill[color={rgb,255:red,60;green,159;blue,144},even odd rule] (1.6985,-3.5400) -- (1.8259,-3.5400) -- (1.8259,-2.8242) -- (1.6887,-2.8242) -- (1.6887,-3.0820) -- (1.6378,-3.0380) -- (1.5904,-3.0168) -- (1.5383,-3.0097) -- (1.4719,-3.0194) -- (1.4146,-3.0483) -- (1.3741,-3.0872) -- (1.3438,-3.1387) -- (1.3256,-3.2025) -- (1.3198,-3.2947) -- (1.3318,-3.3807) -- (1.3554,-3.4401) -- (1.3996,-3.4966) -- (1.4347,-3.5236) -- (1.4730,-3.5416) -- (1.5144,-3.5506) -- (1.5583,-3.5504) -- (1.6019,-3.5395) -- (1.6393,-3.5210) -- (1.6731,-3.4940) -- (1.6985,-3.4638) -- cycle (1.4598,-3.2797) -- (1.4629,-3.2250) -- (1.4753,-3.1815) -- (1.4969,-3.1488) -- (1.5246,-3.1270) -- (1.5625,-3.1153) -- (1.6041,-3.1185) -- (1.6397,-3.1368) -- (1.6679,-3.1705) -- (1.6812,-3.2058) -- (1.6872,-3.2411) -- (1.6891,-3.2925) -- (1.6845,-3.3371) -- (1.6638,-3.3923) -- (1.6431,-3.4167) -- (1.6193,-3.4332) -- (1.5871,-3.4427) -- (1.5540,-3.4414) -- (1.5225,-3.4299) -- (1.5009,-3.4131) -- (1.4827,-3.3896) -- (1.4706,-3.3619) -- cycle;
\fill[color={rgb,255:red,59;green,158;blue,143},even odd rule] (2.3982,-3.3979) -- (2.2615,-3.3750) -- (2.2476,-3.4087) -- (2.2255,-3.4337) -- (2.1968,-3.4467) -- (2.1625,-3.4486) -- (2.1306,-3.4407) -- (2.1031,-3.4232) -- (2.0810,-3.3966) -- (2.0688,-3.3685) -- (2.0618,-3.3203) -- (2.4055,-3.3203) -- (2.4005,-3.2296) -- (2.3819,-3.1549) -- (2.3496,-3.0960) -- (2.3241,-3.0682) -- (2.2950,-3.0458) -- (2.2508,-3.0244) -- (2.2137,-3.0145) -- (2.1589,-3.0097) -- (2.1096,-3.0143) -- (2.0648,-3.0280) -- (2.0244,-3.0508) -- (1.9885,-3.0827) -- (1.9656,-3.1121) -- (1.9425,-3.1573) -- (1.9277,-3.2095) -- (1.9214,-3.2687) -- (1.9243,-3.3363) -- (1.9374,-3.3943) -- (1.9550,-3.4355) -- (1.9808,-3.4740) -- (2.0199,-3.5099) -- (2.0677,-3.5348) -- (2.1244,-3.5486) -- (2.1935,-3.5511) -- (2.2509,-3.5419) -- (2.3003,-3.5216) -- (2.3415,-3.4905) -- (2.3741,-3.4493) -- cycle (2.2693,-3.2363) -- (2.0642,-3.2363) -- (2.0712,-3.1860) -- (2.0895,-3.1516) -- (2.1216,-3.1251) -- (2.1617,-3.1148) -- (2.2015,-3.1206) -- (2.2385,-3.1457) -- (2.2609,-3.1839) -- cycle;
\fill[color={rgb,255:red,55;green,153;blue,139},even odd rule] (2.7961,-2.8242) -- (2.7961,-3.5400) -- (2.9407,-3.5400) -- (2.9407,-3.2700) -- (3.0811,-3.2693) -- (3.1638,-3.2632) -- (3.2132,-3.2511) -- (3.2598,-3.2258) -- (3.3043,-3.1814) -- (3.3332,-3.1251) -- (3.3441,-3.0605) -- (3.3386,-2.9847) -- (3.3211,-2.9339) -- (3.2923,-2.8921) -- (3.2500,-2.8563) -- (3.2076,-2.8371) -- (3.1660,-2.8293) -- (3.1015,-2.8252) -- cycle (2.9407,-3.1484) -- (2.9407,-2.9453) -- (3.1145,-2.9502) -- (3.1537,-2.9655) -- (3.1804,-2.9919) -- (3.1941,-3.0277) -- (3.1933,-3.0691) -- (3.1792,-3.1020) -- (3.1539,-3.1270) -- (3.1259,-3.1398) -- (3.0911,-3.1456) -- cycle;
\fill[color={rgb,255:red,54;green,152;blue,138},even odd rule] (3.4306,-3.2649) -- (3.4337,-3.3266) -- (3.4436,-3.3745) -- (3.4600,-3.4171) -- (3.4882,-3.4610) -- (3.5251,-3.4964) -- (3.5625,-3.5195) -- (3.6204,-3.5415) -- (3.6720,-3.5506) -- (3.7267,-3.5505) -- (3.7794,-3.5408) -- (3.8272,-3.5215) -- (3.8701,-3.4925) -- (3.9074,-3.4543) -- (3.9359,-3.4106) -- (3.9549,-3.3618) -- (3.9644,-3.3080) -- (3.9644,-3.2507) -- (3.9550,-3.1974) -- (3.9362,-3.1491) -- (3.9080,-3.1059) -- (3.8709,-3.0682) -- (3.8280,-3.0396) -- (3.7798,-3.0205) -- (3.7263,-3.0109) -- (3.6603,-3.0118) -- (3.6079,-3.0229) -- (3.5596,-3.0434) -- (3.5172,-3.0729) -- (3.4827,-3.1110) -- (3.4562,-3.1571) -- (3.4369,-3.2144) -- cycle (3.5712,-3.2902) -- (3.5762,-3.2279) -- (3.5916,-3.1853) -- (3.6220,-3.1486) -- (3.6601,-3.1273) -- (3.7112,-3.1222) -- (3.7529,-3.1345) -- (3.7925,-3.1678) -- (3.8154,-3.2124) -- (3.8222,-3.2439) -- (3.8243,-3.2894) -- (3.8130,-3.3559) -- (3.8004,-3.3822) -- (3.7835,-3.4039) -- (3.7584,-3.4239) -- (3.7358,-3.4341) -- (3.7046,-3.4397) -- (3.6722,-3.4373) -- (3.6430,-3.4269) -- (3.6220,-3.4128) -- (3.5991,-3.3881) -- (3.5854,-3.3630) -- (3.5746,-3.3254) -- cycle;
\fill[color={rgb,255:red,53;green,150;blue,137},even odd rule] (4.0730,-3.5400) -- (4.2102,-3.5400) -- (4.2102,-2.8242) -- (4.0730,-2.8242) -- cycle;
\fill[color={rgb,255:red,51;green,147;blue,135},even odd rule] (4.3508,-3.5400) -- (4.4880,-3.5400) -- (4.4880,-3.0214) -- (4.3508,-3.0214) -- cycle;
\fill[color={rgb,255:red,52;green,149;blue,136},even odd rule] (4.3406,-2.9580) -- (4.4275,-2.9580) -- (4.5618,-2.8120) -- (4.4080,-2.8120) -- cycle;
\fill[color={rgb,255:red,50;green,146;blue,134},even odd rule] (4.7727,-3.5932) -- (4.6160,-3.5742) -- (4.6209,-3.6318) -- (4.6465,-3.6812) -- (4.6707,-3.7053) -- (4.6940,-3.7206) -- (4.7570,-3.7418) -- (4.8577,-3.7504) -- (4.9552,-3.7408) -- (4.9948,-3.7278) -- (5.0273,-3.7104) -- (5.0528,-3.6885) -- (5.0704,-3.6645) -- (5.0938,-3.6053) -- (5.1028,-3.5298) -- (5.1038,-3.0214) -- (4.9753,-3.0214) -- (4.9753,-3.0942) -- (4.9330,-3.0497) -- (4.8841,-3.0216) -- (4.8285,-3.0101) -- (4.7719,-3.0140) -- (4.7310,-3.0269) -- (4.6939,-3.0485) -- (4.6606,-3.0786) -- (4.6332,-3.1168) -- (4.6136,-3.1628) -- (4.6040,-3.2023) -- (4.5981,-3.2617) -- (4.5996,-3.3156) -- (4.6071,-3.3623) -- (4.6205,-3.4046) -- (4.6400,-3.4428) -- (4.6702,-3.4817) -- (4.7075,-3.5120) -- (4.7390,-3.5276) -- (4.7736,-3.5369) -- (4.8229,-3.5397) -- (4.8667,-3.5318) -- (4.9066,-3.5134) -- (4.9423,-3.4846) -- (4.9666,-3.4560) -- (4.9636,-3.5812) -- (4.9535,-3.6084) -- (4.9372,-3.6265) -- (4.9178,-3.6364) -- (4.8888,-3.6432) -- (4.8199,-3.6427) -- (4.7858,-3.6264) -- cycle (4.7416,-3.2248) -- (4.7536,-3.1809) -- (4.7746,-3.1484) -- (4.8068,-3.1243) -- (4.8457,-3.1149) -- (4.8874,-3.1202) -- (4.9227,-3.1405) -- (4.9503,-3.1759) -- (4.9626,-3.2108) -- (4.9678,-3.2446) -- (4.9641,-3.3259) -- (4.9425,-3.3794) -- (4.9208,-3.4032) -- (4.8958,-3.4193) -- (4.8677,-3.4278) -- (4.8323,-3.4278) -- (4.8062,-3.4196) -- (4.7830,-3.4039) -- (4.7630,-3.3806) -- (4.7510,-3.3559) -- (4.7416,-3.3172) -- cycle;
\fill[color={rgb,255:red,49;green,144;blue,133},even odd rule] (5.2079,-3.2649) -- (5.2111,-3.3266) -- (5.2209,-3.3745) -- (5.2374,-3.4171) -- (5.2655,-3.4610) -- (5.3024,-3.4964) -- (5.3398,-3.5195) -- (5.3977,-3.5415) -- (5.4493,-3.5506) -- (5.5040,-3.5505) -- (5.5567,-3.5408) -- (5.6045,-3.5215) -- (5.6474,-3.4925) -- (5.6848,-3.4543) -- (5.7133,-3.4106) -- (5.7322,-3.3618) -- (5.7417,-3.3080) -- (5.7417,-3.2507) -- (5.7323,-3.1974) -- (5.7135,-3.1491) -- (5.6853,-3.1059) -- (5.6483,-3.0682) -- (5.6054,-3.0396) -- (5.5571,-3.0205) -- (5.5036,-3.0109) -- (5.4377,-3.0118) -- (5.3853,-3.0229) -- (5.3369,-3.0434) -- (5.2945,-3.0729) -- (5.2601,-3.1110) -- (5.2336,-3.1571) -- (5.2142,-3.2144) -- cycle (5.3485,-3.2902) -- (5.3535,-3.2279) -- (5.3690,-3.1853) -- (5.3993,-3.1486) -- (5.4375,-3.1273) -- (5.4885,-3.1222) -- (5.5303,-3.1345) -- (5.5698,-3.1678) -- (5.5927,-3.2124) -- (5.5995,-3.2439) -- (5.6017,-3.2894) -- (5.5903,-3.3559) -- (5.5778,-3.3822) -- (5.5608,-3.4039) -- (5.5357,-3.4239) -- (5.5132,-3.4341) -- (5.4820,-3.4397) -- (5.4496,-3.4373) -- (5.4203,-3.4269) -- (5.3993,-3.4128) -- (5.3764,-3.3881) -- (5.3627,-3.3630) -- (5.3520,-3.3254) -- cycle;
\fill[color={rgb,255:red,47;green,143;blue,131},even odd rule] (6.3220,-3.5400) -- (6.3201,-3.1646) -- (6.3086,-3.1060) -- (6.2825,-3.0616) -- (6.2524,-3.0372) -- (6.2114,-3.0189) -- (6.1699,-3.0107) -- (6.1225,-3.0111) -- (6.0754,-3.0221) -- (6.0429,-3.0375) -- (6.0035,-3.0677) -- (5.9768,-3.0976) -- (5.9768,-3.0214) -- (5.8494,-3.0214) -- (5.8494,-3.5400) -- (5.9866,-3.5400) -- (5.9869,-3.2745) -- (5.9958,-3.1902) -- (6.0097,-3.1602) -- (6.0324,-3.1367) -- (6.0660,-3.1196) -- (6.0999,-3.1147) -- (6.1286,-3.1190) -- (6.1522,-3.1320) -- (6.1695,-3.1524) -- (6.1790,-3.1780) -- (6.1847,-3.2553) -- (6.1848,-3.5400) -- cycle;
\fill[color={rgb,255:red,46;green,142;blue,130},even odd rule] (6.4296,-3.2649) -- (6.4327,-3.3266) -- (6.4426,-3.3745) -- (6.4591,-3.4171) -- (6.4872,-3.4610) -- (6.5241,-3.4964) -- (6.5615,-3.5195) -- (6.6194,-3.5415) -- (6.6710,-3.5506) -- (6.7257,-3.5505) -- (6.7784,-3.5408) -- (6.8262,-3.5215) -- (6.8691,-3.4925) -- (6.9065,-3.4543) -- (6.9349,-3.4106) -- (6.9539,-3.3618) -- (6.9634,-3.3080) -- (6.9634,-3.2507) -- (6.9540,-3.1974) -- (6.9352,-3.1491) -- (6.9070,-3.1059) -- (6.8699,-3.0682) -- (6.8270,-3.0396) -- (6.7788,-3.0205) -- (6.7253,-3.0109) -- (6.6593,-3.0118) -- (6.6069,-3.0229) -- (6.5586,-3.0434) -- (6.5162,-3.0729) -- (6.4817,-3.1110) -- (6.4552,-3.1571) -- (6.4359,-3.2144) -- cycle (6.5702,-3.2902) -- (6.5752,-3.2279) -- (6.5907,-3.1853) -- (6.6210,-3.1486) -- (6.6591,-3.1273) -- (6.7102,-3.1222) -- (6.7520,-3.1345) -- (6.7915,-3.1678) -- (6.8144,-3.2124) -- (6.8212,-3.2439) -- (6.8233,-3.2894) -- (6.8120,-3.3559) -- (6.7995,-3.3822) -- (6.7825,-3.4039) -- (6.7574,-3.4239) -- (6.7349,-3.4341) -- (6.7037,-3.4397) -- (6.6712,-3.4373) -- (6.6420,-3.4269) -- (6.6210,-3.4128) -- (6.5981,-3.3881) -- (6.5844,-3.3630) -- (6.5736,-3.3254) -- cycle;
\fill[color={rgb,255:red,45;green,140;blue,129},even odd rule] (7.0237,-3.3921) -- (7.0440,-3.4430) -- (7.0763,-3.4853) -- (7.1206,-3.5188) -- (7.1670,-3.5381) -- (7.2340,-3.5502) -- (7.3140,-3.5499) -- (7.3854,-3.5355) -- (7.4327,-3.5125) -- (7.4799,-3.4665) -- (7.5043,-3.4123) -- (7.5081,-3.3779) -- (7.5054,-3.3470) -- (7.4973,-3.3197) -- (7.4839,-3.2960) -- (7.4532,-3.2670) -- (7.4051,-3.2425) -- (7.1940,-3.1818) -- (7.1755,-3.1679) -- (7.1702,-3.1492) -- (7.1776,-3.1308) -- (7.1957,-3.1184) -- (7.2271,-3.1105) -- (7.2634,-3.1084) -- (7.3015,-3.1124) -- (7.3266,-3.1225) -- (7.3470,-3.1413) -- (7.3611,-3.1709) -- (7.4905,-3.1469) -- (7.4629,-3.0874) -- (7.4375,-3.0586) -- (7.4054,-3.0363) -- (7.3634,-3.0207) -- (7.3222,-3.0131) -- (7.2610,-3.0097) -- (7.1968,-3.0142) -- (7.1441,-3.0277) -- (7.1102,-3.0449) -- (7.0679,-3.0864) -- (7.0461,-3.1366) -- (7.0449,-3.1958) -- (7.0545,-3.2280) -- (7.0668,-3.2495) -- (7.1042,-3.2856) -- (7.1439,-3.3051) -- (7.3517,-3.3646) -- (7.3673,-3.3806) -- (7.3692,-3.4073) -- (7.3584,-3.4276) -- (7.3372,-3.4416) -- (7.3029,-3.4507) -- (7.2605,-3.4523) -- (7.2246,-3.4460) -- (7.1937,-3.4292) -- (7.1725,-3.4024) -- (7.1614,-3.3711) -- cycle;
\begin{scope}[transparency group,opacity=0.7]
\draw[color={rgb,255:red,245;green,154;blue,74},line width=2.2pt,line cap=butt] (-0.4500,-1.5173)--(-0.3817,-1.5566);
\draw[color={rgb,255:red,245;green,154;blue,74},line width=2.2pt,line cap=butt] (-0.3817,-1.5566)--(-0.3134,-1.5959);
\draw[color={rgb,255:red,245;green,154;blue,74},line width=2.2pt,line cap=butt] (-0.3134,-1.5959)--(-0.2450,-1.6352);
\draw[color={rgb,255:red,245;green,153;blue,74},line width=2.2pt,line cap=butt] (-0.2450,-1.6352)--(-0.1767,-1.6745);
\draw[color={rgb,255:red,245;green,153;blue,73},line width=2.2pt,line cap=butt] (-0.1767,-1.6745)--(-0.1084,-1.7138);
\draw[color={rgb,255:red,244;green,153;blue,73},line width=2.2pt,line cap=butt] (-0.1084,-1.7138)--(-0.0401,-1.7531);
\draw[color={rgb,255:red,244;green,153;blue,73},line width=2.2pt,line cap=butt] (-0.0401,-1.7531)--(0.0282,-1.7925);
\draw[color={rgb,255:red,244;green,153;blue,73},line width=2.2pt,line cap=butt] (0.0282,-1.7925)--(0.0965,-1.8318);
\draw[color={rgb,255:red,244;green,152;blue,73},line width=2.2pt,line cap=butt] (0.0965,-1.8318)--(0.1649,-1.8711);
\draw[color={rgb,255:red,244;green,152;blue,73},line width=2.2pt,line cap=butt] (0.1649,-1.8711)--(0.2332,-1.9104);
\draw[color={rgb,255:red,244;green,152;blue,73},line width=2.2pt,line cap=butt] (0.2332,-1.9104)--(0.3015,-1.9497);
\draw[color={rgb,255:red,244;green,152;blue,73},line width=2.2pt,line cap=butt] (0.3015,-1.9497)--(0.3698,-1.9890);
\draw[color={rgb,255:red,244;green,152;blue,72},line width=2.2pt,line cap=butt] (0.3698,-1.9890)--(0.4381,-2.0283);
\draw[color={rgb,255:red,244;green,152;blue,72},line width=2.2pt,line cap=butt] (0.4381,-2.0283)--(0.5064,-2.0676);
\draw[color={rgb,255:red,243;green,151;blue,72},line width=2.2pt,line cap=butt] (0.5064,-2.0676)--(0.5748,-2.1069);
\draw[color={rgb,255:red,243;green,151;blue,72},line width=2.2pt,line cap=butt] (0.5748,-2.1069)--(0.6431,-2.1462);
\draw[color={rgb,255:red,243;green,151;blue,72},line width=2.2pt,line cap=butt] (0.6431,-2.1462)--(0.7114,-2.1855);
\draw[color={rgb,255:red,243;green,151;blue,72},line width=2.2pt,line cap=butt] (0.7114,-2.1855)--(0.7797,-2.2248);
\draw[color={rgb,255:red,243;green,151;blue,72},line width=2.2pt,line cap=butt] (0.7797,-2.2248)--(0.8480,-2.2641);
\draw[color={rgb,255:red,243;green,151;blue,71},line width=2.2pt,line cap=butt] (0.8480,-2.2641)--(0.9164,-2.3034);
\draw[color={rgb,255:red,243;green,150;blue,71},line width=2.2pt,line cap=butt] (0.9164,-2.3034)--(0.9847,-2.3427);
\draw[color={rgb,255:red,243;green,150;blue,71},line width=2.2pt,line cap=butt] (0.9847,-2.3427)--(1.0530,-2.3820);
\draw[color={rgb,255:red,243;green,150;blue,71},line width=2.2pt,line cap=butt] (1.0530,-2.3820)--(1.1213,-2.4213);
\draw[color={rgb,255:red,242;green,150;blue,71},line width=2.2pt,line cap=butt] (1.1213,-2.4213)--(1.1896,-2.4606);
\draw[color={rgb,255:red,242;green,150;blue,71},line width=2.2pt,line cap=butt] (1.1896,-2.4606)--(1.2579,-2.4999);
\draw[color={rgb,255:red,242;green,149;blue,71},line width=2.2pt,line cap=butt] (1.2579,-2.4999)--(1.3263,-2.5393);
\draw[color={rgb,255:red,242;green,149;blue,71},line width=2.2pt,line cap=butt] (1.3263,-2.5393)--(1.3946,-2.5786);
\draw[color={rgb,255:red,242;green,149;blue,70},line width=2.2pt,line cap=butt] (1.3946,-2.5786)--(1.4629,-2.6179);
\draw[color={rgb,255:red,242;green,149;blue,70},line width=2.2pt,line cap=butt] (1.4629,-2.6179)--(1.5312,-2.6572);
\draw[color={rgb,255:red,242;green,149;blue,70},line width=2.2pt,line cap=butt] (1.5312,-2.6572)--(1.5995,-2.6965);
\draw[color={rgb,255:red,242;green,149;blue,70},line width=2.2pt,line cap=butt] (1.5995,-2.6965)--(1.6678,-2.7358);
\draw[color={rgb,255:red,242;green,148;blue,70},line width=2.2pt,line cap=butt] (1.6678,-2.7358)--(1.7362,-2.7751);
\draw[color={rgb,255:red,242;green,148;blue,70},line width=2.2pt,line cap=butt] (1.7362,-2.7751)--(1.8045,-2.8144);
\draw[color={rgb,255:red,241;green,148;blue,70},line width=2.2pt,line cap=butt] (1.8045,-2.8144)--(1.8728,-2.8537);
\draw[color={rgb,255:red,241;green,148;blue,70},line width=2.2pt,line cap=butt] (1.8728,-2.8537)--(1.9411,-2.8930);
\draw[color={rgb,255:red,241;green,148;blue,69},line width=2.2pt,line cap=butt] (1.9411,-2.8930)--(2.0094,-2.9323);
\draw[color={rgb,255:red,241;green,147;blue,69},line width=2.2pt,line cap=butt] (2.0094,-2.9323)--(2.0778,-2.9716);
\draw[color={rgb,255:red,241;green,147;blue,69},line width=2.2pt,line cap=butt] (2.0778,-2.9716)--(2.1461,-3.0109);
\draw[color={rgb,255:red,241;green,147;blue,69},line width=2.2pt,line cap=butt] (2.1461,-3.0109)--(2.2047,-2.9608);
\draw[color={rgb,255:red,241;green,147;blue,69},line width=2.2pt,line cap=butt] (2.2047,-2.9608)--(2.2634,-2.9108);
\draw[color={rgb,255:red,241;green,147;blue,69},line width=2.2pt,line cap=butt] (2.2634,-2.9108)--(2.3220,-2.8607);
\draw[color={rgb,255:red,241;green,147;blue,69},line width=2.2pt,line cap=butt] (2.3220,-2.8607)--(2.3806,-2.8106);
\draw[color={rgb,255:red,240;green,146;blue,69},line width=2.2pt,line cap=butt] (2.3806,-2.8106)--(2.4393,-2.7606);
\draw[color={rgb,255:red,240;green,146;blue,68},line width=2.2pt,line cap=butt] (2.4393,-2.7606)--(2.4979,-2.7105);
\draw[color={rgb,255:red,240;green,146;blue,68},line width=2.2pt,line cap=butt] (2.4979,-2.7105)--(2.5566,-2.6604);
\draw[color={rgb,255:red,240;green,146;blue,68},line width=2.2pt,line cap=butt] (2.5566,-2.6604)--(2.6152,-2.6104);
\draw[color={rgb,255:red,240;green,146;blue,68},line width=2.2pt,line cap=butt] (2.6152,-2.6104)--(2.6739,-2.5603);
\draw[color={rgb,255:red,240;green,146;blue,68},line width=2.2pt,line cap=butt] (2.6739,-2.5603)--(2.7325,-2.5102);
\draw[color={rgb,255:red,240;green,145;blue,68},line width=2.2pt,line cap=butt] (2.7325,-2.5102)--(2.7912,-2.4601);
\draw[color={rgb,255:red,240;green,145;blue,68},line width=2.2pt,line cap=butt] (2.7912,-2.4601)--(2.8498,-2.4101);
\draw[color={rgb,255:red,240;green,145;blue,68},line width=2.2pt,line cap=butt] (2.8498,-2.4101)--(2.9084,-2.3600);
\draw[color={rgb,255:red,240;green,145;blue,67},line width=2.2pt,line cap=butt] (2.9084,-2.3600)--(2.9731,-2.3500);
\draw[color={rgb,255:red,239;green,145;blue,67},line width=2.2pt,line cap=butt] (2.9731,-2.3500)--(3.0377,-2.3400);
\draw[color={rgb,255:red,239;green,145;blue,67},line width=2.2pt,line cap=butt] (3.0377,-2.3400)--(3.1024,-2.3299);
\draw[color={rgb,255:red,239;green,144;blue,67},line width=2.2pt,line cap=butt] (3.1024,-2.3299)--(3.1261,-2.4002);
\draw[color={rgb,255:red,239;green,144;blue,67},line width=2.2pt,line cap=butt] (3.1261,-2.4002)--(3.1499,-2.4705);
\draw[color={rgb,255:red,239;green,144;blue,67},line width=2.2pt,line cap=butt] (3.1499,-2.4705)--(3.1736,-2.5408);
\draw[color={rgb,255:red,239;green,144;blue,67},line width=2.2pt,line cap=butt] (3.1736,-2.5408)--(3.1973,-2.6110);
\draw[color={rgb,255:red,239;green,144;blue,67},line width=2.2pt,line cap=butt] (3.1973,-2.6110)--(3.2211,-2.6813);
\draw[color={rgb,255:red,239;green,144;blue,66},line width=2.2pt,line cap=butt] (3.2211,-2.6813)--(3.2448,-2.7516);
\draw[color={rgb,255:red,239;green,143;blue,66},line width=2.2pt,line cap=butt] (3.2448,-2.7516)--(3.2686,-2.8218);
\draw[color={rgb,255:red,239;green,143;blue,66},line width=2.2pt,line cap=butt] (3.2686,-2.8218)--(3.2923,-2.8921);
\draw[color={rgb,255:red,238;green,143;blue,66},line width=2.2pt,line cap=butt] (3.2923,-2.8921)--(3.3646,-2.9119);
\draw[color={rgb,255:red,238;green,143;blue,66},line width=2.2pt,line cap=butt] (3.3646,-2.9119)--(3.4369,-2.9317);
\draw[color={rgb,255:red,238;green,143;blue,66},line width=2.2pt,line cap=butt] (3.4369,-2.9317)--(3.5093,-2.9515);
\draw[color={rgb,255:red,238;green,143;blue,66},line width=2.2pt,line cap=butt] (3.5093,-2.9515)--(3.5816,-2.9713);
\draw[color={rgb,255:red,238;green,142;blue,66},line width=2.2pt,line cap=butt] (3.5816,-2.9713)--(3.6539,-2.9911);
\draw[color={rgb,255:red,238;green,142;blue,66},line width=2.2pt,line cap=butt] (3.6539,-2.9911)--(3.7263,-3.0109);
\draw[color={rgb,255:red,238;green,142;blue,65},line width=2.2pt,line cap=butt] (3.7263,-3.0109)--(3.8031,-3.0043);
\draw[color={rgb,255:red,238;green,142;blue,65},line width=2.2pt,line cap=butt] (3.8031,-3.0043)--(3.8798,-2.9977);
\draw[color={rgb,255:red,238;green,142;blue,65},line width=2.2pt,line cap=butt] (3.8798,-2.9977)--(3.9566,-2.9911);
\draw[color={rgb,255:red,238;green,141;blue,65},line width=2.2pt,line cap=butt] (3.9566,-2.9911)--(4.0334,-2.9844);
\draw[color={rgb,255:red,237;green,141;blue,65},line width=2.2pt,line cap=butt] (4.0334,-2.9844)--(4.1102,-2.9778);
\draw[color={rgb,255:red,237;green,141;blue,65},line width=2.2pt,line cap=butt] (4.1102,-2.9778)--(4.1870,-2.9712);
\draw[color={rgb,255:red,237;green,141;blue,65},line width=2.2pt,line cap=butt] (4.1870,-2.9712)--(4.2638,-2.9646);
\draw[color={rgb,255:red,237;green,141;blue,64},line width=2.2pt,line cap=butt] (4.2638,-2.9646)--(4.3406,-2.9580);
\draw[color={rgb,255:red,237;green,141;blue,64},line width=2.2pt,line cap=butt] (4.3406,-2.9580)--(4.3897,-2.9791);
\draw[color={rgb,255:red,237;green,141;blue,64},line width=2.2pt,line cap=butt] (4.3897,-2.9791)--(4.4389,-3.0003);
\draw[color={rgb,255:red,237;green,140;blue,64},line width=2.2pt,line cap=butt] (4.4389,-3.0003)--(4.4880,-3.0214);
\draw[color={rgb,255:red,237;green,140;blue,64},line width=2.2pt,line cap=butt] (4.4880,-3.0214)--(4.5675,-3.0222);
\draw[color={rgb,255:red,237;green,140;blue,64},line width=2.2pt,line cap=butt] (4.5675,-3.0222)--(4.6470,-3.0229);
\draw[color={rgb,255:red,237;green,140;blue,64},line width=2.2pt,line cap=butt] (4.6470,-3.0229)--(4.7264,-3.0236);
\draw[color={rgb,255:red,237;green,140;blue,64},line width=2.2pt,line cap=butt] (4.7264,-3.0236)--(4.8059,-3.0243);
\draw[color={rgb,255:red,236;green,140;blue,64},line width=2.2pt,line cap=butt] (4.8059,-3.0243)--(4.8853,-3.0250);
\draw[color={rgb,255:red,236;green,139;blue,63},line width=2.2pt,line cap=butt] (4.8853,-3.0250)--(4.9648,-3.0257);
\draw[color={rgb,255:red,236;green,139;blue,63},line width=2.2pt,line cap=butt] (4.9648,-3.0257)--(5.0442,-3.0264);
\draw[color={rgb,255:red,236;green,139;blue,63},line width=2.2pt,line cap=butt] (5.0442,-3.0264)--(5.1237,-3.0271);
\draw[color={rgb,255:red,236;green,139;blue,63},line width=2.2pt,line cap=butt] (5.1237,-3.0271)--(5.2032,-3.0278);
\draw[color={rgb,255:red,236;green,139;blue,63},line width=2.2pt,line cap=butt] (5.2032,-3.0278)--(5.2826,-3.0285);
\draw[color={rgb,255:red,236;green,138;blue,63},line width=2.2pt,line cap=butt] (5.2826,-3.0285)--(5.3621,-3.0293);
\draw[color={rgb,255:red,236;green,138;blue,63},line width=2.2pt,line cap=butt] (5.3621,-3.0293)--(5.4415,-3.0300);
\draw[color={rgb,255:red,236;green,138;blue,63},line width=2.2pt,line cap=butt] (5.4415,-3.0300)--(5.5210,-3.0307);
\draw[color={rgb,255:red,235;green,138;blue,62},line width=2.2pt,line cap=butt] (5.5210,-3.0307)--(5.6005,-3.0314);
\draw[color={rgb,255:red,235;green,138;blue,62},line width=2.2pt,line cap=butt] (5.6005,-3.0314)--(5.6799,-3.0321);
\draw[color={rgb,255:red,235;green,138;blue,62},line width=2.2pt,line cap=butt] (5.6799,-3.0321)--(5.7594,-3.0328);
\draw[color={rgb,255:red,235;green,137;blue,62},line width=2.2pt,line cap=butt] (5.7594,-3.0328)--(5.8388,-3.0335);
\draw[color={rgb,255:red,235;green,137;blue,62},line width=2.2pt,line cap=butt] (5.8388,-3.0335)--(5.9183,-3.0342);
\draw[color={rgb,255:red,235;green,137;blue,62},line width=2.2pt,line cap=butt] (5.9183,-3.0342)--(5.9977,-3.0349);
\draw[color={rgb,255:red,235;green,137;blue,62},line width=2.2pt,line cap=butt] (5.9977,-3.0349)--(6.0772,-3.0356);
\draw[color={rgb,255:red,235;green,137;blue,62},line width=2.2pt,line cap=butt] (6.0772,-3.0356)--(6.1567,-3.0364);
\draw[color={rgb,255:red,235;green,136;blue,61},line width=2.2pt,line cap=butt] (6.1567,-3.0364)--(6.2361,-3.0371);
\draw[color={rgb,255:red,234;green,136;blue,61},line width=2.2pt,line cap=butt] (6.2361,-3.0371)--(6.3156,-3.0378);
\draw[color={rgb,255:red,234;green,136;blue,61},line width=2.2pt,line cap=butt] (6.3156,-3.0378)--(6.3950,-3.0385);
\draw[color={rgb,255:red,234;green,136;blue,61},line width=2.2pt,line cap=butt] (6.3950,-3.0385)--(6.4745,-3.0392);
\draw[color={rgb,255:red,234;green,136;blue,61},line width=2.2pt,line cap=butt] (6.4745,-3.0392)--(6.5540,-3.0399);
\draw[color={rgb,255:red,234;green,136;blue,61},line width=2.2pt,line cap=butt] (6.5540,-3.0399)--(6.6334,-3.0406);
\draw[color={rgb,255:red,234;green,135;blue,61},line width=2.2pt,line cap=butt] (6.6334,-3.0406)--(6.7129,-3.0413);
\draw[color={rgb,255:red,234;green,135;blue,60},line width=2.2pt,line cap=butt] (6.7129,-3.0413)--(6.7923,-3.0420);
\draw[color={rgb,255:red,234;green,135;blue,60},line width=2.2pt,line cap=butt] (6.7923,-3.0420)--(6.8718,-3.0427);
\draw[color={rgb,255:red,234;green,135;blue,60},line width=2.2pt,line cap=butt] (6.8718,-3.0427)--(6.9513,-3.0434);
\draw[color={rgb,255:red,233;green,135;blue,60},line width=2.2pt,line cap=butt] (6.9513,-3.0434)--(7.0307,-3.0442);
\draw[color={rgb,255:red,233;green,134;blue,60},line width=2.2pt,line cap=butt] (7.0307,-3.0442)--(7.1102,-3.0449);
\draw[color={rgb,255:red,233;green,134;blue,60},line width=2.2pt,line cap=butt] (7.1102,-3.0449)--(7.1355,-2.9703);
\draw[color={rgb,255:red,233;green,134;blue,60},line width=2.2pt,line cap=butt] (7.1355,-2.9703)--(7.1607,-2.8957);
\draw[color={rgb,255:red,233;green,134;blue,60},line width=2.2pt,line cap=butt] (7.1607,-2.8957)--(7.1860,-2.8211);
\draw[color={rgb,255:red,233;green,134;blue,59},line width=2.2pt,line cap=butt] (7.1860,-2.8211)--(7.2113,-2.7465);
\draw[color={rgb,255:red,233;green,134;blue,59},line width=2.2pt,line cap=butt] (7.2113,-2.7465)--(7.2366,-2.6719);
\draw[color={rgb,255:red,233;green,133;blue,59},line width=2.2pt,line cap=butt] (7.2366,-2.6719)--(7.2619,-2.5973);
\draw[color={rgb,255:red,233;green,133;blue,59},line width=2.2pt,line cap=butt] (7.2619,-2.5973)--(7.2872,-2.5227);
\draw[color={rgb,255:red,233;green,133;blue,59},line width=2.2pt,line cap=butt] (7.2872,-2.5227)--(7.3125,-2.4481);
\draw[color={rgb,255:red,232;green,133;blue,59},line width=2.2pt,line cap=butt] (7.3125,-2.4481)--(7.3378,-2.3735);
\draw[color={rgb,255:red,232;green,133;blue,59},line width=2.2pt,line cap=butt] (7.3378,-2.3735)--(7.3631,-2.2989);
\draw[color={rgb,255:red,232;green,132;blue,59},line width=2.2pt,line cap=butt] (7.3631,-2.2989)--(7.3884,-2.2243);
\draw[color={rgb,255:red,232;green,132;blue,58},line width=2.2pt,line cap=butt] (7.3884,-2.2243)--(7.4137,-2.1497);
\draw[color={rgb,255:red,232;green,132;blue,58},line width=2.2pt,line cap=butt] (7.4137,-2.1497)--(7.4389,-2.0751);
\draw[color={rgb,255:red,232;green,132;blue,58},line width=2.2pt,line cap=butt] (7.4389,-2.0751)--(7.4642,-2.0005);
\draw[color={rgb,255:red,232;green,132;blue,58},line width=2.2pt,line cap=butt] (7.4642,-2.0005)--(7.4895,-1.9259);
\draw[color={rgb,255:red,232;green,132;blue,58},line width=2.2pt,line cap=butt] (7.4895,-1.9259)--(7.5148,-1.8514);
\draw[color={rgb,255:red,232;green,131;blue,58},line width=2.2pt,line cap=butt] (7.5148,-1.8514)--(7.5401,-1.7768);
\draw[color={rgb,255:red,231;green,131;blue,58},line width=2.2pt,line cap=butt] (7.5401,-1.7768)--(7.5654,-1.7022);
\draw[color={rgb,255:red,231;green,131;blue,58},line width=2.2pt,line cap=butt] (7.5654,-1.7022)--(7.5907,-1.6276);
\draw[color={rgb,255:red,231;green,131;blue,57},line width=2.2pt,line cap=butt] (7.5907,-1.6276)--(7.6160,-1.5530);
\draw[color={rgb,255:red,231;green,131;blue,57},line width=2.2pt,line cap=butt] (7.6160,-1.5530)--(7.6413,-1.4784);
\draw[color={rgb,255:red,231;green,131;blue,57},line width=2.2pt,line cap=butt] (7.6413,-1.4784)--(7.6666,-1.4038);
\draw[color={rgb,255:red,231;green,130;blue,57},line width=2.2pt,line cap=butt] (7.6666,-1.4038)--(7.6919,-1.3292);
\draw[color={rgb,255:red,231;green,130;blue,57},line width=2.2pt,line cap=butt] (7.6919,-1.3292)--(7.7171,-1.2546);
\draw[color={rgb,255:red,231;green,130;blue,57},line width=2.2pt,line cap=butt] (7.7171,-1.2546)--(7.7424,-1.1800);
\draw[color={rgb,255:red,231;green,130;blue,57},line width=2.2pt,line cap=butt] (7.7424,-1.1800)--(7.7747,-1.1091);
\draw[color={rgb,255:red,231;green,130;blue,56},line width=2.2pt,line cap=butt] (7.7747,-1.1091)--(7.8069,-1.0381);
\draw[color={rgb,255:red,230;green,129;blue,56},line width=2.2pt,line cap=butt] (7.8069,-1.0381)--(7.8391,-0.9672);
\draw[color={rgb,255:red,230;green,129;blue,56},line width=2.2pt,line cap=butt] (7.8391,-0.9672)--(7.8713,-0.8962);
\draw[color={rgb,255:red,230;green,129;blue,56},line width=2.2pt,line cap=butt] (7.8713,-0.8962)--(7.9036,-0.8253);
\draw[color={rgb,255:red,230;green,129;blue,56},line width=2.2pt,line cap=butt] (7.9036,-0.8253)--(7.9358,-0.7543);
\draw[color={rgb,255:red,230;green,129;blue,56},line width=2.2pt,line cap=butt] (7.9358,-0.7543)--(7.9680,-0.6834);
\draw[color={rgb,255:red,230;green,129;blue,56},line width=2.2pt,line cap=butt] (7.9680,-0.6834)--(8.0003,-0.6125);
\draw[color={rgb,255:red,230;green,128;blue,56},line width=2.2pt,line cap=butt] (8.0003,-0.6125)--(8.0325,-0.5415);
\draw[color={rgb,255:red,230;green,128;blue,55},line width=2.2pt,line cap=butt] (8.0325,-0.5415)--(8.0647,-0.4706);
\draw[color={rgb,255:red,230;green,128;blue,55},line width=2.2pt,line cap=butt] (8.0647,-0.4706)--(8.0969,-0.3996);
\draw[color={rgb,255:red,229;green,128;blue,55},line width=2.2pt,line cap=butt] (8.0969,-0.3996)--(8.1292,-0.3287);
\draw[color={rgb,255:red,229;green,128;blue,55},line width=2.2pt,line cap=butt] (8.1292,-0.3287)--(8.1614,-0.2577);
\draw[color={rgb,255:red,229;green,127;blue,55},line width=2.2pt,line cap=butt] (8.1614,-0.2577)--(8.1936,-0.1868);
\draw[color={rgb,255:red,229;green,127;blue,55},line width=2.2pt,line cap=butt] (8.1936,-0.1868)--(8.2259,-0.1159);
\draw[color={rgb,255:red,229;green,127;blue,55},line width=2.2pt,line cap=butt] (8.2259,-0.1159)--(8.2581,-0.0449);
\draw[color={rgb,255:red,229;green,127;blue,55},line width=2.2pt,line cap=butt] (8.2581,-0.0449)--(8.2903,0.0260);
\draw[color={rgb,255:red,229;green,127;blue,54},line width=2.2pt,line cap=butt] (8.2903,0.0260)--(8.3225,0.0970);
\draw[color={rgb,255:red,229;green,127;blue,54},line width=2.2pt,line cap=butt] (8.3225,0.0970)--(8.2429,0.0934);
\draw[color={rgb,255:red,229;green,126;blue,54},line width=2.2pt,line cap=butt] (8.2429,0.0934)--(8.1633,0.0898);
\draw[color={rgb,255:red,229;green,126;blue,54},line width=2.2pt,line cap=butt] (8.1633,0.0898)--(8.0837,0.0862);
\draw[color={rgb,255:red,228;green,126;blue,54},line width=2.2pt,line cap=butt] (8.0837,0.0862)--(8.0041,0.0826);
\draw[color={rgb,255:red,228;green,126;blue,54},line width=2.2pt,line cap=butt] (8.0041,0.0826)--(7.9244,0.0790);
\draw[color={rgb,255:red,228;green,126;blue,54},line width=2.2pt,line cap=butt] (7.9244,0.0790)--(7.8448,0.0754);
\draw[color={rgb,255:red,228;green,126;blue,54},line width=2.2pt,line cap=butt] (7.8448,0.0754)--(7.7652,0.0718);
\draw[color={rgb,255:red,228;green,125;blue,53},line width=2.2pt,line cap=butt] (7.7652,0.0718)--(7.6856,0.0682);
\draw[color={rgb,255:red,228;green,125;blue,53},line width=2.2pt,line cap=butt] (7.6856,0.0682)--(7.6060,0.0647);
\draw[color={rgb,255:red,228;green,125;blue,53},line width=2.2pt,line cap=butt] (7.6060,0.0647)--(7.5264,0.0611);
\draw[color={rgb,255:red,228;green,125;blue,53},line width=2.2pt,line cap=butt] (7.5264,0.0611)--(7.4467,0.0575);
\draw[color={rgb,255:red,228;green,125;blue,53},line width=2.2pt,line cap=butt] (7.4467,0.0575)--(7.3671,0.0539);
\draw[color={rgb,255:red,227;green,124;blue,53},line width=2.2pt,line cap=butt] (7.3671,0.0539)--(7.2875,0.0503);
\draw[color={rgb,255:red,227;green,124;blue,53},line width=2.2pt,line cap=butt] (7.2875,0.0503)--(7.2079,0.0467);
\draw[color={rgb,255:red,227;green,124;blue,52},line width=2.2pt,line cap=butt] (7.2079,0.0467)--(7.1283,0.0431);
\draw[color={rgb,255:red,227;green,124;blue,52},line width=2.2pt,line cap=butt] (7.1283,0.0431)--(7.0486,0.0395);
\draw[color={rgb,255:red,227;green,124;blue,52},line width=2.2pt,line cap=butt] (7.0486,0.0395)--(6.9690,0.0359);
\draw[color={rgb,255:red,227;green,124;blue,52},line width=2.2pt,line cap=butt] (6.9690,0.0359)--(6.8894,0.0323);
\draw[color={rgb,255:red,227;green,123;blue,52},line width=2.2pt,line cap=butt] (6.8894,0.0323)--(6.8098,0.0287);
\draw[color={rgb,255:red,227;green,123;blue,52},line width=2.2pt,line cap=butt] (6.8098,0.0287)--(6.7302,0.0251);
\draw[color={rgb,255:red,227;green,123;blue,52},line width=2.2pt,line cap=butt] (6.7302,0.0251)--(6.6506,0.0216);
\draw[color={rgb,255:red,226;green,123;blue,52},line width=2.2pt,line cap=butt] (6.6506,0.0216)--(6.5709,0.0180);
\draw[color={rgb,255:red,226;green,123;blue,51},line width=2.2pt,line cap=butt] (6.5709,0.0180)--(6.4913,0.0144);
\draw[color={rgb,255:red,226;green,122;blue,51},line width=2.2pt,line cap=butt] (6.4913,0.0144)--(6.4117,0.0108);
\draw[color={rgb,255:red,226;green,122;blue,51},line width=2.2pt,line cap=butt] (6.4117,0.0108)--(6.3321,0.0072);
\draw[color={rgb,255:red,226;green,122;blue,51},line width=2.2pt,line cap=butt] (6.3321,0.0072)--(6.2525,0.0036);
\draw[color={rgb,255:red,226;green,122;blue,51},line width=2.2pt,line cap=butt] (6.2525,0.0036)--(6.1729,0.0000);
\draw[color={rgb,255:red,226;green,122;blue,51},line width=2.2pt,line cap=butt] (6.1729,0.0000)--(6.1949,-0.0725);
\draw[color={rgb,255:red,226;green,121;blue,51},line width=2.2pt,line cap=butt] (6.1949,-0.0725)--(6.2169,-0.1450);
\draw[color={rgb,255:red,225;green,121;blue,51},line width=2.2pt,line cap=butt] (6.2169,-0.1450)--(6.2389,-0.2174);
\draw[color={rgb,255:red,225;green,121;blue,50},line width=2.2pt,line cap=butt] (6.2389,-0.2174)--(6.2609,-0.2899);
\draw[color={rgb,255:red,225;green,121;blue,50},line width=2.2pt,line cap=butt] (6.2609,-0.2899)--(6.2830,-0.3624);
\draw[color={rgb,255:red,225;green,120;blue,50},line width=2.2pt,line cap=butt] (6.2830,-0.3624)--(6.3050,-0.4349);
\draw[color={rgb,255:red,225;green,120;blue,50},line width=2.2pt,line cap=butt] (6.3050,-0.4349)--(6.3270,-0.5074);
\draw[color={rgb,255:red,225;green,120;blue,50},line width=2.2pt,line cap=butt] (6.3270,-0.5074)--(6.3490,-0.5799);
\draw[color={rgb,255:red,224;green,120;blue,50},line width=2.2pt,line cap=butt] (6.3490,-0.5799)--(6.3711,-0.6523);
\draw[color={rgb,255:red,224;green,120;blue,50},line width=2.2pt,line cap=butt] (6.3711,-0.6523)--(6.3931,-0.7248);
\draw[color={rgb,255:red,224;green,119;blue,49},line width=2.2pt,line cap=butt] (6.3931,-0.7248)--(6.4151,-0.7973);
\draw[color={rgb,255:red,224;green,119;blue,49},line width=2.2pt,line cap=butt] (6.4151,-0.7973)--(6.3494,-0.8407);
\draw[color={rgb,255:red,224;green,119;blue,49},line width=2.2pt,line cap=butt] (6.3494,-0.8407)--(6.2837,-0.8840);
\draw[color={rgb,255:red,223;green,119;blue,49},line width=2.2pt,line cap=butt] (6.2837,-0.8840)--(6.2180,-0.9274);
\draw[color={rgb,255:red,223;green,118;blue,49},line width=2.2pt,line cap=butt] (6.2180,-0.9274)--(6.1523,-0.9708);
\draw[color={rgb,255:red,223;green,118;blue,49},line width=2.2pt,line cap=butt] (6.1523,-0.9708)--(6.0866,-1.0142);
\draw[color={rgb,255:red,223;green,118;blue,49},line width=2.2pt,line cap=butt] (6.0866,-1.0142)--(6.0210,-1.0575);
\draw[color={rgb,255:red,223;green,118;blue,48},line width=2.2pt,line cap=butt] (6.0210,-1.0575)--(6.0205,-1.1329);
\draw[color={rgb,255:red,223;green,118;blue,48},line width=2.2pt,line cap=butt] (6.0205,-1.1329)--(6.0201,-1.2082);
\draw[color={rgb,255:red,222;green,117;blue,48},line width=2.2pt,line cap=butt] (6.0201,-1.2082)--(6.0196,-1.2836);
\draw[color={rgb,255:red,222;green,117;blue,48},line width=2.2pt,line cap=butt] (6.0196,-1.2836)--(6.0192,-1.3589);
\draw[color={rgb,255:red,222;green,117;blue,48},line width=2.2pt,line cap=butt] (6.0192,-1.3589)--(6.0187,-1.4343);
\draw[color={rgb,255:red,222;green,117;blue,48},line width=2.2pt,line cap=butt] (6.0187,-1.4343)--(6.0183,-1.5096);
\draw[color={rgb,255:red,222;green,116;blue,48},line width=2.2pt,line cap=butt] (6.0183,-1.5096)--(6.0178,-1.5850);
\draw[color={rgb,255:red,222;green,116;blue,48},line width=2.2pt,line cap=butt] (6.0178,-1.5850)--(6.0174,-1.6603);
\draw[color={rgb,255:red,221;green,116;blue,47},line width=2.2pt,line cap=butt] (6.0174,-1.6603)--(6.0169,-1.7357);
\draw[color={rgb,255:red,221;green,116;blue,47},line width=2.2pt,line cap=butt] (6.0169,-1.7357)--(6.0165,-1.8110);
\draw[color={rgb,255:red,221;green,115;blue,47},line width=2.2pt,line cap=butt] (6.0165,-1.8110)--(6.0161,-1.8864);
\draw[color={rgb,255:red,221;green,115;blue,47},line width=2.2pt,line cap=butt] (6.0161,-1.8864)--(5.9446,-1.8799);
\draw[color={rgb,255:red,221;green,115;blue,47},line width=2.2pt,line cap=butt] (5.9446,-1.8799)--(5.8731,-1.8735);
\draw[color={rgb,255:red,221;green,115;blue,47},line width=2.2pt,line cap=butt] (5.8731,-1.8735)--(5.8017,-1.8671);
\draw[color={rgb,255:red,220;green,115;blue,47},line width=2.2pt,line cap=butt] (5.8017,-1.8671)--(5.7302,-1.8607);
\draw[color={rgb,255:red,220;green,114;blue,46},line width=2.2pt,line cap=butt] (5.7302,-1.8607)--(5.6588,-1.8543);
\draw[color={rgb,255:red,220;green,114;blue,46},line width=2.2pt,line cap=butt] (5.6588,-1.8543)--(5.5873,-1.8479);
\draw[color={rgb,255:red,220;green,114;blue,46},line width=2.2pt,line cap=butt] (5.5873,-1.8479)--(5.5159,-1.8414);
\draw[color={rgb,255:red,220;green,114;blue,46},line width=2.2pt,line cap=butt] (5.5159,-1.8414)--(5.4473,-1.8063);
\draw[color={rgb,255:red,220;green,114;blue,46},line width=2.2pt,line cap=butt] (5.4473,-1.8063)--(5.3787,-1.7711);
\draw[color={rgb,255:red,219;green,113;blue,46},line width=2.2pt,line cap=butt] (5.3787,-1.7711)--(5.3068,-1.7816);
\draw[color={rgb,255:red,219;green,113;blue,46},line width=2.2pt,line cap=butt] (5.3068,-1.7816)--(5.2349,-1.7922);
\draw[color={rgb,255:red,219;green,113;blue,46},line width=2.2pt,line cap=butt] (5.2349,-1.7922)--(5.1630,-1.8027);
\draw[color={rgb,255:red,219;green,113;blue,45},line width=2.2pt,line cap=butt] (5.1630,-1.8027)--(5.0911,-1.8132);
\draw[color={rgb,255:red,219;green,112;blue,45},line width=2.2pt,line cap=butt] (5.0911,-1.8132)--(5.0192,-1.8237);
\draw[color={rgb,255:red,219;green,112;blue,45},line width=2.2pt,line cap=butt] (5.0192,-1.8237)--(4.9473,-1.8342);
\draw[color={rgb,255:red,218;green,112;blue,45},line width=2.2pt,line cap=butt] (4.9473,-1.8342)--(4.8810,-1.8360);
\draw[color={rgb,255:red,218;green,112;blue,45},line width=2.2pt,line cap=butt] (4.8810,-1.8360)--(4.8146,-1.8378);
\draw[color={rgb,255:red,218;green,112;blue,45},line width=2.2pt,line cap=butt] (4.8146,-1.8378)--(4.7482,-1.8396);
\draw[color={rgb,255:red,218;green,111;blue,45},line width=2.2pt,line cap=butt] (4.7482,-1.8396)--(4.6819,-1.8414);
\draw[color={rgb,255:red,218;green,111;blue,45},line width=2.2pt,line cap=butt] (4.6819,-1.8414)--(4.6294,-1.7921);
\draw[color={rgb,255:red,218;green,111;blue,44},line width=2.2pt,line cap=butt] (4.6294,-1.7921)--(4.5769,-1.7428);
\draw[color={rgb,255:red,217;green,111;blue,44},line width=2.2pt,line cap=butt] (4.5769,-1.7428)--(4.5244,-1.6935);
\draw[color={rgb,255:red,217;green,111;blue,44},line width=2.2pt,line cap=butt] (4.5244,-1.6935)--(4.4719,-1.6442);
\draw[color={rgb,255:red,217;green,110;blue,44},line width=2.2pt,line cap=butt] (4.4719,-1.6442)--(4.5083,-1.5784);
\draw[color={rgb,255:red,217;green,110;blue,44},line width=2.2pt,line cap=butt] (4.5083,-1.5784)--(4.5448,-1.5126);
\draw[color={rgb,255:red,217;green,110;blue,44},line width=2.2pt,line cap=butt] (4.5448,-1.5126)--(4.5812,-1.4467);
\draw[color={rgb,255:red,217;green,110;blue,44},line width=2.2pt,line cap=butt] (4.5812,-1.4467)--(4.6176,-1.3809);
\draw[color={rgb,255:red,216;green,109;blue,44},line width=2.2pt,line cap=butt] (4.6176,-1.3809)--(4.6541,-1.3151);
\draw[color={rgb,255:red,216;green,109;blue,43},line width=2.2pt,line cap=butt] (4.6541,-1.3151)--(4.6905,-1.2493);
\draw[color={rgb,255:red,216;green,109;blue,43},line width=2.2pt,line cap=butt] (4.6905,-1.2493)--(4.7269,-1.1835);
\draw[color={rgb,255:red,216;green,109;blue,43},line width=2.2pt,line cap=butt] (4.7269,-1.1835)--(4.7633,-1.1177);
\draw[color={rgb,255:red,216;green,109;blue,43},line width=2.2pt,line cap=butt] (4.7633,-1.1177)--(4.7998,-1.0519);
\draw[color={rgb,255:red,216;green,108;blue,43},line width=2.2pt,line cap=butt] (4.7998,-1.0519)--(4.8362,-0.9861);
\draw[color={rgb,255:red,215;green,108;blue,43},line width=2.2pt,line cap=butt] (4.8362,-0.9861)--(4.8726,-0.9202);
\draw[color={rgb,255:red,215;green,108;blue,43},line width=2.2pt,line cap=butt] (4.8726,-0.9202)--(4.9090,-0.8544);
\draw[color={rgb,255:red,215;green,108;blue,42},line width=2.2pt,line cap=butt] (4.9090,-0.8544)--(4.8978,-0.7948);
\draw[color={rgb,255:red,215;green,107;blue,42},line width=2.2pt,line cap=butt] (4.8978,-0.7948)--(4.8865,-0.7352);
\draw[color={rgb,255:red,215;green,107;blue,42},line width=2.2pt,line cap=butt] (4.8865,-0.7352)--(4.8752,-0.6756);
\draw[color={rgb,255:red,215;green,107;blue,42},line width=2.2pt,line cap=butt] (4.8752,-0.6756)--(4.8675,-0.6026);
\draw[color={rgb,255:red,214;green,107;blue,42},line width=2.2pt,line cap=butt] (4.8675,-0.6026)--(4.8597,-0.5296);
\draw[color={rgb,255:red,214;green,107;blue,42},line width=2.2pt,line cap=butt] (4.8597,-0.5296)--(4.8520,-0.4565);
\draw[color={rgb,255:red,214;green,106;blue,42},line width=2.2pt,line cap=butt] (4.8520,-0.4565)--(4.8442,-0.3835);
\draw[color={rgb,255:red,214;green,106;blue,42},line width=2.2pt,line cap=butt] (4.8442,-0.3835)--(4.8365,-0.3105);
\draw[color={rgb,255:red,214;green,106;blue,41},line width=2.2pt,line cap=butt] (4.8365,-0.3105)--(4.8287,-0.2374);
\draw[color={rgb,255:red,214;green,106;blue,41},line width=2.2pt,line cap=butt] (4.8287,-0.2374)--(4.8210,-0.1644);
\draw[color={rgb,255:red,213;green,106;blue,41},line width=2.2pt,line cap=butt] (4.8210,-0.1644)--(4.8132,-0.0914);
\draw[color={rgb,255:red,213;green,105;blue,41},line width=2.2pt,line cap=butt] (4.8132,-0.0914)--(4.8055,-0.0184);
\draw[color={rgb,255:red,213;green,105;blue,41},line width=2.2pt,line cap=butt] (4.8055,-0.0184)--(4.7977,0.0547);
\draw[color={rgb,255:red,213;green,105;blue,41},line width=2.2pt,line cap=butt] (4.7977,0.0547)--(4.7351,0.0456);
\draw[color={rgb,255:red,213;green,105;blue,41},line width=2.2pt,line cap=butt] (4.7351,0.0456)--(4.6725,0.0365);
\draw[color={rgb,255:red,213;green,105;blue,41},line width=2.2pt,line cap=butt] (4.6725,0.0365)--(4.6099,0.0274);
\draw[color={rgb,255:red,213;green,104;blue,41},line width=2.2pt,line cap=butt] (4.6099,0.0274)--(4.5473,0.0184);
\draw[color={rgb,255:red,212;green,104;blue,40},line width=2.2pt,line cap=butt] (4.5473,0.0184)--(4.4866,-0.0255);
\draw[color={rgb,255:red,212;green,104;blue,40},line width=2.2pt,line cap=butt] (4.4866,-0.0255)--(4.4258,-0.0694);
\draw[color={rgb,255:red,212;green,104;blue,40},line width=2.2pt,line cap=butt] (4.4258,-0.0694)--(4.3651,-0.1132);
\draw[color={rgb,255:red,212;green,103;blue,40},line width=2.2pt,line cap=butt] (4.3651,-0.1132)--(4.3043,-0.1571);
\draw[color={rgb,255:red,212;green,103;blue,40},line width=2.2pt,line cap=butt] (4.3043,-0.1571)--(4.2436,-0.2010);
\draw[color={rgb,255:red,212;green,103;blue,40},line width=2.2pt,line cap=butt] (4.2436,-0.2010)--(4.1829,-0.2448);
\draw[color={rgb,255:red,211;green,103;blue,40},line width=2.2pt,line cap=butt] (4.1829,-0.2448)--(4.1221,-0.2887);
\draw[color={rgb,255:red,211;green,103;blue,39},line width=2.2pt,line cap=butt] (4.1221,-0.2887)--(4.0614,-0.3326);
\draw[color={rgb,255:red,211;green,102;blue,39},line width=2.2pt,line cap=butt] (4.0614,-0.3326)--(4.0006,-0.3764);
\draw[color={rgb,255:red,211;green,102;blue,39},line width=2.2pt,line cap=butt] (4.0006,-0.3764)--(3.9399,-0.4203);
\draw[color={rgb,255:red,211;green,102;blue,39},line width=2.2pt,line cap=butt] (3.9399,-0.4203)--(3.8792,-0.4642);
\draw[color={rgb,255:red,211;green,102;blue,39},line width=2.2pt,line cap=butt] (3.8792,-0.4642)--(3.8137,-0.4911);
\draw[color={rgb,255:red,210;green,101;blue,39},line width=2.2pt,line cap=butt] (3.8137,-0.4911)--(3.7482,-0.5180);
\draw[color={rgb,255:red,210;green,101;blue,39},line width=2.2pt,line cap=butt] (3.7482,-0.5180)--(3.6827,-0.5450);
\draw[color={rgb,255:red,210;green,101;blue,39},line width=2.2pt,line cap=butt] (3.6827,-0.5450)--(3.6173,-0.5719);
\draw[color={rgb,255:red,210;green,101;blue,38},line width=2.2pt,line cap=butt] (3.6173,-0.5719)--(3.5518,-0.5988);
\draw[color={rgb,255:red,210;green,101;blue,38},line width=2.2pt,line cap=butt] (3.5518,-0.5988)--(3.4863,-0.6257);
\draw[color={rgb,255:red,210;green,100;blue,38},line width=2.2pt,line cap=butt] (3.4863,-0.6257)--(3.4208,-0.6527);
\draw[color={rgb,255:red,209;green,100;blue,38},line width=2.2pt,line cap=butt] (3.4208,-0.6527)--(3.3554,-0.6796);
\draw[color={rgb,255:red,209;green,100;blue,38},line width=2.2pt,line cap=butt] (3.3554,-0.6796)--(3.3473,-0.6041);
\draw[color={rgb,255:red,209;green,100;blue,38},line width=2.2pt,line cap=butt] (3.3473,-0.6041)--(3.3393,-0.5286);
\draw[color={rgb,255:red,209;green,100;blue,38},line width=2.2pt,line cap=butt] (3.3393,-0.5286)--(3.3313,-0.4531);
\draw[color={rgb,255:red,209;green,99;blue,38},line width=2.2pt,line cap=butt] (3.3313,-0.4531)--(3.3233,-0.3775);
\draw[color={rgb,255:red,209;green,99;blue,37},line width=2.2pt,line cap=butt] (3.3233,-0.3775)--(3.3153,-0.3020);
\draw[color={rgb,255:red,208;green,99;blue,37},line width=2.2pt,line cap=butt] (3.3153,-0.3020)--(3.3073,-0.2265);
\draw[color={rgb,255:red,208;green,99;blue,37},line width=2.2pt,line cap=butt] (3.3073,-0.2265)--(3.2992,-0.1510);
\draw[color={rgb,255:red,208;green,98;blue,37},line width=2.2pt,line cap=butt] (3.2992,-0.1510)--(3.2912,-0.0755);
\draw[color={rgb,255:red,208;green,98;blue,37},line width=2.2pt,line cap=butt] (3.2912,-0.0755)--(3.2832,0.0000);
\draw[color={rgb,255:red,208;green,98;blue,37},line width=2.2pt,line cap=butt] (3.2832,0.0000)--(3.2286,0.0535);
\draw[color={rgb,255:red,208;green,98;blue,37},line width=2.2pt,line cap=butt] (3.2286,0.0535)--(3.1740,0.1071);
\draw[color={rgb,255:red,207;green,98;blue,36},line width=2.2pt,line cap=butt] (3.1740,0.1071)--(3.1194,0.1606);
\draw[color={rgb,255:red,207;green,97;blue,36},line width=2.2pt,line cap=butt] (3.1194,0.1606)--(3.0648,0.2141);
\draw[color={rgb,255:red,207;green,97;blue,36},line width=2.2pt,line cap=butt] (3.0648,0.2141)--(3.0102,0.2677);
\draw[color={rgb,255:red,207;green,97;blue,36},line width=2.2pt,line cap=butt] (3.0102,0.2677)--(2.9556,0.3212);
\draw[color={rgb,255:red,207;green,97;blue,36},line width=2.2pt,line cap=butt] (2.9556,0.3212)--(2.9010,0.3747);
\draw[color={rgb,255:red,206;green,96;blue,36},line width=2.2pt,line cap=butt] (2.9010,0.3747)--(2.8464,0.4283);
\draw[color={rgb,255:red,206;green,96;blue,36},line width=2.2pt,line cap=butt] (2.8464,0.4283)--(2.7918,0.4818);
\draw[color={rgb,255:red,206;green,96;blue,35},line width=2.2pt,line cap=butt] (2.7918,0.4818)--(2.7372,0.5353);
\draw[color={rgb,255:red,206;green,96;blue,35},line width=2.2pt,line cap=butt] (2.7372,0.5353)--(2.6826,0.5889);
\draw[color={rgb,255:red,206;green,95;blue,35},line width=2.2pt,line cap=butt] (2.6826,0.5889)--(2.6241,0.5378);
\draw[color={rgb,255:red,206;green,95;blue,35},line width=2.2pt,line cap=butt] (2.6241,0.5378)--(2.5657,0.4868);
\draw[color={rgb,255:red,205;green,95;blue,35},line width=2.2pt,line cap=butt] (2.5657,0.4868)--(2.5072,0.4357);
\draw[color={rgb,255:red,205;green,95;blue,35},line width=2.2pt,line cap=butt] (2.5072,0.4357)--(2.4487,0.3847);
\draw[color={rgb,255:red,205;green,95;blue,35},line width=2.2pt,line cap=butt] (2.4487,0.3847)--(2.3903,0.3336);
\draw[color={rgb,255:red,205;green,94;blue,35},line width=2.2pt,line cap=butt] (2.3903,0.3336)--(2.3318,0.2825);
\draw[color={rgb,255:red,205;green,94;blue,34},line width=2.2pt,line cap=butt] (2.3318,0.2825)--(2.2733,0.2315);
\draw[color={rgb,255:red,205;green,94;blue,34},line width=2.2pt,line cap=butt] (2.2733,0.2315)--(2.2148,0.1804);
\draw[color={rgb,255:red,204;green,94;blue,34},line width=2.2pt,line cap=butt] (2.2148,0.1804)--(2.1564,0.1294);
\draw[color={rgb,255:red,204;green,93;blue,34},line width=2.2pt,line cap=butt] (2.1564,0.1294)--(2.1587,0.0534);
\draw[color={rgb,255:red,204;green,93;blue,34},line width=2.2pt,line cap=butt] (2.1587,0.0534)--(2.1610,-0.0226);
\draw[color={rgb,255:red,204;green,93;blue,34},line width=2.2pt,line cap=butt] (2.1610,-0.0226)--(2.1634,-0.0985);
\draw[color={rgb,255:red,204;green,93;blue,34},line width=2.2pt,line cap=butt] (2.1634,-0.0985)--(2.1657,-0.1745);
\draw[color={rgb,255:red,203;green,92;blue,33},line width=2.2pt,line cap=butt] (2.1657,-0.1745)--(2.1681,-0.2505);
\draw[color={rgb,255:red,203;green,92;blue,33},line width=2.2pt,line cap=butt] (2.1681,-0.2505)--(2.1704,-0.3265);
\draw[color={rgb,255:red,203;green,92;blue,33},line width=2.2pt,line cap=butt] (2.1704,-0.3265)--(2.1727,-0.4025);
\draw[color={rgb,255:red,203;green,92;blue,33},line width=2.2pt,line cap=butt] (2.1727,-0.4025)--(2.1751,-0.4784);
\draw[color={rgb,255:red,203;green,92;blue,33},line width=2.2pt,line cap=butt] (2.1751,-0.4784)--(2.1774,-0.5544);
\draw[color={rgb,255:red,203;green,91;blue,33},line width=2.2pt,line cap=butt] (2.1774,-0.5544)--(2.1797,-0.6304);
\draw[color={rgb,255:red,202;green,91;blue,33},line width=2.2pt,line cap=butt] (2.1797,-0.6304)--(2.1821,-0.7064);
\draw[color={rgb,255:red,202;green,91;blue,32},line width=2.2pt,line cap=butt] (2.1821,-0.7064)--(2.1105,-0.7019);
\draw[color={rgb,255:red,202;green,91;blue,32},line width=2.2pt,line cap=butt] (2.1105,-0.7019)--(2.0389,-0.6975);
\draw[color={rgb,255:red,202;green,90;blue,32},line width=2.2pt,line cap=butt] (2.0389,-0.6975)--(1.9672,-0.6930);
\draw[color={rgb,255:red,202;green,90;blue,32},line width=2.2pt,line cap=butt] (1.9672,-0.6930)--(1.8956,-0.6886);
\draw[color={rgb,255:red,202;green,90;blue,32},line width=2.2pt,line cap=butt] (1.8956,-0.6886)--(1.8240,-0.6841);
\draw[color={rgb,255:red,201;green,90;blue,32},line width=2.2pt,line cap=butt] (1.8240,-0.6841)--(1.7524,-0.6797);
\draw[color={rgb,255:red,201;green,90;blue,32},line width=2.2pt,line cap=butt] (1.7524,-0.6797)--(1.6808,-0.6753);
\draw[color={rgb,255:red,201;green,89;blue,32},line width=2.2pt,line cap=butt] (1.6808,-0.6753)--(1.6092,-0.6708);
\draw[color={rgb,255:red,201;green,89;blue,31},line width=2.2pt,line cap=butt] (1.6092,-0.6708)--(1.5376,-0.6664);
\draw[color={rgb,255:red,201;green,89;blue,31},line width=2.2pt,line cap=butt] (1.5376,-0.6664)--(1.4905,-0.6107);
\draw[color={rgb,255:red,201;green,89;blue,31},line width=2.2pt,line cap=butt] (1.4905,-0.6107)--(1.4434,-0.5550);
\draw[color={rgb,255:red,200;green,89;blue,31},line width=2.2pt,line cap=butt] (1.4434,-0.5550)--(1.3963,-0.4993);
\draw[color={rgb,255:red,200;green,88;blue,31},line width=2.2pt,line cap=butt] (1.3963,-0.4993)--(1.3492,-0.4437);
\draw[color={rgb,255:red,200;green,88;blue,31},line width=2.2pt,line cap=butt] (1.3492,-0.4437)--(1.3021,-0.3880);
\draw[color={rgb,255:red,200;green,88;blue,31},line width=2.2pt,line cap=butt] (1.3021,-0.3880)--(1.2550,-0.3323);
\draw[color={rgb,255:red,200;green,88;blue,31},line width=2.2pt,line cap=butt] (1.2550,-0.3323)--(1.2080,-0.2766);
\draw[color={rgb,255:red,200;green,87;blue,30},line width=2.2pt,line cap=butt] (1.2080,-0.2766)--(1.1609,-0.2210);
\draw[color={rgb,255:red,199;green,87;blue,30},line width=2.2pt,line cap=butt] (1.1609,-0.2210)--(1.1138,-0.1653);
\draw[color={rgb,255:red,199;green,87;blue,30},line width=2.2pt,line cap=butt] (1.1138,-0.1653)--(1.0498,-0.1322);
\draw[color={rgb,255:red,199;green,87;blue,30},line width=2.2pt,line cap=butt] (1.0498,-0.1322)--(0.9858,-0.0992);
\draw[color={rgb,255:red,199;green,87;blue,30},line width=2.2pt,line cap=butt] (0.9858,-0.0992)--(0.9219,-0.0661);
\draw[color={rgb,255:red,199;green,86;blue,30},line width=2.2pt,line cap=butt] (0.9219,-0.0661)--(0.8579,-0.0331);
\draw[color={rgb,255:red,199;green,86;blue,30},line width=2.2pt,line cap=butt] (0.8579,-0.0331)--(0.7939,0.0000);
\draw[color={rgb,255:red,199;green,86;blue,30},line width=2.2pt,line cap=butt] (0.7939,0.0000)--(0.7183,0.0000);
\draw[color={rgb,255:red,198;green,86;blue,29},line width=2.2pt,line cap=butt] (0.7183,0.0000)--(0.6934,-0.0663);
\draw[color={rgb,255:red,198;green,85;blue,29},line width=2.2pt,line cap=butt] (0.6934,-0.0663)--(0.6685,-0.1326);
\draw[color={rgb,255:red,198;green,85;blue,29},line width=2.2pt,line cap=butt] (0.6685,-0.1326)--(0.6437,-0.1989);
\draw[color={rgb,255:red,198;green,85;blue,29},line width=2.2pt,line cap=butt] (0.6437,-0.1989)--(0.6188,-0.2652);
\draw[color={rgb,255:red,198;green,85;blue,29},line width=2.2pt,line cap=butt] (0.6188,-0.2652)--(0.5939,-0.3316);
\draw[color={rgb,255:red,198;green,85;blue,29},line width=2.2pt,line cap=butt] (0.5939,-0.3316)--(0.5691,-0.3979);
\draw[color={rgb,255:red,197;green,84;blue,29},line width=2.2pt,line cap=butt] (0.5691,-0.3979)--(0.5442,-0.4642);
\draw[color={rgb,255:red,197;green,84;blue,29},line width=2.2pt,line cap=butt] (0.5442,-0.4642)--(0.4919,-0.5196);
\draw[color={rgb,255:red,197;green,84;blue,28},line width=2.2pt,line cap=butt] (0.4919,-0.5196)--(0.4395,-0.5750);
\draw[color={rgb,255:red,197;green,84;blue,28},line width=2.2pt,line cap=butt] (0.4395,-0.5750)--(0.3872,-0.6305);
\draw[color={rgb,255:red,197;green,84;blue,28},line width=2.2pt,line cap=butt] (0.3872,-0.6305)--(0.3349,-0.6859);
\draw[color={rgb,255:red,197;green,83;blue,28},line width=2.2pt,line cap=butt] (0.3349,-0.6859)--(0.2826,-0.7413);
\draw[color={rgb,255:red,196;green,83;blue,28},line width=2.2pt,line cap=butt] (0.2826,-0.7413)--(0.2302,-0.7967);
\draw[color={rgb,255:red,196;green,83;blue,28},line width=2.2pt,line cap=butt] (0.2302,-0.7967)--(0.1779,-0.8522);
\draw[color={rgb,255:red,196;green,83;blue,28},line width=2.2pt,line cap=butt] (0.1779,-0.8522)--(0.1256,-0.9076);
\draw[color={rgb,255:red,196;green,82;blue,27},line width=2.2pt,line cap=butt] (0.1256,-0.9076)--(0.0733,-0.9630);
\draw[color={rgb,255:red,196;green,82;blue,27},line width=2.2pt,line cap=butt] (0.0733,-0.9630)--(0.0209,-1.0185);
\draw[color={rgb,255:red,195;green,82;blue,27},line width=2.2pt,line cap=butt] (0.0209,-1.0185)--(-0.0314,-1.0739);
\draw[color={rgb,255:red,195;green,82;blue,27},line width=2.2pt,line cap=butt] (-0.0314,-1.0739)--(-0.0837,-1.1293);
\draw[color={rgb,255:red,195;green,81;blue,27},line width=2.2pt,line cap=butt] (-0.0837,-1.1293)--(-0.1360,-1.1847);
\draw[color={rgb,255:red,195;green,81;blue,27},line width=2.2pt,line cap=butt] (-0.1360,-1.1847)--(-0.1884,-1.2402);
\draw[color={rgb,255:red,195;green,81;blue,27},line width=2.2pt,line cap=butt] (-0.1884,-1.2402)--(-0.2407,-1.2956);
\draw[color={rgb,255:red,195;green,81;blue,26},line width=2.2pt,line cap=butt] (-0.2407,-1.2956)--(-0.2930,-1.3510);
\draw[color={rgb,255:red,194;green,81;blue,26},line width=2.2pt,line cap=butt] (-0.2930,-1.3510)--(-0.3453,-1.4065);
\draw[color={rgb,255:red,194;green,80;blue,26},line width=2.2pt,line cap=butt] (-0.3453,-1.4065)--(-0.3977,-1.4619);
\draw[color={rgb,255:red,194;green,80;blue,26},line width=2.2pt,line cap=butt] (-0.3977,-1.4619)--(-0.4500,-1.5173);
\end{scope}
\fill[white,draw=Ink,line width=1.2pt] (-0.4500,-1.5173) circle (0.045);
\end{scope}
\end{tikzpicture}};
  \node[anchor=north west,inner sep=0,text width=22cm] at ([xshift=1.2cm,yshift=-10.2cm]current page.north west)
    {\bfseries\sz{32}{36}\hi Qual é o menor caminho que visita todas as regiões?};
  \node[anchor=north west,inner sep=0,text width=14cm] at ([xshift=1.2cm,yshift=-11.8cm]current page.north west)
    {\sffamily\bfseries\sz{26}{30}\ink Gabriel Freire Ushijima\par\vspace{2pt}\mdseries\sz{22}{26}\muted Orientador: Ernesto G. Birgin · IME-USP};
  \node[anchor=south east,inner sep=0,align=right] at ([xshift=-1.0cm,yshift=0.4cm]current page.south east) {\parbox{4.6cm}{\raggedleft\includegraphics[width=3.8cm]{../poster/figures/fapesp-logo.png}\par\vspace{3pt}\sffamily\sz{16}{19}\muted Processo 2025/13861-1}};
\end{tikzpicture}
\newpage

%%%%%%%%%%%%%%%%%%%%%%%%%%%%%%%% 2. O problema
\begin{tikzpicture}[remember picture,overlay]
  \fill[white] (current page.south west) rectangle (current page.north east);
  \slidetitle{O Problema de Visita de Polígonos}
  \node[anchor=north west,inner sep=0,text width=10.6cm] at ([xshift=1.2cm,yshift=-2.6cm]current page.north west) {%
    {\bfseries\sz{28}{32}\hi Dados}\par\vspace{2pt}
    {\sz{26}{31}\ink um ponto inicial $s$, um ponto final $t$ e polígonos $P_1,\dots,P_k$ no plano\par}\vspace{0.55cm}
    {\bfseries\sz{28}{32}\hi Objetivo}\par\vspace{2pt}
    {\sz{26}{31}\ink o caminho de $s$ a $t$ de \textbf{menor comprimento} que \textbf{toca} cada $P_i$\par}};
  \node[anchor=north east,inner sep=0] at ([xshift=-1.1cm,yshift=-2.45cm]current page.north east) {%
    \begin{tikzpicture}[x=9.4cm,y=9.4cm]\fill[Pale,rounded corners=0.35cm] (-0.05,-0.05) rectangle (1.12,1.05);\def\uspaspect{1.0717}
\fill[color={rgb,255:red,60;green,159;blue,144}] (0.2773,0.4664) -- (0.2754,0.4541) -- (0.2813,0.4533) -- (0.2804,0.4472) -- (0.3042,0.4433) -- (0.3020,0.4280) -- (0.2961,0.4292) -- (0.2975,0.4386) -- (0.2855,0.4404) -- (0.2825,0.4221) -- (0.2945,0.4203) -- (0.2954,0.4253) -- (0.3014,0.4243) -- (0.3007,0.4197) -- (0.3187,0.4169) -- (0.3241,0.4530) -- (0.3189,0.4539) -- (0.3223,0.4779) -- (0.3164,0.4788) -- (0.3155,0.4731) -- (0.3096,0.4738) -- (0.3106,0.4795) -- (0.3042,0.4805) -- (0.3011,0.4626) -- cycle;
\fill[color={rgb,255:red,20;green,45;blue,56}] (0.3370,0.4131) -- (0.3595,0.4099) -- (0.3616,0.4235) -- (0.3576,0.4240) -- (0.3647,0.4687) -- (0.3460,0.4720) -- cycle;
\fill[color={rgb,255:red,20;green,48;blue,58}] (0.3465,0.4860) -- (0.3473,0.4943) -- (0.3315,0.4965) -- (0.3315,0.4976) -- (0.3276,0.4981) -- (0.3275,0.4972) -- (0.3138,0.4995) -- (0.3139,0.5007) -- (0.3103,0.5012) -- (0.3102,0.4999) -- (0.2936,0.5026) -- (0.2926,0.4964) -- (0.2950,0.4962) -- (0.2946,0.4941) -- cycle;
\fill[color={rgb,255:red,50;green,146;blue,134}] (0.2753,0.3475) -- (0.2727,0.3305) -- (0.2686,0.3311) -- (0.2661,0.3144) -- (0.2615,0.3154) -- (0.2623,0.3241) -- (0.2634,0.3243) -- (0.2662,0.3435) -- (0.2696,0.3456) -- (0.2698,0.3488) -- (0.2625,0.3502) -- (0.2548,0.3500) -- (0.2480,0.3484) -- (0.2409,0.3450) -- (0.2547,0.3191) -- (0.2565,0.3200) -- (0.2574,0.3189) -- (0.2575,0.3141) -- (0.2699,0.3114) -- (0.2677,0.2971) -- (0.3445,0.2853) -- (0.3464,0.2985) -- (0.3507,0.2981) -- (0.3538,0.3181) -- (0.3493,0.3187) -- (0.3519,0.3358) -- cycle;
\fill[color={rgb,255:red,47;green,143;blue,131}] (0.4283,0.4247) -- (0.4074,0.3855) -- (0.4728,0.3506) -- (0.4937,0.3901) -- cycle;
\fill[color={rgb,255:red,46;green,141;blue,130}] (0.5189,0.3881) -- (0.5379,0.3266) -- (0.5412,0.3275) -- (0.5602,0.3636) -- (0.5592,0.3670) -- cycle;
\fill[color={rgb,255:red,48;green,144;blue,132}] (0.4540,0.3129) -- (0.4486,0.3028) -- (0.5352,0.2571) -- (0.5402,0.2671) -- cycle;
\fill[color={rgb,255:red,44;green,138;blue,128}] (0.6360,0.2188) -- (0.6444,0.2344) -- (0.6346,0.2394) -- (0.6320,0.2343) -- (0.6260,0.2375) -- (0.6286,0.2426) -- (0.6239,0.2454) -- (0.6294,0.2563) -- (0.6501,0.2451) -- (0.6557,0.2561) -- (0.6351,0.2668) -- (0.6377,0.2723) -- (0.6588,0.2610) -- (0.6648,0.2721) -- (0.6436,0.2833) -- (0.6461,0.2884) -- (0.6674,0.2771) -- (0.6733,0.2877) -- (0.6518,0.2991) -- (0.6545,0.3047) -- (0.6758,0.2932) -- (0.6815,0.3035) -- (0.6601,0.3151) -- (0.6572,0.3098) -- (0.6304,0.3239) -- (0.6278,0.3188) -- (0.6491,0.3077) -- (0.6462,0.3022) -- (0.6253,0.3132) -- (0.6225,0.3081) -- (0.6435,0.2970) -- (0.6409,0.2914) -- (0.6194,0.3026) -- (0.6166,0.2972) -- (0.6380,0.2863) -- (0.6354,0.2806) -- (0.6137,0.2923) -- (0.6107,0.2864) -- (0.6324,0.2750) -- (0.6272,0.2646) -- (0.5946,0.2816) -- (0.5861,0.2656) -- (0.6182,0.2483) -- (0.6129,0.2381) -- (0.6291,0.2293) -- (0.6263,0.2239) -- cycle;
\fill[color={rgb,255:red,43;green,137;blue,127}] (0.7339,0.2080) -- (0.7446,0.2290) -- (0.7499,0.2263) -- (0.7698,0.2650) -- (0.7594,0.2704) -- (0.7427,0.2377) -- (0.7377,0.2402) -- (0.7493,0.2626) -- (0.7427,0.2659) -- (0.7315,0.2437) -- (0.7274,0.2458) -- (0.7440,0.2779) -- (0.7319,0.2839) -- (0.7134,0.2476) -- (0.7081,0.2503) -- (0.7103,0.2545) -- (0.7048,0.2572) -- (0.6939,0.2358) -- (0.6886,0.2385) -- (0.6970,0.2549) -- (0.6908,0.2580) -- (0.6797,0.2363) -- (0.6966,0.2278) -- (0.7018,0.2382) -- (0.7066,0.2358) -- (0.7012,0.2250) -- (0.7124,0.2193) -- (0.7150,0.2245) -- (0.7204,0.2218) -- (0.7175,0.2162) -- cycle;
\fill[color={rgb,255:red,36;green,128;blue,120}] (0.7635,0.2274) -- (0.7485,0.1991) -- (0.7764,0.1844) -- (0.7879,0.2059) -- (0.7922,0.2034) -- (0.7833,0.1868) -- (0.7957,0.1801) -- (0.8041,0.1957) -- (0.8081,0.1936) -- (0.7998,0.1782) -- (0.8121,0.1717) -- (0.8329,0.2106) -- (0.8196,0.2177) -- (0.8109,0.2012) -- (0.8077,0.2032) -- (0.8163,0.2193) -- (0.8037,0.2261) -- (0.7951,0.2107) -- (0.7852,0.2158) -- (0.7873,0.2195) -- (0.7932,0.2163) -- (0.8003,0.2291) -- (0.7868,0.2364) -- (0.7779,0.2198) -- cycle;
\fill[color={rgb,255:red,34;green,126;blue,118}] (0.8797,0.1949) -- (0.8738,0.1840) -- (0.8725,0.1847) -- (0.8636,0.1678) -- (0.8652,0.1670) -- (0.8593,0.1556) -- (0.9262,0.1208) -- (0.9321,0.1317) -- (0.9337,0.1308) -- (0.9425,0.1473) -- (0.9408,0.1481) -- (0.9468,0.1597) -- cycle;
\fill[color={rgb,255:red,33;green,125;blue,117}] (0.9910,0.2129) -- (1.0191,0.1979) -- (1.0318,0.2211) -- (1.0037,0.2363) -- (1.0047,0.2397) -- (1.0043,0.2436) -- (1.0026,0.2468) -- (0.9997,0.2491) -- (0.9963,0.2502) -- (0.9924,0.2499) -- (0.9891,0.2481) -- (0.9870,0.2452) -- (0.9260,0.2781) -- (0.9135,0.2550) -- (0.9746,0.2218) -- (0.9684,0.2101) -- (0.9846,0.2012) -- cycle;
\fill[color={rgb,255:red,28;green,104;blue,101}] (0.6518,0.5140) -- (0.6523,0.5137) -- (0.6498,0.5092) -- (0.6494,0.5094) -- (0.6365,0.4850) -- (0.6464,0.4796) -- (0.6594,0.5040) -- (0.6590,0.5042) -- (0.6615,0.5088) -- (0.6620,0.5087) -- (0.6954,0.5714) -- (0.6949,0.5719) -- (0.6852,0.5768) -- cycle;
\fill[color={rgb,255:red,29;green,107;blue,103}] (0.6856,0.5350) -- (0.7013,0.5269) -- (0.7044,0.5326) -- (0.7022,0.5336) -- (0.7040,0.5369) -- (0.7062,0.5358) -- (0.7140,0.5504) -- (0.6982,0.5585) -- cycle;
\fill[color={rgb,255:red,29;green,110;blue,106}] (0.7236,0.5476) -- (0.7494,0.5337) -- (0.7239,0.4860) -- (0.7162,0.4897) -- (0.7265,0.5089) -- (0.7139,0.5154) -- (0.6975,0.4848) -- (0.7100,0.4783) -- (0.7149,0.4873) -- (0.7225,0.4834) -- (0.7081,0.4560) -- (0.6820,0.4700) -- (0.6793,0.4645) -- (0.7127,0.4468) -- (0.7265,0.4728) -- (0.7284,0.4719) -- (0.7477,0.5083) -- (0.7457,0.5093) -- (0.7597,0.5352) -- (0.7262,0.5529) -- cycle;
\fill[color={rgb,255:red,27;green,97;blue,96}] (0.5909,0.5562) -- (0.6166,0.5800) -- (0.6218,0.5745) -- (0.6289,0.5809) -- (0.6236,0.5866) -- (0.6350,0.5972) -- (0.6280,0.6047) -- (0.6167,0.5941) -- (0.5906,0.6217) -- (0.5839,0.6154) -- (0.5884,0.6105) -- (0.5664,0.5681) -- (0.5750,0.5636) -- (0.5860,0.5842) -- (0.5896,0.5826) -- (0.5919,0.5867) -- (0.5883,0.5886) -- (0.5958,0.6027) -- (0.5994,0.5988) -- (0.5964,0.5960) -- (0.5996,0.5925) -- (0.6026,0.5954) -- (0.6097,0.5876) -- (0.6041,0.5822) -- (0.6011,0.5853) -- (0.5978,0.5822) -- (0.6007,0.5791) -- (0.5839,0.5635) -- cycle;
\fill[color={rgb,255:red,27;green,91;blue,91}] (0.5361,0.6285) -- (0.5310,0.6310) -- (0.5203,0.6105) -- (0.5187,0.6111) -- (0.5160,0.6055) -- (0.5097,0.6035) -- (0.5121,0.5958) -- (0.5138,0.5874) -- (0.5221,0.5903) -- (0.5266,0.5991) -- (0.5259,0.5993) -- (0.5285,0.6046) -- (0.5303,0.6051) -- (0.5312,0.6065) -- (0.5313,0.6079) -- (0.5308,0.6088) -- (0.5418,0.6299) -- (0.5379,0.6318) -- cycle;
\fill[color={rgb,255:red,26;green,88;blue,89}] (0.5036,0.6153) -- (0.4939,0.6205) -- (0.4744,0.5841) -- (0.4843,0.5789) -- cycle;
\fill[color={rgb,255:red,26;green,85;blue,87}] (0.4627,0.5386) -- (0.4680,0.5487) -- (0.4650,0.5502) -- (0.4689,0.5577) -- (0.4622,0.5611) -- (0.4641,0.5644) -- (0.4705,0.5610) -- (0.4731,0.5659) -- (0.4667,0.5693) -- (0.4669,0.5700) -- (0.4384,0.5850) -- (0.4561,0.6180) -- (0.4802,0.6054) -- (0.4903,0.6246) -- (0.5041,0.6399) -- (0.5067,0.6384) -- (0.5118,0.6483) -- (0.4772,0.6665) -- (0.4720,0.6565) -- (0.5023,0.6407) -- (0.4882,0.6254) -- (0.4792,0.6078) -- (0.4775,0.6086) -- (0.4789,0.6111) -- (0.4722,0.6146) -- (0.4813,0.6315) -- (0.4753,0.6348) -- (0.4747,0.6336) -- (0.4709,0.6358) -- (0.4703,0.6346) -- (0.4131,0.6648) -- (0.4048,0.6490) -- (0.4289,0.6363) -- (0.4280,0.6347) -- (0.4136,0.6423) -- (0.3962,0.6093) -- (0.4102,0.6015) -- (0.4099,0.6006) -- (0.3812,0.6156) -- (0.3727,0.5997) -- (0.4033,0.5833) -- (0.4016,0.5796) -- (0.4006,0.5803) -- (0.3996,0.5782) -- (0.3987,0.5786) -- (0.3985,0.5781) -- (0.3968,0.5789) -- (0.3941,0.5741) -- (0.3964,0.5730) -- (0.3959,0.5721) -- (0.3919,0.5743) -- (0.3930,0.5765) -- (0.3935,0.5764) -- (0.3953,0.5799) -- (0.3918,0.5817) -- (0.3887,0.5758) -- (0.3798,0.5806) -- (0.3791,0.5791) -- (0.3756,0.5808) -- (0.3703,0.5709) -- (0.3980,0.5563) -- (0.4034,0.5664) -- (0.3966,0.5699) -- (0.3976,0.5719) -- (0.4028,0.5691) -- (0.4057,0.5745) -- (0.4032,0.5762) -- (0.4062,0.5818) -- (0.4344,0.5669) -- (0.4295,0.5576) -- (0.4327,0.5559) -- (0.4377,0.5652) -- (0.4446,0.5614) -- (0.4416,0.5556) -- (0.4422,0.5544) -- (0.4409,0.5518) -- (0.4532,0.5452) -- (0.4525,0.5440) -- cycle;
\fill[color={rgb,255:red,21;green,54;blue,63}] (0.2178,0.6379) -- (0.2218,0.6459) -- (0.2135,0.6504) -- (0.2149,0.6537) -- (0.2237,0.6491) -- (0.2252,0.6522) -- (0.2260,0.6518) -- (0.2224,0.6448) -- (0.2272,0.6422) -- (0.2331,0.6583) -- (0.2338,0.6579) -- (0.2351,0.6604) -- (0.2317,0.6623) -- (0.2274,0.6542) -- (0.2267,0.6545) -- (0.2280,0.6570) -- (0.2191,0.6614) -- (0.2219,0.6667) -- (0.2305,0.6623) -- (0.2329,0.6664) -- (0.2344,0.6656) -- (0.2363,0.6695) -- (0.2264,0.6749) -- (0.2273,0.6767) -- (0.2308,0.6749) -- (0.2330,0.6784) -- (0.2217,0.6845) -- (0.2196,0.6810) -- (0.2253,0.6778) -- (0.2243,0.6759) -- (0.2051,0.6859) -- (0.2041,0.6838) -- (0.2020,0.6848) -- (0.1990,0.6790) -- (0.2201,0.6678) -- (0.2170,0.6625) -- (0.1961,0.6737) -- (0.1917,0.6660) -- (0.2133,0.6546) -- (0.2115,0.6514) -- (0.1905,0.6627) -- (0.1862,0.6543) -- cycle;
\fill[color={rgb,255:red,22;green,57;blue,65}] (0.2374,0.6641) -- (0.2309,0.6519) -- (0.2525,0.6406) -- (0.2506,0.6368) -- (0.2483,0.6375) -- (0.2438,0.6373) -- (0.2400,0.6361) -- (0.2362,0.6340) -- (0.2333,0.6308) -- (0.2312,0.6271) -- (0.2304,0.6233) -- (0.2304,0.6186) -- (0.2316,0.6145) -- (0.2343,0.6109) -- (0.2373,0.6080) -- (0.2415,0.6061) -- (0.2456,0.6054) -- (0.2500,0.6057) -- (0.2545,0.6073) -- (0.2582,0.6102) -- (0.2610,0.6140) -- (0.2624,0.6182) -- (0.2627,0.6228) -- (0.2617,0.6272) -- (0.2597,0.6308) -- (0.2564,0.6343) -- (0.2579,0.6377) -- (0.2781,0.6271) -- (0.2843,0.6392) -- (0.2832,0.6405) -- (0.2667,0.6495) -- (0.2663,0.6487) -- (0.2643,0.6499) -- (0.2658,0.6526) -- (0.2739,0.6484) -- (0.2766,0.6537) -- (0.2986,0.6421) -- (0.2991,0.6431) -- (0.3275,0.6282) -- (0.3309,0.6345) -- (0.2726,0.6658) -- (0.2743,0.6694) -- (0.2986,0.6568) -- (0.3067,0.6723) -- (0.2528,0.7007) -- (0.2447,0.6853) -- (0.2688,0.6725) -- (0.2665,0.6687) -- (0.2250,0.6906) -- (0.2218,0.6846) -- (0.2389,0.6753) -- (0.2389,0.6739) -- (0.2551,0.6653) -- (0.2521,0.6599) -- (0.2596,0.6559) -- (0.2582,0.6531) -- cycle;
\fill[color={rgb,255:red,21;green,51;blue,61}] (0.2848,0.6399) -- (0.2663,0.6043) -- (0.2756,0.5997) -- (0.2828,0.6139) -- (0.3038,0.6030) -- (0.3016,0.5983) -- (0.2848,0.6067) -- (0.2800,0.5974) -- (0.3058,0.5838) -- (0.3120,0.5958) -- (0.3076,0.5980) -- (0.3121,0.6069) -- (0.3202,0.6025) -- (0.3234,0.6089) -- (0.3150,0.6133) -- (0.3186,0.6200) -- (0.2978,0.6309) -- (0.2987,0.6327) -- cycle;
\fill[color={rgb,255:red,22;green,60;blue,68}] (0.3104,0.7315) -- (0.3509,0.8065) -- (0.3314,0.8171) -- (0.3354,0.8246) -- (0.3539,0.8148) -- (0.3682,0.8410) -- (0.3699,0.8443) -- (0.3847,0.8365) -- (0.3902,0.8466) -- (0.3637,0.8609) -- (0.3582,0.8505) -- (0.3565,0.8474) -- (0.3217,0.8660) -- (0.3070,0.8388) -- (0.3255,0.8292) -- (0.3220,0.8229) -- (0.3045,0.8321) -- (0.2640,0.7562) -- cycle;
\fill[color={rgb,255:red,23;green,67;blue,73}] (0.3371,0.9405) -- (0.3620,0.9273) -- (0.3938,0.9867) -- (0.3690,1.0000) -- cycle;
\fill[color={rgb,255:red,24;green,70;blue,75}] (0.3642,0.8885) -- (0.3995,0.8697) -- (0.4313,0.9291) -- (0.3963,0.9479) -- cycle;
\fill[color={rgb,255:red,24;green,73;blue,77}] (0.4284,0.8710) -- (0.4464,0.8616) -- (0.4493,0.8672) -- (0.4563,0.8637) -- (0.4522,0.8559) -- (0.4669,0.8481) -- (0.4707,0.8553) -- (0.4926,0.8439) -- (0.4897,0.8380) -- (0.5269,0.8182) -- (0.5707,0.9005) -- (0.5333,0.9203) -- (0.5319,0.9177) -- (0.4883,0.9407) -- (0.4896,0.9432) -- (0.4717,0.9526) -- cycle;
\fill[color={rgb,255:red,24;green,76;blue,80}] (0.5653,0.8608) -- (0.5646,0.8604) -- (0.5592,0.8494) -- (0.5604,0.8485) -- (0.5453,0.8198) -- (0.5438,0.8207) -- (0.5415,0.8161) -- (0.5837,0.7942) -- (0.5835,0.7937) -- (0.5900,0.7902) -- (0.6089,0.8257) -- (0.6510,0.8035) -- (0.6526,0.8066) -- (0.6547,0.8054) -- (0.6700,0.8341) -- (0.7108,0.8124) -- (0.6948,0.7818) -- (0.6961,0.7810) -- (0.6956,0.7799) -- (0.7221,0.7657) -- (0.7382,0.7960) -- (0.7375,0.7965) -- (0.7418,0.8037) -- (0.6219,0.8683) -- (0.5801,0.8897) -- cycle;
\fill[color={rgb,255:red,25;green,82;blue,84}] (0.5161,0.6661) -- (0.5263,0.6614) -- (0.5440,0.6994) -- (0.5338,0.7039) -- cycle;
\fill[color={rgb,255:red,28;green,100;blue,98}] (0.6382,0.6146) -- (0.6441,0.6257) -- (0.6360,0.6302) -- (0.6378,0.6333) -- (0.6459,0.6289) -- (0.6477,0.6322) -- (0.6399,0.6373) -- (0.6419,0.6403) -- (0.6495,0.6352) -- (0.6568,0.6457) -- (0.6522,0.6487) -- (0.6565,0.6548) -- (0.6531,0.6571) -- (0.6650,0.6743) -- (0.6447,0.6884) -- (0.6327,0.6712) -- (0.6296,0.6732) -- (0.6255,0.6673) -- (0.6144,0.6748) -- (0.6071,0.6644) -- (0.6241,0.6527) -- (0.6212,0.6485) -- (0.6109,0.6549) -- (0.6087,0.6516) -- (0.6188,0.6453) -- (0.6164,0.6407) -- (0.6002,0.6495) -- (0.5942,0.6386) -- cycle;
\fill[color={rgb,255:red,59;green,158;blue,143}] (0.1804,0.4236) -- (0.1810,0.4310) -- (0.2116,0.4275) -- (0.2125,0.4365) -- (0.2038,0.4375) -- (0.2046,0.4450) -- (0.2133,0.4442) -- (0.2142,0.4532) -- (0.1783,0.4571) -- (0.1781,0.4539) -- (0.1765,0.4539) -- (0.1769,0.4619) -- (0.1802,0.4618) -- (0.1803,0.4657) -- (0.2101,0.4635) -- (0.2112,0.4739) -- (0.1810,0.4770) -- (0.1808,0.4732) -- (0.1788,0.4731) -- (0.1789,0.4768) -- (0.1746,0.4770) -- (0.1740,0.4666) -- (0.1782,0.4666) -- (0.1782,0.4642) -- (0.1747,0.4645) -- (0.1738,0.4514) -- (0.1776,0.4510) -- (0.1761,0.4366) -- (0.1714,0.4371) -- (0.1720,0.4423) -- (0.1711,0.4426) -- (0.1714,0.4453) -- (0.1673,0.4457) -- (0.1680,0.4527) -- (0.1661,0.4530) -- (0.1665,0.4559) -- (0.1338,0.4595) -- (0.1341,0.4634) -- (0.1290,0.4639) -- (0.1290,0.4623) -- (0.1222,0.4626) -- (0.1213,0.4520) -- (0.1325,0.4514) -- (0.1321,0.4472) -- (0.0999,0.4502) -- (0.0998,0.4471) -- (0.0814,0.4485) -- (0.0815,0.4491) -- (0.0771,0.4493) -- (0.0767,0.4409) -- (0.1639,0.4319) -- (0.1640,0.4343) -- (0.1661,0.4342) -- (0.1660,0.4328) -- (0.1704,0.4323) -- (0.1707,0.4351) -- (0.1759,0.4346) -- (0.1747,0.4240) -- cycle;
\fill[color={rgb,255:red,57;green,155;blue,141}] (0.1671,0.3627) -- (0.1679,0.3697) -- (0.1531,0.3715) -- (0.1539,0.3782) -- (0.1688,0.3766) -- (0.1694,0.3837) -- (0.0439,0.3972) -- (0.0431,0.3894) -- (0.0589,0.3877) -- (0.0584,0.3819) -- (0.0424,0.3834) -- (0.0417,0.3758) -- cycle;
\fill[color={rgb,255:red,58;green,156;blue,142}] (0.1916,0.3598) -- (0.1921,0.3629) -- (0.1931,0.3629) -- (0.1935,0.3663) -- (0.1882,0.3668) -- (0.1895,0.3791) -- (0.1802,0.3802) -- (0.1798,0.3762) -- (0.1815,0.3761) -- (0.1815,0.3749) -- (0.1794,0.3752) -- (0.1792,0.3740) -- (0.1751,0.3743) -- (0.1737,0.3621) -- (0.1786,0.3616) -- (0.1789,0.3641) -- (0.1832,0.3637) -- (0.1830,0.3607) -- cycle;
\fill[color={rgb,255:red,55;green,153;blue,139}] (0.0498,0.3612) -- (0.0482,0.3437) -- (0.0773,0.3408) -- (0.0736,0.3023) -- (0.0768,0.3021) -- (0.0763,0.2960) -- (0.0788,0.2958) -- (0.0771,0.2784) -- (0.0935,0.2768) -- (0.0951,0.2953) -- (0.0790,0.2968) -- (0.0800,0.3066) -- (0.0775,0.3069) -- (0.0808,0.3403) -- (0.0928,0.3391) -- (0.0921,0.3313) -- (0.1042,0.3302) -- (0.1040,0.3268) -- (0.1060,0.3267) -- (0.1056,0.3218) -- (0.1126,0.3212) -- (0.1122,0.3161) -- (0.1250,0.3149) -- (0.1259,0.3274) -- (0.1086,0.3290) -- (0.1089,0.3323) -- (0.1020,0.3329) -- (0.1023,0.3385) -- (0.1269,0.3363) -- (0.1284,0.3532) -- cycle;
\fill[color={rgb,255:red,52;green,149;blue,136}] (0.0971,0.2920) -- (0.0962,0.2844) -- (0.1069,0.2832) -- (0.1060,0.2743) -- (0.1196,0.2730) -- (0.1210,0.2881) -- (0.1563,0.2845) -- (0.1581,0.3013) -- (0.1462,0.3025) -- (0.1464,0.3047) -- (0.1419,0.3052) -- (0.1417,0.3032) -- (0.1216,0.3053) -- (0.1214,0.3028) -- (0.1201,0.3028) -- (0.1191,0.2936) -- (0.1066,0.2949) -- (0.1066,0.2939) -- (0.1031,0.2941) -- (0.1028,0.2914) -- cycle;
\fill[color={rgb,255:red,51;green,147;blue,135}] (0.1221,0.2724) -- (0.1240,0.2722) -- (0.1238,0.2704) -- (0.0896,0.2740) -- (0.0893,0.2709) -- (0.0779,0.2717) -- (0.0774,0.2673) -- (0.1319,0.2614) -- (0.1309,0.2499) -- (0.1292,0.2500) -- (0.1294,0.2525) -- (0.1286,0.2527) -- (0.1291,0.2573) -- (0.1248,0.2577) -- (0.1239,0.2485) -- (0.1307,0.2479) -- (0.1304,0.2433) -- (0.1399,0.2423) -- (0.1398,0.2417) -- (0.1412,0.2415) -- (0.1408,0.2368) -- (0.1503,0.2358) -- (0.1503,0.2363) -- (0.1528,0.2362) -- (0.1528,0.2355) -- (0.1623,0.2344) -- (0.1633,0.2451) -- (0.1640,0.2451) -- (0.1647,0.2519) -- (0.1671,0.2518) -- (0.1654,0.2342) -- (0.1811,0.2325) -- (0.1844,0.2641) -- (0.1685,0.2658) -- (0.1682,0.2641) -- (0.1660,0.2644) -- (0.1661,0.2659) -- (0.1560,0.2669) -- (0.1551,0.2581) -- (0.1541,0.2583) -- (0.1537,0.2539) -- (0.1521,0.2541) -- (0.1535,0.2675) -- (0.1494,0.2679) -- (0.1495,0.2697) -- (0.1607,0.2686) -- (0.1613,0.2742) -- (0.1227,0.2783) -- cycle;
\fill[color={rgb,255:red,54;green,152;blue,138}] (0.0327,0.3360) -- (0.0309,0.3360) -- (0.0310,0.3377) -- (0.0262,0.3381) -- (0.0244,0.3200) -- (0.0053,0.3219) -- (0.0046,0.3146) -- (0.0086,0.3141) -- (0.0071,0.3011) -- (0.0131,0.3004) -- (0.0129,0.2974) -- (0.0193,0.2966) -- (0.0204,0.3094) -- (0.0367,0.3075) -- (0.0349,0.2909) -- (0.0408,0.2904) -- (0.0425,0.3073) -- (0.0698,0.3044) -- (0.0705,0.3109) -- (0.0289,0.3152) -- (0.0305,0.3304) -- (0.0412,0.3300) -- (0.0418,0.3369) -- (0.0329,0.3386) -- cycle;
\fill[color={rgb,255:red,53;green,150;blue,137}] (0.0000,0.2783) -- (0.0086,0.2775) -- (0.0090,0.2795) -- (0.0106,0.2794) -- (0.0092,0.2671) -- (0.0213,0.2660) -- (0.0240,0.2910) -- (0.0105,0.2924) -- (0.0112,0.2998) -- (0.0025,0.3008) -- cycle;
\fill[color={rgb,255:red,41;green,135;blue,125}] (0.6609,0.1694) -- (0.7210,0.1378) -- (0.7342,0.1626) -- (0.6741,0.1943) -- cycle;
\fill[color={rgb,255:red,40;green,134;blue,124}] (0.6574,0.1288) -- (0.6497,0.1323) -- (0.6575,0.1477) -- (0.6558,0.1487) -- (0.6589,0.1539) -- (0.6535,0.1569) -- (0.6390,0.1299) -- (0.6466,0.1263) -- (0.6481,0.1292) -- (0.6554,0.1251) -- (0.6472,0.1097) -- (0.6796,0.0926) -- (0.6993,0.1293) -- (0.6669,0.1464) -- cycle;
\fill[color={rgb,255:red,38;green,131;blue,122}] (0.7290,0.1457) -- (0.7503,0.1345) -- (0.7580,0.1492) -- (0.7371,0.1601) -- cycle;
\fill[color={rgb,255:red,39;green,132;blue,123}] (0.7022,0.1115) -- (0.7353,0.0939) -- (0.7335,0.0907) -- (0.7006,0.1081) -- (0.6951,0.0980) -- (0.7281,0.0805) -- (0.7259,0.0766) -- (0.6926,0.0942) -- (0.6869,0.0837) -- (0.7204,0.0661) -- (0.7184,0.0625) -- (0.6849,0.0801) -- (0.6796,0.0702) -- (0.7131,0.0525) -- (0.7107,0.0482) -- (0.6772,0.0660) -- (0.6720,0.0562) -- (0.7054,0.0385) -- (0.7031,0.0339) -- (0.6700,0.0515) -- (0.6626,0.0377) -- (0.6958,0.0202) -- (0.6939,0.0162) -- (0.6602,0.0339) -- (0.6553,0.0245) -- (0.6647,0.0195) -- (0.6650,0.0201) -- (0.6890,0.0074) -- (0.6919,0.0127) -- (0.6966,0.0106) -- (0.6988,0.0133) -- (0.7005,0.0179) -- (0.7345,0.0000) -- (0.7419,0.0135) -- (0.7078,0.0316) -- (0.7101,0.0359) -- (0.7444,0.0175) -- (0.7498,0.0275) -- (0.7154,0.0458) -- (0.7174,0.0494) -- (0.7515,0.0311) -- (0.7571,0.0418) -- (0.7229,0.0597) -- (0.7243,0.0628) -- (0.7581,0.0446) -- (0.7641,0.0558) -- (0.7302,0.0737) -- (0.7322,0.0770) -- (0.7658,0.0591) -- (0.7713,0.0697) -- (0.7377,0.0876) -- (0.7397,0.0913) -- (0.7729,0.0736) -- (0.7805,0.0880) -- (0.7475,0.1057) -- (0.7465,0.1092) -- (0.7426,0.1080) -- (0.7096,0.1251) -- cycle;
\fill[color={rgb,255:red,37;green,129;blue,121}] (0.7484,0.1233) -- (0.8684,0.0595) -- (0.8817,0.0844) -- (0.7613,0.1482) -- cycle;
\fill[color={rgb,255:red,32;green,123;blue,116}] (0.9905,0.2697) -- (0.9988,0.2854) -- (0.9694,0.3006) -- (0.9709,0.3036) -- (1.0072,0.2844) -- (1.0096,0.2895) -- (0.9956,0.2971) -- (0.9982,0.3017) -- (1.0001,0.3007) -- (1.0021,0.3048) -- (1.0003,0.3057) -- (1.0057,0.3159) -- (0.9799,0.3293) -- (0.9822,0.3338) -- (0.9386,0.3570) -- (0.9358,0.3517) -- (0.9499,0.3443) -- (0.9278,0.3029) -- cycle;
\fill[color={rgb,255:red,45;green,140;blue,129}] (0.6081,0.4733) -- (0.5852,0.4310) -- (0.5862,0.4193) -- (0.5901,0.4119) -- (0.5976,0.4057) -- (0.6725,0.3636) -- (0.6822,0.3599) -- (0.6920,0.3619) -- (0.7034,0.3746) -- (0.7224,0.4127) -- cycle;
\fill[color={rgb,255:red,30;green,116;blue,110}] (0.7319,0.6646) -- (0.7281,0.6577) -- (0.7402,0.6508) -- (0.7377,0.6466) -- (0.7260,0.6532) -- (0.7219,0.6462) -- (0.7557,0.6277) -- (0.7593,0.6348) -- (0.7472,0.6415) -- (0.7496,0.6459) -- (0.7620,0.6396) -- (0.7656,0.6468) -- (0.7514,0.6542) -- (0.7521,0.6557) -- (0.7665,0.6484) -- (0.7674,0.6498) -- (0.7726,0.6479) -- (0.7716,0.6460) -- (0.7753,0.6444) -- (0.7767,0.6471) -- (0.7813,0.6453) -- (0.7801,0.6434) -- (0.8059,0.6298) -- (0.8113,0.6399) -- (0.7855,0.6536) -- (0.7822,0.6476) -- (0.7782,0.6494) -- (0.8038,0.6976) -- (0.8252,0.6865) -- (0.8237,0.6838) -- (0.8156,0.6879) -- (0.8110,0.6794) -- (0.8288,0.6705) -- (0.8170,0.6475) -- (0.8256,0.6431) -- (0.8417,0.6745) -- (0.8256,0.6824) -- (0.8290,0.6879) -- (0.8454,0.6789) -- (0.8492,0.6860) -- (0.8372,0.6926) -- (0.8392,0.6969) -- (0.8345,0.6994) -- (0.8351,0.7009) -- (0.8380,0.6995) -- (0.8452,0.7131) -- (0.8535,0.7087) -- (0.8460,0.6950) -- (0.8543,0.6905) -- (0.8565,0.6950) -- (0.8552,0.6960) -- (0.8639,0.7121) -- (0.8475,0.7207) -- (0.8458,0.7172) -- (0.8418,0.7194) -- (0.8435,0.7229) -- (0.8269,0.7315) -- (0.8163,0.7106) -- (0.8339,0.7014) -- (0.8332,0.7000) -- (0.8301,0.7016) -- (0.8277,0.6973) -- (0.8154,0.7039) -- (0.8116,0.6969) -- (0.8267,0.6889) -- (0.8254,0.6872) -- (0.8042,0.6986) -- (0.8226,0.7336) -- (0.8187,0.7354) -- (0.8164,0.7312) -- (0.8142,0.7324) -- (0.8157,0.7355) -- (0.7765,0.7558) -- (0.7706,0.7443) -- (0.8096,0.7239) -- (0.8135,0.7315) -- (0.8159,0.7302) -- (0.8013,0.7022) -- (0.7955,0.7052) -- (0.7974,0.7096) -- (0.7824,0.7176) -- (0.7802,0.7132) -- (0.7762,0.7156) -- (0.7786,0.7196) -- (0.7632,0.7279) -- (0.7577,0.7170) -- (0.7718,0.7096) -- (0.7681,0.7030) -- (0.7744,0.6994) -- (0.7778,0.7061) -- (0.7920,0.6985) -- (0.7947,0.7041) -- (0.8007,0.7010) -- (0.7730,0.6488) -- (0.7480,0.6606) -- (0.7460,0.6571) -- cycle;
\fill[color={rgb,255:red,31;green,119;blue,113}] (0.8465,0.5338) -- (0.8528,0.5459) -- (0.8500,0.5472) -- (0.8567,0.5602) -- (0.8451,0.5664) -- (0.8468,0.5694) -- (0.8437,0.5709) -- (0.8443,0.5716) -- (0.8400,0.5738) -- (0.8381,0.5700) -- (0.8335,0.5725) -- (0.8367,0.5783) -- (0.8593,0.5666) -- (0.8643,0.5761) -- (0.8412,0.5885) -- (0.8454,0.5963) -- (0.8393,0.5997) -- (0.8421,0.6039) -- (0.8655,0.5916) -- (0.8612,0.5831) -- (0.8685,0.5787) -- (0.8774,0.5963) -- (0.8458,0.6128) -- (0.8280,0.5804) -- (0.8219,0.5838) -- (0.8275,0.5944) -- (0.8078,0.6046) -- (0.8122,0.6133) -- (0.8140,0.6127) -- (0.8152,0.6150) -- (0.8137,0.6157) -- (0.8148,0.6183) -- (0.8268,0.6121) -- (0.8278,0.6141) -- (0.8295,0.6132) -- (0.8315,0.6173) -- (0.8298,0.6180) -- (0.8311,0.6203) -- (0.8106,0.6310) -- (0.7950,0.6016) -- (0.7897,0.6010) -- (0.7756,0.6083) -- (0.7809,0.6186) -- (0.7960,0.6104) -- (0.8006,0.6193) -- (0.7769,0.6319) -- (0.7622,0.6039) -- (0.7854,0.5922) -- (0.7897,0.5997) -- (0.7954,0.6004) -- (0.8118,0.5917) -- (0.8085,0.5860) -- (0.8199,0.5797) -- (0.8213,0.5823) -- (0.8335,0.5756) -- (0.8326,0.5726) -- (0.8213,0.5507) -- (0.8303,0.5460) -- (0.8288,0.5429) -- (0.8323,0.5415) -- (0.8336,0.5441) -- (0.8453,0.5381) -- (0.8437,0.5353) -- cycle;
\fill[color={rgb,255:red,30;green,113;blue,108}] (0.6848,0.5889) -- (0.7115,0.5746) -- (0.7174,0.5856) -- (0.7228,0.5828) -- (0.7170,0.5719) -- (0.7440,0.5577) -- (0.7500,0.5698) -- (0.7550,0.5667) -- (0.7488,0.5548) -- (0.7759,0.5404) -- (0.7904,0.5676) -- (0.7631,0.5815) -- (0.7568,0.5694) -- (0.7517,0.5722) -- (0.7581,0.5845) -- (0.7312,0.5985) -- (0.7240,0.5853) -- (0.7188,0.5882) -- (0.7260,0.6016) -- (0.6989,0.6159) -- cycle;
\fill[color={rgb,255:red,27;green,94;blue,94}] (0.4705,0.5390) -- (0.4499,0.5015) -- (0.5495,0.4495) -- (0.5699,0.4868) -- (0.5234,0.5114) -- (0.5289,0.5201) -- (0.5640,0.5015) -- (0.5693,0.5106) -- (0.5342,0.5295) -- (0.5381,0.5354) -- (0.5726,0.5171) -- (0.5782,0.5270) -- (0.5843,0.5380) -- (0.5480,0.5571) -- (0.5155,0.5745) -- (0.5132,0.5695) -- (0.5021,0.5750) -- (0.5026,0.5759) -- (0.5037,0.5762) -- (0.5048,0.5798) -- (0.5017,0.5810) -- (0.4994,0.5765) -- (0.4944,0.5790) -- (0.4779,0.5483) -- (0.4899,0.5424) -- (0.5059,0.5707) -- (0.5122,0.5678) -- (0.4955,0.5357) -- (0.5249,0.5202) -- (0.5211,0.5131) -- cycle;
\fill[color={rgb,255:red,25;green,79;blue,82}] (0.4269,0.7299) -- (0.4241,0.7244) -- (0.4410,0.7158) -- (0.4438,0.7209) -- (0.4470,0.7195) -- (0.4438,0.7132) -- (0.4594,0.7049) -- (0.4715,0.7273) -- (0.4781,0.7241) -- (0.4670,0.7023) -- (0.4768,0.6975) -- (0.4884,0.7199) -- (0.4720,0.7284) -- (0.4805,0.7444) -- (0.4831,0.7431) -- (0.4770,0.7310) -- (0.4911,0.7239) -- (0.4977,0.7366) -- (0.4808,0.7451) -- (0.4819,0.7470) -- (0.4983,0.7392) -- (0.5165,0.7743) -- (0.5033,0.7809) -- (0.4857,0.7465) -- (0.4824,0.7480) -- (0.4999,0.7808) -- (0.4588,0.8025) -- (0.4280,0.7451) -- (0.4445,0.7364) -- (0.4559,0.7578) -- (0.4651,0.7529) -- (0.4472,0.7197) -- (0.4427,0.7214) -- cycle;
\fill[color={rgb,255:red,31;green,122;blue,115}] (1.0409,0.3563) -- (1.0448,0.3634) -- (1.0413,0.3651) -- (1.0477,0.3773) -- (1.0405,0.3810) -- (1.0423,0.3847) -- (1.0337,0.3892) -- (1.0318,0.3856) -- (1.0211,0.3913) -- (1.0136,0.3769) -- (1.0068,0.3804) -- (1.0046,0.3763) -- (1.0034,0.3773) -- (1.0263,0.4209) -- (1.0368,0.4152) -- (1.0299,0.4022) -- (1.0520,0.3906) -- (1.0575,0.4021) -- (1.0601,0.4008) -- (1.0717,0.4228) -- (1.0476,0.4357) -- (1.0372,0.4161) -- (1.0266,0.4214) -- (1.0271,0.4222) -- (1.0288,0.4213) -- (1.0314,0.4264) -- (0.9877,0.4497) -- (0.9854,0.4442) -- (0.9922,0.4409) -- (0.9900,0.4367) -- (0.9789,0.4426) -- (0.9616,0.4096) -- (0.9942,0.3927) -- (1.0053,0.4140) -- (1.0088,0.4121) -- (1.0080,0.4105) -- (1.0123,0.4083) -- (1.0170,0.4177) -- (1.0131,0.4198) -- (1.0124,0.4185) -- (1.0088,0.4205) -- (1.0137,0.4293) -- (1.0212,0.4254) -- (1.0192,0.4217) -- (1.0231,0.4197) -- (1.0251,0.4233) -- (1.0262,0.4226) -- (1.0019,0.3761) -- (1.0043,0.3753) -- cycle;
\fill[color={rgb,255:red,23;green,63;blue,70}] (0.2496,0.9040) -- (0.2735,0.8921) -- (0.2642,0.8741) -- (0.2405,0.8861) -- cycle;
\begin{scope}[transparency group,opacity=0.7]
\draw[color={rgb,255:red,245;green,154;blue,74},line width=1.8pt] (0.3257,0.4503)--(0.2972,0.4446);
\draw[color={rgb,255:red,245;green,153;blue,73},line width=1.8pt] (0.2972,0.4446)--(0.2686,0.4389);
\draw[color={rgb,255:red,244;green,153;blue,73},line width=1.8pt] (0.2686,0.4389)--(0.2401,0.4332);
\draw[color={rgb,255:red,244;green,152;blue,73},line width=1.8pt] (0.2401,0.4332)--(0.2116,0.4275);
\draw[color={rgb,255:red,244;green,151;blue,72},line width=1.8pt] (0.2116,0.4275)--(0.1959,0.4039);
\draw[color={rgb,255:red,243;green,151;blue,72},line width=1.8pt] (0.1959,0.4039)--(0.1802,0.3802);
\draw[color={rgb,255:red,243;green,150;blue,71},line width=1.8pt] (0.1802,0.3802)--(0.1542,0.3654);
\draw[color={rgb,255:red,243;green,150;blue,71},line width=1.8pt] (0.1542,0.3654)--(0.1281,0.3505);
\draw[color={rgb,255:red,242;green,149;blue,71},line width=1.8pt] (0.1281,0.3505)--(0.1021,0.3356);
\draw[color={rgb,255:red,242;green,149;blue,70},line width=1.8pt] (0.1021,0.3356)--(0.0761,0.3207);
\draw[color={rgb,255:red,241;green,148;blue,70},line width=1.8pt] (0.0761,0.3207)--(0.0500,0.3058);
\draw[color={rgb,255:red,241;green,148;blue,69},line width=1.8pt] (0.0500,0.3058)--(0.0240,0.2910);
\draw[color={rgb,255:red,241;green,147;blue,69},line width=1.8pt] (0.0240,0.2910)--(0.0487,0.2878);
\draw[color={rgb,255:red,241;green,147;blue,69},line width=1.8pt] (0.0487,0.2878)--(0.0734,0.2846);
\draw[color={rgb,255:red,240;green,146;blue,68},line width=1.8pt] (0.0734,0.2846)--(0.0980,0.2815);
\draw[color={rgb,255:red,240;green,146;blue,68},line width=1.8pt] (0.0980,0.2815)--(0.1227,0.2783);
\draw[color={rgb,255:red,240;green,145;blue,68},line width=1.8pt] (0.1227,0.2783)--(0.1503,0.2812);
\draw[color={rgb,255:red,239;green,145;blue,67},line width=1.8pt] (0.1503,0.2812)--(0.1779,0.2841);
\draw[color={rgb,255:red,239;green,144;blue,67},line width=1.8pt] (0.1779,0.2841)--(0.2055,0.2869);
\draw[color={rgb,255:red,239;green,143;blue,66},line width=1.8pt] (0.2055,0.2869)--(0.2331,0.2898);
\draw[color={rgb,255:red,238;green,143;blue,66},line width=1.8pt] (0.2331,0.2898)--(0.2608,0.2927);
\draw[color={rgb,255:red,238;green,142;blue,66},line width=1.8pt] (0.2608,0.2927)--(0.2884,0.2956);
\draw[color={rgb,255:red,238;green,142;blue,65},line width=1.8pt] (0.2884,0.2956)--(0.3160,0.2985);
\draw[color={rgb,255:red,237;green,141;blue,65},line width=1.8pt] (0.3160,0.2985)--(0.3436,0.3013);
\draw[color={rgb,255:red,237;green,141;blue,65},line width=1.8pt] (0.3436,0.3013)--(0.3712,0.3042);
\draw[color={rgb,255:red,237;green,140;blue,64},line width=1.8pt] (0.3712,0.3042)--(0.3988,0.3071);
\draw[color={rgb,255:red,237;green,140;blue,64},line width=1.8pt] (0.3988,0.3071)--(0.4264,0.3100);
\draw[color={rgb,255:red,236;green,139;blue,63},line width=1.8pt] (0.4264,0.3100)--(0.4540,0.3129);
\draw[color={rgb,255:red,236;green,139;blue,63},line width=1.8pt] (0.4540,0.3129)--(0.4634,0.3317);
\draw[color={rgb,255:red,236;green,138;blue,63},line width=1.8pt] (0.4634,0.3317)--(0.4728,0.3506);
\draw[color={rgb,255:red,235;green,138;blue,62},line width=1.8pt] (0.4728,0.3506)--(0.4997,0.3569);
\draw[color={rgb,255:red,235;green,137;blue,62},line width=1.8pt] (0.4997,0.3569)--(0.5266,0.3632);
\draw[color={rgb,255:red,235;green,137;blue,62},line width=1.8pt] (0.5266,0.3632)--(0.5525,0.3692);
\draw[color={rgb,255:red,235;green,136;blue,61},line width=1.8pt] (0.5525,0.3692)--(0.5785,0.3752);
\draw[color={rgb,255:red,234;green,136;blue,61},line width=1.8pt] (0.5785,0.3752)--(0.6044,0.3813);
\draw[color={rgb,255:red,234;green,135;blue,61},line width=1.8pt] (0.6044,0.3813)--(0.6303,0.3873);
\draw[color={rgb,255:red,234;green,135;blue,60},line width=1.8pt] (0.6303,0.3873)--(0.6385,0.3622);
\draw[color={rgb,255:red,233;green,134;blue,60},line width=1.8pt] (0.6385,0.3622)--(0.6468,0.3370);
\draw[color={rgb,255:red,233;green,134;blue,59},line width=1.8pt] (0.6468,0.3370)--(0.6550,0.3118);
\draw[color={rgb,255:red,233;green,133;blue,59},line width=1.8pt] (0.6550,0.3118)--(0.6632,0.2867);
\draw[color={rgb,255:red,232;green,133;blue,59},line width=1.8pt] (0.6632,0.2867)--(0.6714,0.2615);
\draw[color={rgb,255:red,232;green,132;blue,58},line width=1.8pt] (0.6714,0.2615)--(0.6797,0.2363);
\draw[color={rgb,255:red,232;green,132;blue,58},line width=1.8pt] (0.6797,0.2363)--(0.6846,0.2096);
\draw[color={rgb,255:red,231;green,131;blue,58},line width=1.8pt] (0.6846,0.2096)--(0.6895,0.1828);
\draw[color={rgb,255:red,231;green,131;blue,57},line width=1.8pt] (0.6895,0.1828)--(0.6944,0.1561);
\draw[color={rgb,255:red,231;green,130;blue,57},line width=1.8pt] (0.6944,0.1561)--(0.6993,0.1293);
\draw[color={rgb,255:red,231;green,130;blue,57},line width=1.8pt] (0.6993,0.1293)--(0.7096,0.1251);
\draw[color={rgb,255:red,230;green,129;blue,56},line width=1.8pt] (0.7096,0.1251)--(0.7355,0.1366);
\draw[color={rgb,255:red,230;green,129;blue,56},line width=1.8pt] (0.7355,0.1366)--(0.7613,0.1482);
\draw[color={rgb,255:red,230;green,128;blue,56},line width=1.8pt] (0.7613,0.1482)--(0.7867,0.1599);
\draw[color={rgb,255:red,229;green,128;blue,55},line width=1.8pt] (0.7867,0.1599)--(0.8121,0.1717);
\draw[color={rgb,255:red,229;green,127;blue,55},line width=1.8pt] (0.8121,0.1717)--(0.8346,0.1794);
\draw[color={rgb,255:red,229;green,127;blue,55},line width=1.8pt] (0.8346,0.1794)--(0.8572,0.1872);
\draw[color={rgb,255:red,229;green,126;blue,54},line width=1.8pt] (0.8572,0.1872)--(0.8797,0.1949);
\draw[color={rgb,255:red,228;green,126;blue,54},line width=1.8pt] (0.8797,0.1949)--(0.8910,0.2150);
\draw[color={rgb,255:red,228;green,126;blue,54},line width=1.8pt] (0.8910,0.2150)--(0.9022,0.2350);
\draw[color={rgb,255:red,228;green,125;blue,53},line width=1.8pt] (0.9022,0.2350)--(0.9135,0.2550);
\draw[color={rgb,255:red,228;green,125;blue,53},line width=1.8pt] (0.9135,0.2550)--(0.9215,0.2808);
\draw[color={rgb,255:red,227;green,124;blue,53},line width=1.8pt] (0.9215,0.2808)--(0.9295,0.3066);
\draw[color={rgb,255:red,227;green,124;blue,52},line width=1.8pt] (0.9295,0.3066)--(0.9375,0.3323);
\draw[color={rgb,255:red,227;green,123;blue,52},line width=1.8pt] (0.9375,0.3323)--(0.9456,0.3581);
\draw[color={rgb,255:red,226;green,123;blue,51},line width=1.8pt] (0.9456,0.3581)--(0.9536,0.3839);
\draw[color={rgb,255:red,226;green,122;blue,51},line width=1.8pt] (0.9536,0.3839)--(0.9616,0.4096);
\draw[color={rgb,255:red,225;green,121;blue,51},line width=1.8pt] (0.9616,0.4096)--(0.9410,0.4315);
\draw[color={rgb,255:red,225;green,121;blue,50},line width=1.8pt] (0.9410,0.4315)--(0.9204,0.4533);
\draw[color={rgb,255:red,224;green,120;blue,50},line width=1.8pt] (0.9204,0.4533)--(0.8998,0.4751);
\draw[color={rgb,255:red,224;green,119;blue,49},line width=1.8pt] (0.8998,0.4751)--(0.8792,0.4969);
\draw[color={rgb,255:red,223;green,118;blue,49},line width=1.8pt] (0.8792,0.4969)--(0.8586,0.5187);
\draw[color={rgb,255:red,223;green,118;blue,48},line width=1.8pt] (0.8586,0.5187)--(0.8380,0.5405);
\draw[color={rgb,255:red,222;green,117;blue,48},line width=1.8pt] (0.8380,0.5405)--(0.8175,0.5623);
\draw[color={rgb,255:red,221;green,116;blue,47},line width=1.8pt] (0.8175,0.5623)--(0.7969,0.5841);
\draw[color={rgb,255:red,221;green,115;blue,47},line width=1.8pt] (0.7969,0.5841)--(0.7763,0.6059);
\draw[color={rgb,255:red,220;green,115;blue,47},line width=1.8pt] (0.7763,0.6059)--(0.7557,0.6277);
\draw[color={rgb,255:red,220;green,114;blue,46},line width=1.8pt] (0.7557,0.6277)--(0.7459,0.6028);
\draw[color={rgb,255:red,219;green,113;blue,46},line width=1.8pt] (0.7459,0.6028)--(0.7360,0.5778);
\draw[color={rgb,255:red,219;green,112;blue,45},line width=1.8pt] (0.7360,0.5778)--(0.7262,0.5529);
\draw[color={rgb,255:red,218;green,112;blue,45},line width=1.8pt] (0.7262,0.5529)--(0.6984,0.5584);
\draw[color={rgb,255:red,218;green,111;blue,44},line width=1.8pt] (0.6984,0.5584)--(0.6783,0.5772);
\draw[color={rgb,255:red,217;green,110;blue,44},line width=1.8pt] (0.6783,0.5772)--(0.6582,0.5959);
\draw[color={rgb,255:red,217;green,110;blue,44},line width=1.8pt] (0.6582,0.5959)--(0.6382,0.6146);
\draw[color={rgb,255:red,216;green,109;blue,43},line width=1.8pt] (0.6382,0.6146)--(0.6140,0.6064);
\draw[color={rgb,255:red,216;green,108;blue,43},line width=1.8pt] (0.6140,0.6064)--(0.5899,0.5982);
\draw[color={rgb,255:red,215;green,108;blue,43},line width=1.8pt] (0.5899,0.5982)--(0.5657,0.5899);
\draw[color={rgb,255:red,215;green,107;blue,42},line width=1.8pt] (0.5657,0.5899)--(0.5416,0.5817);
\draw[color={rgb,255:red,214;green,106;blue,42},line width=1.8pt] (0.5416,0.5817)--(0.5175,0.5735);
\draw[color={rgb,255:red,214;green,106;blue,41},line width=1.8pt] (0.5175,0.5735)--(0.5138,0.5874);
\draw[color={rgb,255:red,213;green,105;blue,41},line width=1.8pt] (0.5138,0.5874)--(0.5036,0.6153);
\draw[color={rgb,255:red,213;green,105;blue,41},line width=1.8pt] (0.5036,0.6153)--(0.5098,0.6407);
\draw[color={rgb,255:red,212;green,104;blue,40},line width=1.8pt] (0.5098,0.6407)--(0.5161,0.6661);
\draw[color={rgb,255:red,212;green,103;blue,40},line width=1.8pt] (0.5161,0.6661)--(0.5162,0.6931);
\draw[color={rgb,255:red,211;green,103;blue,40},line width=1.8pt] (0.5162,0.6931)--(0.5163,0.7202);
\draw[color={rgb,255:red,211;green,102;blue,39},line width=1.8pt] (0.5163,0.7202)--(0.5164,0.7472);
\draw[color={rgb,255:red,210;green,101;blue,39},line width=1.8pt] (0.5164,0.7472)--(0.5165,0.7743);
\draw[color={rgb,255:red,210;green,101;blue,38},line width=1.8pt] (0.5165,0.7743)--(0.5290,0.7952);
\draw[color={rgb,255:red,209;green,100;blue,38},line width=1.8pt] (0.5290,0.7952)--(0.5415,0.8161);
\draw[color={rgb,255:red,209;green,99;blue,38},line width=1.8pt] (0.5415,0.8161)--(0.5190,0.8300);
\draw[color={rgb,255:red,208;green,99;blue,37},line width=1.8pt] (0.5190,0.8300)--(0.4966,0.8439);
\draw[color={rgb,255:red,208;green,98;blue,37},line width=1.8pt] (0.4966,0.8439)--(0.4742,0.8578);
\draw[color={rgb,255:red,207;green,97;blue,36},line width=1.8pt] (0.4742,0.8578)--(0.4518,0.8717);
\draw[color={rgb,255:red,207;green,97;blue,36},line width=1.8pt] (0.4518,0.8717)--(0.4293,0.8856);
\draw[color={rgb,255:red,206;green,96;blue,36},line width=1.8pt] (0.4293,0.8856)--(0.4069,0.8995);
\draw[color={rgb,255:red,206;green,95;blue,35},line width=1.8pt] (0.4069,0.8995)--(0.3845,0.9134);
\draw[color={rgb,255:red,205;green,95;blue,35},line width=1.8pt] (0.3845,0.9134)--(0.3620,0.9273);
\draw[color={rgb,255:red,205;green,94;blue,34},line width=1.8pt] (0.3620,0.9273)--(0.3376,0.9140);
\draw[color={rgb,255:red,204;green,93;blue,34},line width=1.8pt] (0.3376,0.9140)--(0.3131,0.9007);
\draw[color={rgb,255:red,204;green,93;blue,34},line width=1.8pt] (0.3131,0.9007)--(0.2887,0.8874);
\draw[color={rgb,255:red,203;green,92;blue,33},line width=1.8pt] (0.2887,0.8874)--(0.2642,0.8741);
\draw[color={rgb,255:red,203;green,91;blue,33},line width=1.8pt] (0.2642,0.8741)--(0.2642,0.8447);
\draw[color={rgb,255:red,202;green,91;blue,32},line width=1.8pt] (0.2642,0.8447)--(0.2641,0.8152);
\draw[color={rgb,255:red,201;green,90;blue,32},line width=1.8pt] (0.2641,0.8152)--(0.2640,0.7857);
\draw[color={rgb,255:red,201;green,89;blue,31},line width=1.8pt] (0.2640,0.7857)--(0.2640,0.7562);
\draw[color={rgb,255:red,200;green,88;blue,31},line width=1.8pt] (0.2640,0.7562)--(0.2571,0.7346);
\draw[color={rgb,255:red,200;green,88;blue,31},line width=1.8pt] (0.2571,0.7346)--(0.2502,0.7129);
\draw[color={rgb,255:red,200;green,87;blue,30},line width=1.8pt] (0.2502,0.7129)--(0.2433,0.6912);
\draw[color={rgb,255:red,199;green,87;blue,30},line width=1.8pt] (0.2433,0.6912)--(0.2363,0.6695);
\draw[color={rgb,255:red,199;green,86;blue,30},line width=1.8pt] (0.2363,0.6695)--(0.2500,0.6445);
\draw[color={rgb,255:red,198;green,85;blue,29},line width=1.8pt] (0.2500,0.6445)--(0.2637,0.6195);
\draw[color={rgb,255:red,198;green,85;blue,29},line width=1.8pt] (0.2637,0.6195)--(0.2773,0.5945);
\draw[color={rgb,255:red,197;green,84;blue,28},line width=1.8pt] (0.2773,0.5945)--(0.2910,0.5696);
\draw[color={rgb,255:red,196;green,83;blue,28},line width=1.8pt] (0.2910,0.5696)--(0.3047,0.5446);
\draw[color={rgb,255:red,196;green,82;blue,27},line width=1.8pt] (0.3047,0.5446)--(0.3183,0.5196);
\draw[color={rgb,255:red,195;green,82;blue,27},line width=1.8pt] (0.3183,0.5196)--(0.3320,0.4946);
\draw[color={rgb,255:red,195;green,81;blue,27},line width=1.8pt] (0.3320,0.4946)--(0.3457,0.4696);
\draw[color={rgb,255:red,194;green,80;blue,26},line width=1.8pt] (0.3457,0.4696)--(0.3257,0.4503);
\end{scope}
\fill[white,draw=Ink,line width=1pt] (0.3257,0.4503) circle (0.012);

      \node[anchor=south west,inner sep=1pt,font=\sffamily\bfseries\sz{24}{28},text=Ink] (st) at (0.035,0.80) {$s=t$};\draw[Muted,line width=0.8pt] (st.south east)--(\uspsx,\uspsy);
    \end{tikzpicture}};
  \node[anchor=south west,inner sep=0.28cm,fill=Pale,rounded corners=4pt,text width=11.0cm] at ([xshift=1.2cm,yshift=1.6cm]current page.south west)
    {\sz{20}{24}\muted\textbf{Exemplo:} drone do {\bfseries\color{GoldText}IME} e 51 regiões da USP\par\vspace{0.12cm}
     {\bfseries\sz{32}{36}\hi $\approx 10^{104}$ rotas possíveis}\par{\sz{22}{26}\ink Testar todas é inviável.\par}\vspace{0.1cm}
     {\sz{24}{28}\ink Nosso solver: rota ótima em \textbf{0,23 s}\par}};
  \footerbar{2}
\end{tikzpicture}
\newpage

%%%%%%%%%%%%%%%%%%%%%%%%%%%%%%%% 3. Método (busca real: caso 17, 5 regiões)
\def\traceU{1{,}00}
\def\traceLroot{0{,}96}\def\traceLb{0{,}99}\def\traceLc{1{,}42}\def\traceLd{1{,}00}
\def\miniRoot{%
\filldraw[fill=Rule!60,draw=Ink,line width=0.5pt,line join=round] (4.0379,0.3541) -- (4.3991,1.6275) -- (3.6586,1.8453) -- (3.2601,0.5474) -- cycle;
\filldraw[fill=Rule!60,draw=Ink,line width=0.5pt,line join=round] (2.7222,0.2204) -- (3.3852,1.7663) -- (2.8931,1.9747) -- (2.2298,0.4291) -- cycle;
\filldraw[fill=Rule!60,draw=Ink,line width=0.5pt,line join=round] (2.1055,0.5668) -- (2.2813,1.0260) -- (1.9458,1.1528) -- (1.7700,0.6936) -- cycle;
\filldraw[fill=Rule!60,draw=Ink,line width=0.5pt,line join=round] (1.8882,0.0003) -- (2.0938,0.5482) -- (1.7872,0.6613) -- (1.5127,0.0000) -- cycle;
\filldraw[fill=Rule!60,draw=Ink,line width=0.5pt,line join=round] (0.2997,0.4998) -- (0.4500,1.1697) -- (0.1506,1.2361) -- (0.0000,0.5662) -- cycle;
\draw[Accent,line width=1.8pt,line join=round,line cap=round] (0.0000,-0.0000) -- (4.3991,1.9747);
\fill[white,draw=Ink,line width=0.8pt] (0.0000,0.0000) circle (0.07);
\fill[Ink] (4.3991,1.9747) circle (0.07);
}
\def\miniB{%
\filldraw[fill=Rule!60,draw=Ink,line width=0.5pt,line join=round] (4.0379,0.3541) -- (4.3991,1.6275) -- (3.6586,1.8453) -- (3.2601,0.5474) -- cycle;
\filldraw[fill=Rule!60,draw=Ink,line width=0.5pt,line join=round] (2.7222,0.2204) -- (3.3852,1.7663) -- (2.8931,1.9747) -- (2.2298,0.4291) -- cycle;
\filldraw[fill=Rule!60,draw=Ink,line width=0.5pt,line join=round] (2.1055,0.5668) -- (2.2813,1.0260) -- (1.9458,1.1528) -- (1.7700,0.6936) -- cycle;
\filldraw[fill=Rule!60,draw=Ink,line width=0.5pt,line join=round] (1.8882,0.0003) -- (2.0938,0.5482) -- (1.7872,0.6613) -- (1.5127,0.0000) -- cycle;
\filldraw[fill=Teal,draw=Ink,line width=0.5pt,line join=round] (0.2997,0.4998) -- (0.4500,1.1697) -- (0.1506,1.2361) -- (0.0000,0.5662) -- cycle;
\draw[Accent,line width=1.8pt,line join=round,line cap=round] (0.0000,-0.0000) -- (0.2997,0.4998) -- (4.3991,1.9747);
\fill[white,draw=Ink,line width=0.8pt] (0.0000,0.0000) circle (0.07);
\fill[Ink] (4.3991,1.9747) circle (0.07);
}
\def\miniC{%
\filldraw[fill=Rule!60,draw=Ink,line width=0.5pt,line join=round] (4.0379,0.3541) -- (4.3991,1.6275) -- (3.6586,1.8453) -- (3.2601,0.5474) -- cycle;
\filldraw[fill=Rule!60,draw=Ink,line width=0.5pt,line join=round] (2.7222,0.2204) -- (3.3852,1.7663) -- (2.8931,1.9747) -- (2.2298,0.4291) -- cycle;
\filldraw[fill=Rule!60,draw=Ink,line width=0.5pt,line join=round] (2.1055,0.5668) -- (2.2813,1.0260) -- (1.9458,1.1528) -- (1.7700,0.6936) -- cycle;
\filldraw[fill=Teal,draw=Ink,line width=0.5pt,line join=round] (1.8882,0.0003) -- (2.0938,0.5482) -- (1.7872,0.6613) -- (1.5127,0.0000) -- cycle;
\filldraw[fill=Teal,draw=Ink,line width=0.5pt,line join=round] (0.2997,0.4998) -- (0.4500,1.1697) -- (0.1506,1.2361) -- (0.0000,0.5662) -- cycle;
\draw[Muted,dashed,line width=1.8pt,line join=round,line cap=round] (0.0000,0.0000) -- (1.5127,0.0000) -- (0.3965,0.9311) -- (4.3991,1.9747);
\fill[white,draw=Ink,line width=0.8pt] (0.0000,0.0000) circle (0.07);
\fill[Ink] (4.3991,1.9747) circle (0.07);
}
\def\miniD{%
\filldraw[fill=Teal!35,draw=Ink,line width=0.5pt,line join=round] (4.0379,0.3541) -- (4.3991,1.6275) -- (3.6586,1.8453) -- (3.2601,0.5474) -- cycle;
\filldraw[fill=Teal!35,draw=Ink,line width=0.5pt,line join=round] (2.7222,0.2204) -- (3.3852,1.7663) -- (2.8931,1.9747) -- (2.2298,0.4291) -- cycle;
\filldraw[fill=Teal!35,draw=Ink,line width=0.5pt,line join=round] (2.1055,0.5668) -- (2.2813,1.0260) -- (1.9458,1.1528) -- (1.7700,0.6936) -- cycle;
\filldraw[fill=Teal,draw=Ink,line width=0.5pt,line join=round] (1.8882,0.0003) -- (2.0938,0.5482) -- (1.7872,0.6613) -- (1.5127,0.0000) -- cycle;
\filldraw[fill=Teal,draw=Ink,line width=0.5pt,line join=round] (0.2997,0.4998) -- (0.4500,1.1697) -- (0.1506,1.2361) -- (0.0000,0.5662) -- cycle;
\draw[Accent,line width=1.8pt,line join=round,line cap=round] (0.0000,-0.0000) -- (0.2997,0.4998) -- (1.7872,0.6613) -- (4.3991,1.9747);
\fill[white,draw=Ink,line width=0.8pt] (0.0000,0.0000) circle (0.07);
\fill[Ink] (4.3991,1.9747) circle (0.07);
}

\newcommand{\mapbox}[2]{\scalebox{0.85}{\begin{tikzpicture}\draw[#2,rounded corners=3pt,fill=white] (-0.15,-0.15) rectangle (4.55,2.13);\begin{scope}#1\end{scope}\end{tikzpicture}}}
\newcommand{\stepdot}[1]{\node[circle,fill=Accent,text=white,minimum size=0.8cm,inner sep=0,font=\sffamily\bfseries\sz{22}{24}] at #1}
\begin{tikzpicture}[remember picture,overlay,
    mframe/.style={draw=Rule,line width=1.2pt},
    mok/.style={draw=Teal,line width=2.4pt},
    mcut/.style={draw=Muted!60,line width=1.2pt,dashed},
    ed/.style={line width=1.4pt,Ink,-{Stealth[length=3mm]}},
    trcap/.style={font=\sffamily\sz{20}{24},inner sep=0,align=center}]
  \fill[white] (current page.south west) rectangle (current page.north east);
  \slidetitle{Como evitar testar todas as rotas?}
  % busca real (caso de 5 regiões): raiz -> um passo -> {descartado, ótimo}
  \node[anchor=north west,inner sep=0] (nr) at ([xshift=1.0cm,yshift=-4.2cm]current page.north west) {\mapbox{\miniRoot}{mframe}};
  \node[anchor=north west,inner sep=0] (nb) at ([xshift=5.9cm,yshift=-4.2cm]current page.north west) {\mapbox{\miniB}{mframe}};
  \node[anchor=north west,inner sep=0] (nc) at ([xshift=10.8cm,yshift=-2.55cm]current page.north west) {\mapbox{\miniC}{mcut}};
  \node[anchor=north west,inner sep=0] (nd) at ([xshift=10.8cm,yshift=-5.9cm]current page.north west) {\mapbox{\miniD}{mok}};
  \node[trcap,anchor=north,text=Ink] at ([yshift=-3pt]nr.south) {$\traceLroot$: rota inicial};
  \node[trcap,anchor=north,text=Ink] at ([yshift=-3pt]nb.south) {$\ge\traceLb$};
  \node[trcap,anchor=north,text=Muted] at ([yshift=-3pt]nc.south) {$\traceLc$: descartado};
  \node[trcap,anchor=north,text=Accent] at ([yshift=-3pt]nd.south) {\bfseries $\traceLd$: rota ótima};
  \draw[ed] (nr.east) -- (nb.west);
  \draw[ed] (nb.east) -- (nc.west);
  \draw[ed] (nb.east) -- (nd.west);
  \stepdot{(nr.north west)} {1};
  \stepdot{(nb.north west)} {2};
  \stepdot{(nc.north west)} {3};
  \node[anchor=north west,inner sep=0,font=\sffamily\sz{16}{19},text=Muted] at ([xshift=1.0cm,yshift=-8.3cm]current page.north west) {Comprimentos relativos ao da rota ótima ($=1$)};
  % texto: três passos em linguagem simples
  \foreach \h/\t [count=\i from 0] in {{Rota inicial}/{Vizinho mais próximo: rápido, mas não ótimo ($\traceLroot$).},
                                       {Acrescentar}/{Região por região. O mínimo da rota parcial só cresce.},
                                       {Descartar}/{Rota parcial pior que a melhor conhecida: ramo abandonado.}} {    \node[anchor=north west,inner sep=0,text width=7.9cm] at ([xshift=16.4cm,yshift={-2.35cm-\i*2.2cm}]current page.north west) {\bfseries\sz{26}{30}\hi \the\numexpr\i+1\relax\ \ink\h\par\vspace{2pt}\sz{21}{25}\muted\mdseries \t};
  }
  % força bruta × nosso método (instância de 60 regiões)
  \draw[Rule,line width=0.8pt] ([xshift=1.2cm,yshift=-9.35cm]current page.north west) -- ([xshift=-1.2cm,yshift=-9.35cm]current page.north east);
  \node[anchor=north west,inner sep=0,font=\sffamily\sz{19}{22},text=Muted] at ([xshift=9.2cm,yshift=-9.7cm]current page.north west) {Cálculos};
  \node[anchor=north west,inner sep=0,font=\sffamily\sz{19}{22},text=Muted] at ([xshift=1.2cm,yshift=-9.7cm]current page.north west) {Instância de 60 regiões};
  \node[anchor=north west,inner sep=0,font=\sffamily\sz{19}{22},text=Muted] at ([xshift=17.9cm,yshift=-9.7cm]current page.north west) {Tempo};
  \node[anchor=north west,inner sep=0,font=\sffamily\bfseries\sz{26}{30},text=Muted] at ([xshift=1.2cm,yshift=-10.6cm]current page.north west) {Força bruta};
  \node[anchor=north west,inner sep=0,font=\sffamily\bfseries\sz{26}{30},text=Muted] at ([xshift=9.2cm,yshift=-10.6cm]current page.north west) {$\approx 10^{117}$};
  \node[anchor=north west,inner sep=0,font=\sffamily\bfseries\sz{26}{30},text=Muted] at ([xshift=17.9cm,yshift=-10.6cm]current page.north west) {$\approx 10^{100}$ anos$^{*}$};
  \node[anchor=north west,inner sep=0,font=\sffamily\bfseries\sz{28}{32},text=Teal] at ([xshift=1.2cm,yshift=-11.55cm]current page.north west) {Nosso método};
  \node[anchor=north west,inner sep=0,font=\sffamily\bfseries\sz{28}{32},text=Teal] at ([xshift=9.2cm,yshift=-11.55cm]current page.north west) {$\approx$ 1,4 mil};
  \node[anchor=north west,inner sep=0,font=\sffamily\bfseries\sz{28}{32},text=Accent] at ([xshift=17.9cm,yshift=-11.55cm]current page.north west) {0,1 s};
  \node[anchor=north west,inner sep=0,font=\sffamily\sz{14}{16},text=Muted] at ([xshift=1.2cm,yshift=-12.5cm]current page.north west) {* supondo 1 bilhão de rotas testadas por segundo (a idade do universo é $\approx 10^{10}$ anos)};
  \footerbar{3}
\end{tikzpicture}
\newpage

%%%%%%%%%%%%%%%%%%%%%%%%%%%%%%%% 4. Resultado
% OSM: 314 casos
\def\histOSM{%
(0.625,1)%
(0.875,9)%
(1.125,36)%
(1.375,43)%
(1.625,52)%
(1.875,46)%
(2.125,50)%
(2.375,42)%
(2.625,18)%
(2.875,10)%
(3.125,6)%
(3.375,1)%
}
% random: 160 casos
\def\histrandom{%
(0.625,2)%
(0.875,0)%
(1.125,20)%
(1.375,45)%
(1.625,44)%
(1.875,35)%
(2.125,11)%
(2.375,1)%
(2.625,1)%
(2.875,1)%
(3.125,0)%
(3.375,0)%
}
% tessellation: 78 casos
\def\histvoronoi{%
(0.625,3)%
(0.875,7)%
(1.125,24)%
(1.375,39)%
(1.625,5)%
(1.875,0)%
(2.125,0)%
(2.375,0)%
(2.625,0)%
(2.875,0)%
(3.125,0)%
(3.375,0)%
}
\def\histmedian{1.6136}
\def\histmediantext{41{,}1}
\def\histmean{2.0379}  % média aritmética 109.1x
\def\histmeantext{109{,}1}\def\histcount{552}
\def\histnOSM{314}
\def\histnRANDOM{160}
\def\histnVORONOI{78}
% mediana OSM: 65.4x
% mediana random: 38.3x
% mediana tessellation: 18.8x

\begin{tikzpicture}[remember picture,overlay]
  \fill[white] (current page.south west) rectangle (current page.north east);
  \slidetitle{Solver melhorado: em média, \hi 112$\times$ \ink mais rápido}
  \node[anchor=north west,inner sep=0,text width=22.8cm] at ([xshift=1.2cm,yshift=-1.95cm]current page.north west) {\sffamily\sz{20}{24}\muted Comparação com \textbf{Fekete et al.} (grupo da TU Braunschweig, Alemanha, 2026), que usam um solver genérico (Gurobi). O nosso é um solver geométrico exato.};
  \node[anchor=north west,inner sep=0] at ([xshift=0.9cm,yshift=-3.25cm]current page.north west) {%
  \begin{tikzpicture}
  % largura do eixo 15,5 cm / 4,2 unidades = 3,69 cm por unidade log10; cada bin (0,25) = 26,15 pt
  \begin{axis}[ybar stacked,scale only axis,bar width=26.15pt,width=15.5cm,height=6.8cm,
      xlabel={Speedup = tempo de Fekete et al. / nosso tempo},ylabel={Instâncias},
      label style={font=\sffamily\sz{20}{22}},tick label style={font=\sffamily\sz{20}{22}},xticklabel style={yshift=-6pt},yticklabel style={xshift=-6pt},xlabel style={yshift=-18pt},ylabel style={yshift=16pt},
      xmin=-0.6,xmax=3.6,ymin=0,ymax=165,
      xtick={0,1,2,3},xticklabels={$1\times$,$10\times$,$100\times$,$1000\times$},
      ytick={0,40,80,120,160},ymajorgrids=true,grid style={Rule,line width=0.5pt},axis line style={Ink},
      every axis plot/.append style={draw=white,line width=0.6pt}]
    \addplot[fill=Teal] coordinates {\histOSM};
    \addplot[fill=Accent] coordinates {\histrandom};
    \addplot[fill=Ink] coordinates {\histvoronoi};
    % legenda própria, na ordem visual do gráfico (Voronoi em cima), com quadrados da altura do texto
    \node[anchor=north east,draw=Rule,fill=white,fill opacity=0.92,text opacity=1,inner sep=5pt,font=\sffamily\sz{20}{23},align=left] at (rel axis cs:0.985,0.76) {%
      \tikz[x=1cm,y=1cm,baseline=-0.6ex]\fill[Ink] (0,-0.2) rectangle (0.5,0.3);\ Voronoi (\histnVORONOI)\\[2pt]
      \tikz[x=1cm,y=1cm,baseline=-0.6ex]\fill[Accent] (0,-0.2) rectangle (0.5,0.3);\ Aleatórias (\histnRANDOM)\\[2pt]
      \tikz[x=1cm,y=1cm,baseline=-0.6ex]\fill[Teal] (0,-0.2) rectangle (0.5,0.3);\ OSM (\histnOSM)};
    \draw[Muted,dashed,line width=1.4pt] (axis cs:0,0) -- (axis cs:0,165);
    \draw[Ink,dashed,line width=2.2pt] (axis cs:\histmedian,0) -- (axis cs:\histmedian,148);
    \node[anchor=south east,font=\sffamily\bfseries\sz{20}{23},text=Ink,inner sep=2pt] at (axis cs:\histmedian,146) {mediana $\histmediantext\times$};
    \draw[Accent,line width=2.4pt] (axis cs:\histmean,0) -- (axis cs:\histmean,148);
    \node[anchor=south west,font=\sffamily\bfseries\sz{20}{23},text=Accent,inner sep=2pt] at (axis cs:\histmean,146) {média $\histmeantext\times$};
    \node[anchor=center,rotate=90,font=\sffamily\sz{18}{21},text=Muted,inner sep=0] at (axis cs:-0.3,82) {\hbox to 4cm{\hss Fekete mais rápido: 0\hss}};
  \end{axis}
  \end{tikzpicture}};
  \node[anchor=north west,inner sep=0,text width=5.4cm] at ([xshift=19.3cm,yshift=-3.3cm]current page.north west)
    {{\bfseries\sz{60}{62}\hi 112$\times$}\par\vspace{2pt}{\sz{24}{28}\ink\bfseries média do speedup\par\vspace{0.25cm}\mdseries\muted (mediana: \histmediantext$\times$)}\par\vspace{0.7cm}
     {\sz{18}{22}\muted No pôster: média 10,4$\times$}};
  \node[anchor=south west,inner sep=0,text width=22.8cm] at ([xshift=1.2cm,yshift=1.3cm]current page.south west)
    {\sffamily\sz{16}{19}\muted Solver atualizado após o pôster. \histcount\ casos fechados por ambos; nos 5 restantes o Fekete et al. não terminou (nosso solver é mais rápido nos 558). Uma thread, sem limite de tempo, gap alvo de 0,1\,\%. Intel i9-12900K.};
  \footerbar{4}
\end{tikzpicture}
\newpage

%%%%%%%%%%%%%%%%%%%%%%%%%%%%%%%% 5. Convite ao pôster e ao app
\begin{tikzpicture}[remember picture,overlay]
  \fill[white] (current page.south west) rectangle (current page.north east);
  \slidetitle{Venha ver o pôster}
  \node[anchor=north west,inner sep=0,text width=13.6cm] at ([xshift=1.2cm,yshift=-3.0cm]current page.north west) {%
    \foreach \n/\q in {1/{Por que nosso solver geométrico exato vence o solver genérico (Gurobi) de Fekete et al.?},
                       2/{Em quais tipos de instância vencemos menos?},
                       3/{Como a busca evita testar $10^{117}$ combinações?}} {%
      \parbox[t]{1.1cm}{\bfseries\sz{34}{38}\hi\n}\parbox[t]{12.3cm}{\sz{26}{31}\ink\q}\par\vspace{0.7cm}}};
  \node[anchor=north,inner sep=0,text width=8.4cm,align=center] at ([xshift=-5.4cm,yshift=-2.55cm]current page.north east)
    {\sffamily\bfseries\sz{22}{26}\ink Para ir além, no app:\par\mdseries\sz{20}{24}\muted centenas de instâncias e a simulação passo a passo da busca};
  \node[anchor=north,inner sep=0] (qr) at ([xshift=-5.4cm,yshift=-4.55cm]current page.north east) {\qrcode[level=H,height=7.9cm]{\appurl}};
  % marca do app sobre o centro do QR (nível H de correção), como no pôster
  \begin{scope}[shift={(qr.center)}]
    \fill[white,rounded corners=2pt] (-1.3,-1.3) rectangle (1.3,1.3);
    \begin{scope}[shift={(-1.16,-1.16)},scale=0.03625]
      \fill[Ink,rounded corners=2.5] (0,0) rectangle (64,64);
      \draw[Teal!60!white,line width=4.4pt,line join=round] (32,54) -- (53,14) -- (11,14) -- cycle;
      \draw[Route,line width=3.85pt,line cap=round] (10,28) -- (54,35);
      \fill[white] (10,28) circle[radius=4]; \fill[white] (54,35) circle[radius=4];
    \end{scope}
  \end{scope}
  \footerbar{5}
\end{tikzpicture}
\end{document}
